\documentclass[12pt,answers,letterpaper]{exam}
\usepackage[utf8]{inputenc}
\usepackage[T1]{fontenc}
\usepackage[spanish,es-nodecimaldot]{babel}
\usepackage{mlmodern}
\usepackage{microtype}
\usepackage{amsmath}
\usepackage{graphicx}
\usepackage{needspace}
\usepackage{setspace}
\usepackage{parskip}
\usepackage{tikz}
\usepackage[hidelinks]{hyperref}
\onehalfspacing
\hfuzz=0pt
\AtBeginDocument{\shorthandoff{<>}}

\newlength{\anchools}

\renewcommand{\solutiontitle}{\noindent\textbf{Respuesta:}\par\noindent}
\qformat{\textbf{\large Pregunta \thequestion}\hfill}

\pagestyle{foot}
\footer{}{\thepage}{}

\newcommand{\figura}[1]{\begin{center}[Figura: #1]\end{center}}

% Numeración continua de preguntas entre unidades
\newcommand{\tema}[1]{\addcontentsline{toc}{subsection}{Pregunta \the\numexpr\value{question}+1\relax. #1}}
\newcounter{unitnum}
\newcommand{\unidadtitulo}[1]{\stepcounter{unitnum}\section*{Unidad \Roman{unitnum}: #1}\addcontentsline{toc}{section}{Unidad \Roman{unitnum}: #1}}
\newcounter{globalq}
\newenvironment{unidad}
  {\begin{questions}\setcounter{question}{\value{globalq}}}
  {\setcounter{globalq}{\value{question}}\end{questions}}


\def\bpsc{1}\def\bpe{0.7}
% Gráfico de funciones de reacción en Bertrand (c1,c2: costos; m1,m2: precios de monopolio)
\newcommand{\bertrandplot}[9]{%
\begin{center}
\begin{tikzpicture}[scale=\bpsc,x=0.4cm,y=0.4cm,>=stealth,font=\small,line cap=round]
  \def\e{\bpe}
  \pgfmathsetmacro{\xmax}{#4+3}
  \draw[->] (0,0) -- (\xmax,0) node[right] {$P_1$};
  \draw[->] (0,0) -- (0,\xmax) node[above] {$P_2$};
  \draw[densely dotted,gray] (0,0) -- (\xmax-0.5,\xmax-0.5);
  \draw[densely dashed,gray!70] (#3,0) -- (#3,#3) -- (0,#3);
  \draw[densely dashed,gray!70] (#4,0) -- (#4,#4) -- (0,#4);
  \draw (#1,0.15) -- (#1,-0.15) node[below] {#5};
  \draw (#3,0.15) -- (#3,-0.15) node[below] {#7};
  \draw (0.15,#2) -- (-0.15,#2) node[left] {#6};
  \draw (0.15,#4) -- (-0.15,#4) node[left] {#8};
  % Firma 1: P1(P2)
  \draw[thick,violet] (#1,0) -- (#1,#1+\e) -- (#3-\e,#3) -- (#3,#3) -- (#3,\xmax-0.5);
  % Firma 2: P2(P1)
  \draw[thick,blue!70] (0,#2) -- (#2+\e,#2) -- (#4,#4-\e) -- (#4,#4) -- (\xmax-0.5,#4);
  \fill[red] (#2-#9,#2) circle (2.5pt) node[below right,red] {$N$};
  \node[violet] at (#3-4,#3+1.6) {$P_1(P_2)$};
  \node[blue!70] at (#4+1.8,#4-5) {$P_2(P_1)$};
\end{tikzpicture}
\end{center}}


% Cuadrados de valoraciones (V_x,V_y) para venta separada, atada y empaquetamiento
\newcommand{\cuadrado}[1]{%
\begin{tikzpicture}[x=3.4cm,y=3.4cm,font=\scriptsize]
  \draw[->] (0,0) -- (1.12,0) node[right] {$V_x$};
  \draw[->] (0,0) -- (0,1.12) node[above] {$V_y$};
  \draw (0,0) rectangle (1,1);
  #1
\end{tikzpicture}}

\begin{document}

\vspace*{-1.5cm}
\begin{center}
  \includegraphics[height=2cm]{escudo.jpg}\\[4pt]
  {\Large\bfseries MIC355. Ejercicios\\
  2024-2025}
\end{center}

\tableofcontents
\newpage

Compilado de ejercicios de Microeconomía II, curso de Organización Industrial. Hecho como insumo para ayudantías del mismo curso o cursos relacionados.\\
Contacto: joamartine@fen.uchile.cl

\newpage

\unidadtitulo{Competencia à la Cournot}
\begin{unidad}
%-------------------------------------------------------------------- Pizza
\needspace{8\baselineskip}
\tema{Pizza con anchoas (N firmas simétricas)}
\question Desde tiempos inmemorables en todo Chile se comercia la rica pizza con anchoas. Debido a su alta demanda la producción de esta mítica pizza se ha masificado, pero lamentablemente producto de la complejidad de su preparación, los comercios han optado por vender una misma pizza con anchoas sin diferenciación alguna, por lo que no les queda más opción que competir por cantidad. La demanda inversa es $P = 1000 - Q$, donde $Q$ es la cantidad total de pizzas con anchoas transadas en el mercado. Existen $N$ empresas simétricas compitiendo, donde cada una tiene un costo variable de producción igual a \$40 por unidad.
\begin{parts}
\part Obtenga la función de reacción de una empresa. ¿Depende la función de cómo se reparta la producción entre sus rivales o solo del total de producción del resto? Explique.
\needspace{5\baselineskip}
\begin{solution}
Cada empresa $i$ maximiza sus beneficios. Separando $q_i$ de $Q$:
\[
  \max_{q_i}\ \pi_i = (P-c)\,q_i = \Bigl(1000 - q_i - \sum_{j\neq i} q_j - c\Bigr) q_i .
\]
CPO:
\[
\begin{gathered}
  \frac{\partial \pi_i}{\partial q_i} = 1000 - 2q_i - \sum_{j\neq i} q_j - c = 0 \\
  \Longrightarrow\; q_i(Q_{-i}) = \frac{1000 - c - Q_{-i}}{2},
\end{gathered}
\]
donde $Q_{-i}=\sum_{j\neq i}q_j$ es la cantidad producida por las $N-1$ firmas con las que compite $i$. Reemplazando $c=40$:
\[
  q_i(Q_{-i}) = \frac{960 - Q_{-i}}{2}.
\]
La función de reacción depende solo del \emph{total} producido por el resto ($Q_{-i}$) y no de cómo se reparte entre los rivales, porque la demanda es lineal en $Q$ y a la firma $i$ solo le importa cuánto baja el precio por la cantidad total que hay en el mercado.

\end{solution}

\part Calcule el precio de equilibrio y los beneficios de cada empresa en función de $N$. Explique de modo intuitivo cómo varían estos parámetros a medida que aumenta $N$.
\needspace{5\baselineskip}
\begin{solution}
Partimos de la función de reacción:
\[
  q_i = \frac{960 - \sum_{j\neq i} q_j}{2}
  \;\Longrightarrow\;
  2q_i = 960 - \sum_{j\neq i} q_j
  \;\Longrightarrow\;
  q_i + \sum_{j=1}^{N} q_j = 960 .
\]
Como las empresas son simétricas (mismos costos), $q_i = q$ para todo $i$, y entonces $Q = Nq$:
\[
  q + Nq = 960
  \quad\Longrightarrow\quad
  q^* = \frac{960}{1+N}, \qquad Q^* = \frac{960N}{1+N}.
\]
El precio y los beneficios de equilibrio son
\[
  P^* = 1000 - \frac{960N}{1+N} = \frac{1000 + 40N}{1+N},
  \qquad
  \pi^* = (P^*-c)\,q^* = \Bigl(\frac{960}{1+N}\Bigr)^2 .
\]
Cuando $N$ aumenta, la producción de cada empresa $q^*$ disminuye, pero la cantidad total $Q^*$ aumenta. Como hay más cantidad en el mercado, el precio de equilibrio baja. Como cada empresa vende menos y a un precio menor, los beneficios individuales también disminuyen. En el límite, cuando $N\to\infty$, el equilibrio se asimila al de competencia perfecta: $P^*\to c = 40$ y $\pi^*\to 0$.

\end{solution}

\part Suponga que una de las empresas realiza una innovación en la preparación de esta mítica pizza, lo que le permite reducir su costo variable a \$20. Explique mediante el uso de funciones de reacción cómo cambia el nivel de producción y de beneficios de la firma de menor costo y del resto con el nuevo equilibrio.
\needspace{5\baselineskip}
\begin{solution}
Con costos distintos, la función de reacción de la empresa $i$ es
\[
  q_i(Q_{-i}) = \frac{1000 - c_i - Q_{-i}}{2}.
\]
Esta función está negativamente correlacionada con el costo $c_i$, por lo que al bajar el costo de \$40 a \$20 la función de reacción de la firma innovadora se desplaza hacia afuera y produce más. Las demás empresas producirán menos que en el caso simétrico: como la firma eficiente produce más, su rival ajusta su producción a la baja. En Cournot las cantidades son \emph{sustitutos estratégicos}: si por razones exógenas una empresa aumenta su producción, el rival reacciona disminuyendo la suya. Finalmente, la empresa más eficiente obtiene mayores beneficios que las empresas menos eficientes.

\end{solution}
\end{parts}

%-------------------------------------------------------------------- Asimetría de costos
\needspace{8\baselineskip}
\tema{Cournot con asimetría de costos}
\question Las firmas $i \in \{1,2\}$ compiten à la Cournot. La firma 1 es más eficiente que su competencia, por lo que $c_1 < c_2$.
\begin{parts}
\part Encuentre las funciones de reacción para ambas firmas. Suponga una demanda lineal cualquiera, $P = A - Q$.
\needspace{5\baselineskip}
\begin{solution}
Planteamos el problema de maximización de la firma 1:
\[
  \max_{q_1}\ \pi_1 = (P - c_1)q_1 = (A - q_1 - q_2 - c_1)q_1 .
\]
CPO:
\[
\begin{gathered}
  \frac{\partial \pi_1}{\partial q_1} = A - 2q_1 - q_2 - c_1 = 0 \\
  \Longrightarrow\; q_1^*(q_2) = \frac{A - q_2 - c_1}{2}.
\end{gathered}
\]
El proceso es simétrico para la firma 2:
\[
  q_2^*(q_1) = \frac{A - q_1 - c_2}{2}.
\]

\end{solution}

\part Dadas las funciones de reacción, encuentre la producción de cada una de ellas y dé una explicación intuitiva del resultado.
\needspace{5\baselineskip}
\begin{solution}
Reemplazamos una función de reacción en la otra:
\[
\begin{gathered}  q_1 = \frac{A - c_1}{2} - \frac{1}{2}\cdot\frac{A - q_1 - c_2}{2}
  \;\Longrightarrow\;\\
  4q_1 = 2A - 2c_1 - A + q_1 + c_2
  \;\Longrightarrow\;
  3q_1 = A - 2c_1 + c_2 ,\end{gathered}
\]
de modo que
\[
  q_1^* = \frac{A - 2c_1 + c_2}{3}, \qquad q_2^* = \frac{A - 2c_2 + c_1}{3}.
\]
Como la firma 1 es más eficiente, $q_1^* > q_2^*$:
\[
  \frac{A - 2c_1 + c_2}{3} > \frac{A - 2c_2 + c_1}{3}
  \iff c_2 > c_1 .
\]
La diferencia de producción es
\[
  q_1^* - q_2^* = c_2 - c_1 ,
\]
es decir, la firma eficiente produce más y la brecha crece una unidad por cada unidad adicional de asimetría de costos. Por ejemplo, si la firma 1 fuera el doble de eficiente ($c_2 = 2c_1$), la diferencia sería $q_1^*-q_2^* = c_1$. Un menor costo marginal desplaza hacia afuera la función de reacción de la firma 1, lo que la lleva a producir más y obliga a la rival a producir menos.

\end{solution}
\end{parts}

%-------------------------------------------------------------------- Comente: n firmas
\needspace{8\baselineskip}
\tema{Comente: variación del número de firmas}
\question \textbf{Comente.} Discuta qué ocurre en competencia à la Cournot al variar el número de empresas $n$ que componen el mercado. ¿Qué sentido tiene desde la función de reacción? ¿Qué sucedería al tender a infinito?
\needspace{5\baselineskip}
\begin{solution}
Planteamos el problema matemáticamente. Las firmas están indexadas como $i\in\{1,2,\dots,n\}$ y por simplicidad asumimos que son simétricas, con costo marginal común $c$. La cantidad total y la demanda son
\[
  Q = \sum_{i=1}^{n} q_i = n\cdot q, \qquad Q(P) = A - P .
\]
A mayor cantidad total, menor precio de mercado: si todas las empresas producen lo mismo, mientras más empresas haya, menor será el precio. La demanda inversa es
\[
  P = A - q_1 - q_2 - \dots - q_n = A - n\cdot q .
\]
La intuición es que las empresas se cuidan de no producir mucho para no desplomar el precio (interdependencia propia de estos modelos). Para formalizar la respuesta de la firma 1 ante las demás firmas $-1$:
\[
\begin{gathered}
  \max_{q_1}\ \pi_1 = (A - q_1 - q_2 - \dots - q_n - c)\,q_1 \\
  \Longrightarrow\; \text{CPO: } A - 2q_1 - q_2 - \dots - q_n - c = 0,
\end{gathered}
\]
\[
  q_1^*(q_{-1}) = \frac{A - c - \sum_{i=2}^{n} q_i}{2}.
\]
Esta es la respuesta de la firma 1 ante las decisiones de las demás. Dado que las firmas son simétricas y todas juegan al mismo tiempo, todas producirán lo mismo, por lo que $q_i=q$ y la función de reacción de una empresa cualquiera es
\[
  q = \frac{A - c - q(n-1)}{2}.
\]
Despejando,
\[
  q^* = \frac{A-c}{1+n}, \qquad P^* = \frac{A + nc}{1+n}.
\]
Al aumentar $n$ cada firma produce menos, pero la producción total $Q^* = \frac{n(A-c)}{1+n}$ aumenta y el precio baja. Cuando $n\to\infty$ nos encontramos en competencia perfecta: $P^* \to c$.

\end{solution}

%-------------------------------------------------------------------- Comente: número de firmas
\needspace{8\baselineskip}
\tema{Comente: número de firmas}
\question \textbf{Comente.} Cuando se compite en precios o en cantidades, el número de firmas es irrelevante para el análisis, porque siempre existirán beneficios. \emph{(Hint: asuma firmas simétricas.)}
\needspace{5\baselineskip}
\begin{solution}
Falso. En el caso de Cournot, mientras más firmas hay en el mercado, mayor es la competencia y por lo tanto menor es el precio. En el límite, si existen infinitas firmas el precio tenderá al costo marginal y los beneficios se disiparán. Matemáticamente, asumiendo una demanda del estilo $Q = A - P$ y un costo marginal $c$ para cada firma, el resultado es
\[
  P^* = \frac{A + Nc}{N+1}, \qquad \pi_i^* = \frac{(A-c)^2}{(N+1)^2}.
\]
En cambio, en la competencia à la Bertrand basta que existan dos firmas simétricas para alcanzar el equilibrio competitivo ($P = c$, $\pi = 0$). Por lo tanto, el número de firmas \emph{sí} importa en Cournot, y en Bertrand lo relevante es pasar de una a dos firmas.

\end{solution}

%-------------------------------------------------------------------- Comente: reducción de costos de una firma
\needspace{8\baselineskip}
\tema{Comente: reducción de costos de una firma en Cournot}
\question \textbf{Comente.} Si dos firmas compiten en cantidades (à la Cournot) y solo una de ellas ve reducidos sus costos, la otra firma mantiene su producción inalterada, ya que sus costos no han cambiado.
\needspace{5\baselineskip}
\begin{solution}
Falso. Si una firma ve reducidos sus costos, tenderá a aumentar su producción, y en reacción la otra firma producirá menos. A pesar de que la otra no cambia sus parámetros de costos, cambia su nivel de producción debido a que enfrenta a un competidor más agresivo (las cantidades son sustitutos estratégicos).

\end{solution}

%-------------------------------------------------------------------- Comente: Cournot es Nash
\needspace{8\baselineskip}
\tema{Comente: el equilibrio de Cournot es un equilibrio de Nash}
\question \textbf{Comente.} El equilibrio que se obtiene del modelo de Cournot es un equilibrio de Nash.
\needspace{5\baselineskip}
\begin{solution}
Verdadero. Recordemos que el equilibrio en el modelo de Cournot se obtiene al intersectar las funciones de reacción de cada firma. Por definición, las funciones de reacción nos indican la mejor estrategia (cantidad) que se puede realizar dada la estrategia (cantidad) elegida por el otro jugador. Por lo tanto, cuando todas las firmas están actuando bajo sus funciones de reacción, ninguna de ellas tiene incentivos a desviarse, por lo que este equilibrio es un equilibrio de Nash.

\end{solution}

%-------------------------------------------------------------------- Cournot N firmas (500-Q)
\needspace{8\baselineskip}
\tema{Cournot con $N$ firmas e innovación (variante numérica)}
\question Considere un mercado con $N$ empresas simétricas donde cada una vende un producto homogéneo y compiten en cantidades. Las empresas se enfrentan a una demanda $P = 500 - Q$ con un costo marginal igual a 10.
\begin{parts}
\part Obtenga la función de reacción de una empresa. ¿Depende la función de cómo se reparta la producción entre sus rivales o solo del total de producción del resto? Explique.
\needspace{5\baselineskip}
\begin{solution}
Cada empresa maximiza $\pi_i = \bigl(500 - q_i - \sum_{j\neq i}q_j - 10\bigr)q_i$. La CPO es $490 - 2q_i - \sum_{j\neq i}q_j = 0$, de donde
\[
  q_i(Q_{-i}) = \frac{490 - Q_{-i}}{2},
\]
con $Q_{-i}$ la cantidad producida por las $N-1$ firmas con las que compite $i$. La función de reacción depende solo del total producido por las rivales.

\end{solution}

\part Calcule el precio de equilibrio y los beneficios de cada empresa en función de $N$. Explique de modo intuitivo cómo varían estos parámetros a medida que aumenta $N$.
\needspace{5\baselineskip}
\begin{solution}
Con simetría, $q_i(1+N) = 490$:
\[
\begin{gathered}
  q^* = \frac{490}{1+N}, \quad Q^* = \frac{490N}{1+N}, \\ P^* = 500 - \frac{490N}{1+N} = \frac{500 + 10N}{1+N}, \quad \pi^* = \Bigl(\frac{490}{1+N}\Bigr)^2 .
\end{gathered}
\]
La producción de cada empresa y los beneficios son decrecientes en $N$, mientras que la cantidad de mercado es creciente y el precio decreciente. En el límite $N \to \infty$ el equilibrio se asimila al de competencia perfecta: $P \to 10$ y $\pi \to 0$.

\end{solution}

\part Suponga que una de las empresas realiza una innovación que le permite reducir su costo marginal a 5. Explique mediante el uso de funciones de reacción cómo cambia el nivel de producción y de beneficios de la firma de menor costo y del resto con el nuevo equilibrio.
\needspace{5\baselineskip}
\begin{solution}
La función de reacción de la empresa $i$ es $q_i(Q_{-i}) = \frac{500 - c_i - Q_{-i}}{2}$, decreciente en el costo $c_i$. Si el costo baja de 10 a 5, la firma innovadora produce más; las demás producen menos que en el caso simétrico (las cantidades son sustitutos estratégicos); y la firma más eficiente obtiene mayores beneficios que las menos eficientes.

\end{solution}
\end{parts}

\end{unidad}

\unidadtitulo{Competencia à la Bertrand}
\begin{unidad}
%-------------------------------------------------------------------- Bertrand matemático
\needspace{8\baselineskip}
\tema{Bertrand con costos simétricos y asimétricos}
\question Existen dos firmas que producen un bien homogéneo. Las firmas compiten à la Bertrand en un mercado donde la demanda viene dada por $Q = 125 - 5P$. Los costos totales de ambas firmas vienen dados por la función $C(q) = 3q$.
\begin{parts}
\part Obtenga el equilibrio de Nash bajo los supuestos normales de competencia Bertrand, específicamente el precio, la cantidad producida y los beneficios de cada firma. Además, deje expresadas las funciones de reacción de cada firma y grafíquelas.
\needspace{5\baselineskip}
\begin{solution}
Bajo los supuestos normales de competencia Bertrand, las firmas compiten en precios y se llega a un equilibrio donde el precio es igual al costo marginal de las firmas: $P_1^* = P_2^* = 3$. A ese precio, la cantidad demandada por el mercado es $Q = 125 - 5\cdot 3 = 110$, por lo que la producción de equilibrio de cada firma es $q_1^* = q_2^* = 55$, y los beneficios de ambas firmas son $\pi_1^* = \pi_2^* = 0$.

Las funciones de reacción son (con $P^M = 14$ el precio de monopolio, que resulta de $IMg = 25 - \tfrac{2}{5}Q = 3$, es decir $Q=55$ y $P=14$):
\[
  P_i(P_j) =
  \begin{cases}
    P^M = 14 & \text{si } P_j > P^M \\
    P_j - \epsilon & \text{si } 3 < P_j \leq P^M \\
    3 & \text{si } P_j \leq 3
  \end{cases}
\]
El gráfico correspondiente es (con $\epsilon$ exagerado para que se distingan las curvas):
\bertrandplot{3}{3}{14}{14}{3}{3}{14}{14}{0}

\end{solution}

\part Suponga ahora, solo para este apartado, que los costos totales de la firma 1 son $C(q) = 3q$, mientras que los costos totales de la firma 2 son $C(q) = 4q$. ¿Cómo cambian los resultados obtenidos en la parte (a)? Obtenga los nuevos parámetros de equilibrio y las nuevas funciones de reacción de ambas firmas. Grafique estas últimas.
\needspace{5\baselineskip}
\begin{solution}
Puesto que ahora la firma 1 es más eficiente que la firma 2, el nuevo precio de equilibrio es igual al costo marginal de la firma ineficiente menos un épsilon: $P^* = 4 - \epsilon$. A ese precio la firma ineficiente es incapaz de ofrecer un precio menor, y la firma eficiente se lleva toda la demanda. Por lo tanto, el mercado demanda $Q = 125 - 5\cdot 4 = 105$, y la producción de equilibrio es $q_1^* = 105$ y $q_2^* = 0$. Los beneficios de la firma eficiente son $\pi_1^* \approx 105$ (pues $(4-\epsilon-3)\cdot 105$) y los de la firma ineficiente son $\pi_2^* = 0$.

Las funciones de reacción de cada firma ahora son (con $P_1^M = 14$ y $P_2^M = 14{,}5$):
\[
  P_1(P_2) =
  \begin{cases}
    P_1^M = 14 & \text{si } P_2 > P_1^M \\
    P_2 - \epsilon & \text{si } 3 < P_2 \leq P_1^M \\
    3 & \text{si } P_2 \leq 3
  \end{cases}
\]
\[
  P_2(P_1) =
  \begin{cases}
    P_2^M = 14{,}5 & \text{si } P_1 > P_2^M \\
    P_1 - \epsilon & \text{si } 4 < P_1 \leq P_2^M \\
    4 & \text{si } P_1 \leq 4
  \end{cases}
\]
Y el nuevo gráfico es:
\bertrandplot{3}{4}{14}{14.5}{3}{4}{14}{14,5}{0.7}

\end{solution}
\end{parts}

%-------------------------------------------------------------------- Capacidad
\needspace{8\baselineskip}
\tema{Bertrand con restricciones de capacidad}
\question Suponga un mercado con dos firmas que compiten en precios, ofreciendo un producto homogéneo. Considere además que la demanda de mercado corresponde a $Q(p) = 90 - 3p$, y las funciones de costo de las firmas son idénticas e iguales a $C(q) = 5q$. Finalmente, es ampliamente conocido que las firmas poseen una capacidad máxima $k$ a producir.
\begin{parts}
\part ¿Qué sucede si $k \geq 75$? ¿Qué sucede en caso de que $k < 75$?
\needspace{5\baselineskip}
\begin{solution}
Dado que en equilibrio de Bertrand tenemos $p = c = 5$, la demanda será
\[
  Q(p=5) = 90 - 3\cdot 5 = 75 .
\]
Por lo que si $k \geq 75$ cada firma es capaz de abastecer a todo el mercado y se genera competencia à la Bertrand. Si $k < 75$, una sola firma no es capaz de abastecer a todo el mercado y habrá espacio para actuar sobre la demanda residual.

\end{solution}

\part Ahora considere la situación en que $k = 15$. ¿Qué sucede con las firmas en el mercado?
\needspace{5\baselineskip}
\begin{solution}
En este tipo de casos con capacidad simétrica las empresas se adelantan a lo que hará la otra firma. Desde el punto de vista de la firma 1, se asume que la firma 2 vende toda su capacidad ($k=15$) a un precio menor y, por racionamiento eficiente, los más dispuestos a pagar compran esa oferta. Queda entonces una demanda residual sobre la cual la firma 1 maximiza beneficios:
\[
  Q^r = 90 - k - 3p = 75 - 3p
\]
\[
\begin{gathered}  \max_{p_1}\ \pi_1 = (p_1 - 5)(75 - 3p_1)
  \quad\Longrightarrow\quad\\
  \text{CPO: } 75 - 3p_1 - 3(p_1-5) = 90 - 6p_1 = 0
  \quad\Longrightarrow\quad p_1 = 15 .\end{gathered}
\]
La demanda residual a ese precio es $Q^r = 75 - 3\cdot 15 = 30$. Sin embargo, esto no es factible: como $k = 15 < 30$, a pesar de que la demanda sea 30 no se puede satisfacer por completo. Dado que no se puede ofrecer 30, se ofrece lo máximo (15), por lo que lo racional es subir el precio hasta que solo se vendan 15 unidades. Debemos encontrar el precio donde la cantidad demandada residual es igual a 15:
\[
  75 - 3p = 15 \quad\Longrightarrow\quad p = 20 .
\]
La firma 1 ofrece entonces su capacidad máxima $k=15$ a un precio de 20. La firma 2 hace el mismo razonamiento, pensando que la firma 1 venderá toda su capacidad. Por tanto, se concluye que las dos firmas venden 15 unidades cada una a un precio de 20 (con $p>c=5$ y beneficios positivos).

\end{solution}

\part ¿Qué ocurriría en caso de que las empresas pudieran primero elegir sus capacidades y luego competir por precios?
\needspace{5\baselineskip}
\begin{solution}
Si las empresas eligen primero sus capacidades y luego sus precios, elegirían capacidades iguales a las cantidades de Cournot y precios iguales al precio de mercado en competencia à la Cournot (resultado de Kreps y Scheinkman). Esto implica que, con restricciones de capacidad (un supuesto muy realista), los precios están por encima del costo marginal y las empresas ganan beneficios positivos. En realidad, lo que observamos en el mercado es idéntico a lo que observaríamos si compitieran en cantidades: suponer que las empresas elegían cantidades no era algo tan erróneo como parecía inicialmente.
\begin{itemize}
  \item Cuando hay restricciones de capacidad se suaviza la competencia: los precios de equilibrio no son tan bajos, $p > c$, y las empresas tienen beneficios positivos.
  \item Las empresas evitan acumular demasiada capacidad para suavizar la competencia en precios; es como un compromiso de que no bajarán mucho los precios.
  \item Ejemplos donde la elección de capacidad es importante: hoteles, líneas aéreas.
  \item El resultado del juego en 2 etapas coincide con el de Cournot si las capacidades se interpretan como cantidades.
\end{itemize}

\end{solution}
\end{parts}

%-------------------------------------------------------------------- Comente: Bertrand costos asimétricos
\needspace{8\baselineskip}
\tema{Comente: Bertrand con costos asimétricos}
\question \textbf{Comente.} Suponga que existen dos firmas, 1 y 2, que compiten à la Bertrand. La firma 1 tiene un $C_{mg} = A$ y la firma 2 tiene un costo marginal $C_{mg} = B$, donde se cumple que $A > B$. Por lo tanto, la competencia en precios llevará a que el precio de equilibrio sea $P = B$.
\needspace{5\baselineskip}
\begin{solution}
Falso. La firma más eficiente, en este caso la firma 2, colocará un precio igual al costo marginal de la firma ineficiente menos un $\epsilon$, pues de esta forma la eliminará del mercado. Colocar un precio diferente a este no sería una estrategia racional para la firma. En este caso, el precio de equilibrio es $P = A - \epsilon$.

\end{solution}

%-------------------------------------------------------------------- Comente: Paradoja de Bertrand
\needspace{8\baselineskip}
\tema{Comente: Paradoja de Bertrand}
\question \textbf{Comente.} Explique de forma detallada en qué consiste la Paradoja de Bertrand.
\needspace{5\baselineskip}
\begin{solution}
La paradoja de Bertrand describe la situación en que basta que 2 firmas simétricas compitan en precios (à la Bertrand) para obtener un equilibrio competitivo, es decir, un equilibrio de Nash donde el precio de mercado es igual al costo marginal. Esto ocurre porque la competencia en precios hará que las firmas recorten sus precios hasta llegar al mínimo permitido, su costo marginal, obteniendo beneficios iguales a 0.

\end{solution}

%-------------------------------------------------------------------- Comente: disipación de rentas
\needspace{8\baselineskip}
\tema{Comente: disipación de rentas}
\question \textbf{Comente.} La competencia en precios en un mercado lleva a la disipación total de las rentas.
\needspace{5\baselineskip}
\begin{solution}
Incierto. La afirmación es verdadera siempre que las firmas que compiten sean simétricas en cuanto a costos.

En un caso alternativo en el que dos firmas compiten pero el costo marginal de la firma 1 es menor que el de la firma 2 ($c_1<c_2$), el equilibrio de Nash se logra con la firma 1 abasteciendo a todo el mercado y cobrando un precio infinitesimalmente inferior al costo marginal de la firma 2.

Dado que la firma 2 no puede fijar un precio menor que $p_2 = c_2$, la mejor estrategia de la firma 1 según su función de reacción es la situación intermedia: $p_1 = p_2 - \epsilon$ con $\epsilon\to 0$. Con demanda $P = A - Q$, la función de reacción de la firma 1 es
\[
  p_1^*(p_2) =
  \begin{cases}
    p^M = \dfrac{A + c_1}{2} & \text{si } p_2 > p^M \\[2mm]
    p_2 - \epsilon & \text{si } c_1 < p_2 \leq p^M \\[1mm]
    c_1 & \text{si } p_2 \leq c_1
  \end{cases}
\]
En este caso las rentas de la firma 1 son
\[
  \pi_1 = (p - c_1)\,q_1(p) = (c_2 - \epsilon - c_1)\bigl(A - (c_2 - \epsilon)\bigr)
  \;\longrightarrow\; (c_2 - c_1)(A - c_2) > 0 ,
\]
por lo que las rentas no se disipan totalmente.

\end{solution}

%-------------------------------------------------------------------- Bertrand 240-3P
\needspace{8\baselineskip}
\tema{Bertrand con costos simétricos y asimétricos (variante numérica)}
\question Existen dos firmas que producen un bien homogéneo. Las firmas compiten à la Bertrand en un mercado donde la demanda viene dada por $Q = 240 - 3P$. Los costos totales de ambas firmas vienen dados por la función $C(q) = 2q$.
\begin{parts}
\part Obtenga el equilibrio de Nash bajo los supuestos normales de competencia Bertrand, específicamente el precio, la cantidad producida y los beneficios de cada firma. Además, deje expresadas las funciones de reacción de cada firma y grafíquelas.
\needspace{5\baselineskip}
\begin{solution}
El equilibrio es $P_1^* = P_2^* = 2$. A ese precio la cantidad demandada es $Q = 240 - 3\cdot 2 = 234$, por lo que $q_1^* = q_2^* = 117$ y $\pi_1^* = \pi_2^* = 0$.

El precio de monopolio resulta de $IMg = 80 - \frac{2}{3}Q = 2$, es decir $Q = 117$ y $P^M = 41$. Las funciones de reacción son
\[
  P_i(P_j) =
  \begin{cases}
    P^M = 41 & \text{si } P_j > P^M \\
    P_j - \epsilon & \text{si } 2 < P_j \leq P^M \\
    2 & \text{si } P_j \leq 2
  \end{cases}
\]
Su gráfico tiene la misma forma que en la pregunta anterior, con $c = 2$ y $P^M = 41$:
{\def\bpsc{0.3}\def\bpe{2.2}\bertrandplot{2}{2}{41}{41}{2}{2}{41}{41}{0}}

\end{solution}

\part Suponga ahora, solo para este apartado, que los costos totales de la firma 1 son $C(q) = 2q$, mientras que los costos totales de la firma 2 son $C(q) = 6q$. ¿Cómo cambian los resultados obtenidos en la parte (a)? Obtenga los nuevos parámetros de equilibrio y las nuevas funciones de reacción de ambas firmas.
\needspace{5\baselineskip}
\begin{solution}
Ahora la firma 1 es más eficiente que la firma 2, por lo que el precio de equilibrio es igual al costo marginal de la firma ineficiente menos un épsilon: $P^* = 6 - \epsilon$. A ese precio la firma ineficiente es incapaz de ofrecer un precio menor y la firma eficiente se lleva toda la demanda: $Q = 240 - 3\cdot 6 = 222$, $q_1^* = 222$, $q_2^* = 0$. Los beneficios son $\pi_1^* \approx (6-2)\cdot 222 = 888$ y $\pi_2^* = 0$.

Las funciones de reacción son (con $P_1^M = 41$ y $P_2^M = 43$, que resultan de $IMg = c_i$):
\[
  P_1(P_2) =
  \begin{cases}
    41 & \text{si } P_2 > 41 \\
    P_2 - \epsilon & \text{si } 2 < P_2 \leq 41 \\
    2 & \text{si } P_2 \leq 2
  \end{cases}
  \qquad
  P_2(P_1) =
  \begin{cases}
    43 & \text{si } P_1 > 43 \\
    P_1 - \epsilon & \text{si } 6 < P_1 \leq 43 \\
    6 & \text{si } P_1 \leq 6
  \end{cases}
\]

\end{solution}
\end{parts}

%-------------------------------------------------------------------- Bertrand 8-P
\needspace{8\baselineskip}
\tema{Bertrand: cuántas firmas operan (variante numérica)}
\question Dos firmas compiten en un mercado con un bien homogéneo bajo un modelo de competencia à la Bertrand. La función de demanda del mercado está dada por $Q(P) = 8 - P$, para $0 < P < 8$, y es cero para cualquier otro precio. Los costos de producción son idénticos para ambas firmas, con una función de costos totales $C_i(q_i) = 2q_i$, con $i = 1,2$.
\begin{parts}
\part Obtenga el equilibrio de Nash, específicamente el precio, la cantidad producida y los beneficios de cada firma. Además, deje expresadas las funciones de reacción de cada firma y grafíquelas.
\needspace{5\baselineskip}
\begin{solution}
El precio de equilibrio es igual al costo marginal: $P_1^* = P_2^* = 2$. La cantidad demandada es $Q = 6$, por lo que $q_1^* = q_2^* = 3$ y $\pi_1^* = \pi_2^* = 0$. El precio de monopolio es $P^M = 5$ (de maximizar $(P-2)(8-P)$). Las funciones de reacción son
\[
  P_i(P_j) =
  \begin{cases}
    5 & \text{si } P_j > 5 \\
    P_j - \epsilon & \text{si } 2 < P_j \leq 5 \\
    2 & \text{si } P_j \leq 2
  \end{cases}
\]
\bertrandplot{2}{2}{5}{5}{2}{2}{5}{5}{0}

\end{solution}

\part Suponga ahora que la estructura de costos de las empresas es $C_1(q_1) = q_1$ y $C_2(q_2) = 3q_2$. Escriba las funciones de beneficios para ambos competidores.
\needspace{5\baselineskip}
\begin{solution}
Para los casos en que los precios son inferiores al costo marginal los beneficios son cero (nadie vende a pérdida). Luego,
\[
  \pi_1 =
  \begin{cases}
    (8 - P_1)(P_1 - 1) & \text{si } P_1 < P_2 \\
    \frac{1}{2}(8 - P_1)(P_1 - 1) & \text{si } P_1 = P_2 (> 1) \\
    0 & \text{si } P_1 > P_2
  \end{cases}
\]
\[
  \pi_2 =
  \begin{cases}
    (8 - P_2)(P_2 - 3) & \text{si } P_2 < P_1 \\
    \frac{1}{2}(8 - P_2)(P_2 - 3) & \text{si } P_2 = P_1 (> 3) \\
    0 & \text{si } P_2 > P_1
  \end{cases}
\]

\end{solution}

\part ¿Cuántas empresas operarán en el mercado? Obtenga los nuevos parámetros de equilibrio y las nuevas funciones de reacción de ambas firmas. Grafique estas últimas.
\needspace{5\baselineskip}
\begin{solution}
En equilibrio solo la empresa 1 estará en el mercado, porque es más eficiente que la empresa 2 y puede fijar un precio levemente menor al costo marginal de la empresa 2 y aun así tener beneficios estrictamente positivos. El precio de equilibrio es $P_1 = 3 - \epsilon$, de modo que $q_1 = Q = 8 - 3 + \epsilon = 5 + \epsilon$. Los beneficios son $\pi_1 = (2 - \epsilon)(5 + \epsilon) \approx 10$ y $\pi_2 = 0$.

Los precios de monopolio son $P_1^M = 4{,}5$ (con $c_1 = 1$) y $P_2^M = 5{,}5$ (con $c_2 = 3$), y las funciones de reacción son
\[
  P_1(P_2) =
  \begin{cases}
    4{,}5 & \text{si } P_2 > 4{,}5 \\
    P_2 - \epsilon & \text{si } 1 < P_2 \leq 4{,}5 \\
    1 & \text{si } P_2 \leq 1
  \end{cases}
  \qquad
  P_2(P_1) =
  \begin{cases}
    5{,}5 & \text{si } P_1 > 5{,}5 \\
    P_1 - \epsilon & \text{si } 3 < P_1 \leq 5{,}5 \\
    3 & \text{si } P_1 \leq 3
  \end{cases}
\]
\bertrandplot{1}{3}{4.5}{5.5}{1}{3}{4,5}{5,5}{0.7}

\end{solution}
\end{parts}

%-------------------------------------------------------------------- Capacidad asimétrica
\needspace{8\baselineskip}
\tema{Bertrand con capacidades distintas (40 y 30)}
\question Dos empresas, $A$ y $B$, compiten en precios en el mercado de un bien homogéneo. La demanda agregada del mercado está dada por $Q = 100 - P$, donde $P$ es el precio más bajo ofrecido en el mercado. Cada empresa produce a un costo marginal de 10, pero tienen capacidades limitadas $k_A = 40$ y $k_B = 30$.
\begin{parts}
\part ¿Cuál sería el equilibrio de Bertrand sin restricciones de capacidad? ¿Por qué este resultado no se mantiene bajo las restricciones $k_A$, $k_B$?
\needspace{5\baselineskip}
\begin{solution}
Sin restricciones de capacidad estaríamos frente al modelo más básico de Bertrand, donde el precio es igual al costo marginal ($P = 10$) y la demanda a ese precio es $90$. Bajo restricciones de capacidad, la cantidad máxima que podrían proveer las empresas no suma más de $70$. Dado que con el precio anterior la demanda efectiva ($90$) superaría a la factible ($70$), ese resultado no se mantiene.

\end{solution}

\part Considerando las restricciones de capacidad, determine los precios de equilibrio y la cantidad que vende cada empresa.
\needspace{5\baselineskip}
\begin{solution}
Para determinar el precio de equilibrio, el objetivo es vaciar el mercado para la capacidad factible: $k_A + k_B = 100 - P$, es decir, $70 = 100 - P$, y entonces $P = 30$. Cada empresa vende toda su capacidad: $A$ vende $40$ y $B$ vende $30$.

\end{solution}

\part Analice si alguna de las empresas tiene incentivos a desviarse y reducir su precio.
\needspace{5\baselineskip}
\begin{solution}
Si una de las empresas intentara reducir su precio por debajo de $P = 30$, ya vende toda su capacidad, por lo que solo disminuiría su margen sin lograr más ventas. Si, por otro lado, una empresa sube el precio, los consumidores comprarían primero de la empresa con precio más bajo hasta agotar su capacidad, y la demanda residual que le quedaría a la de precio mayor ($100 - P - k_{rival}$) la llevaría a vender menos de su capacidad a un precio mayor; en este caso, no hay incentivo a subir el precio dado que el margen más alto no compensa la pérdida de ventas. Por lo tanto, $P = 30$ se mantiene como equilibrio.

\end{solution}

\part ¿Cómo cambiaría el resultado si ambas empresas tuvieran una capacidad de 50 unidades?
\needspace{5\baselineskip}
\begin{solution}
Si ambas empresas pudieran producir hasta 50 unidades cada una, la capacidad conjunta sería $100 \geq 90$, por lo que podrían abastecer todo el mercado a $P = 10$. En este caso volvemos al resultado del modelo de Bertrand sin restricciones, con $P = 10$ (ninguna empresa tiene incentivos a subir su precio, pues la competencia las empuja a fijar un precio igual al costo marginal). Con restricciones de capacidad, en cambio, el precio es $P = 30$, significativamente mayor que el costo marginal: las restricciones de capacidad limitan la competencia y permiten a las empresas obtener un precio superior al costo marginal.

\end{solution}
\end{parts}

\end{unidad}

\unidadtitulo{Cournot versus Bertrand}
\begin{unidad}
%-------------------------------------------------------------------- Comente: primero en mover
\needspace{8\baselineskip}
\tema{Comente: primer movimiento}
\question \textbf{Comente.} En modelos de competencia imperfecta, siempre es preferible fijar primero el precio (o la cantidad) que el rival.
\needspace{5\baselineskip}
\begin{solution}
Falso. La decisión depende de la variable con la cual se compita. Si son cantidades, es decir, à la Cournot, efectivamente es preferible mover primero, pero si se compite en precios, à la Bertrand, se obtienen mayores beneficios no moviendo en primer lugar.

\end{solution}

%-------------------------------------------------------------------- Comente: asimétricas
\needspace{8\baselineskip}
\tema{Comente: firmas asimétricas}
\question \textbf{Comente.} En un mercado con dos firmas asimétricas, la firma ineficiente siempre prefiere competir bajo Cournot, mientras que la firma eficiente siempre preferirá Bertrand.
\needspace{5\baselineskip}
\begin{solution}
Bertrand representa un nivel de competencia mayor. Si la superioridad en costos de la firma eficiente respecto de la ineficiente es superior a cierto umbral, la firma eficiente preferirá un nivel de competencia mayor (Bertrand). Si la diferencia en costos está por debajo de cierto umbral, la firma eficiente preferirá competir bajo Cournot. Es verdadero que la firma ineficiente siempre preferirá competir bajo Cournot, pues compitiendo bajo Bertrand obtendrá beneficios iguales a cero. Por lo tanto, la afirmación es falsa en cuanto a la firma eficiente.

\end{solution}

%-------------------------------------------------------------------- Comente: comparar
\needspace{8\baselineskip}
\tema{Comente: Cournot vs Bertrand}
\question \textbf{Comente.} Compare las competencias à la Cournot y Bertrand. ¿Qué mercados son mejor explicados por cada modelo?
\needspace{5\baselineskip}
\begin{solution}
La característica más evidente es la variable que manejan las empresas para competir (cantidades en Cournot, precios en Bertrand). En Cournot la competencia es menos fuerte que en Bertrand.

Si la capacidad y la producción se pueden ajustar rápidamente, el modelo de Bertrand es más apropiado (software, seguros, banca). Si la capacidad no se puede ajustar en el corto plazo, el modelo de Cournot es más apropiado (cemento, autos, electricidad, \emph{commodities} en general).

\end{solution}

\end{unidad}

\unidadtitulo{Stackelberg y juegos secuenciales}
\begin{unidad}
%-------------------------------------------------------------------- Comente: Stackelberg
\needspace{8\baselineskip}
\tema{Comente: Stackelberg vs Cournot}
\question \textbf{Comente.} Considere 2 empresas. Si comparamos con los resultados obtenidos en una competencia simultánea à la Cournot, tendremos que tanto los consumidores como la firma líder estarán mejor bajo una competencia secuencial tipo Stackelberg.
\needspace{5\baselineskip}
\begin{solution}
Verdadero. Ante un juego secuencial tipo Stackelberg, la firma líder expandirá su producción respecto al juego simultáneo debido a que sabe que la firma seguidora ajustará su producción a la baja para no deprimir demasiado el precio. Lo anterior se traduce en mayores niveles de producción en el agregado, menores precios y, por lo tanto, mayores beneficios para los consumidores y para la firma líder (respecto al juego simultáneo).

\end{solution}

%-------------------------------------------------------------------- Stackelberg numérico
\needspace{8\baselineskip}
\tema{Cournot y Stackelberg con costos asimétricos}
\question La curva de demanda inversa de un bien está dada por $P(Q) = 240 - 2Q$. Existen dos firmas con funciones de costos $CT_1(q_1) = 40q_1$ y $CT_2(q_2) = 80q_2$.
\begin{parts}
\part Si ambas firmas compiten à la Cournot, determine las funciones de reacción de ambas. ¿Cuál es el nivel de producción y el precio de equilibrio en este caso? ¿Cuáles son los beneficios de ambas empresas?
\needspace{5\baselineskip}
\begin{solution}
Cada firma maximiza sus utilidades:
\begin{gather*}
  \max_{q_1}\ \pi_1 = (240 - 2Q - 40)q_1 \;\Rightarrow\; \frac{\partial\pi_1}{\partial q_1} = 240 - 2q_2 - 4q_1 - 40 = 0 \\ \;\Rightarrow\; q_1(q_2) = \frac{200 - 2q_2}{4}, \\
  \max_{q_2}\ \pi_2 = (240 - 2Q - 80)q_2 \;\Rightarrow\; \frac{\partial\pi_2}{\partial q_2} = 240 - 2q_1 - 4q_2 - 80 = 0 \\ \;\Rightarrow\; q_2(q_1) = \frac{160 - 2q_1}{4}.
\end{gather*}
Reemplazando una en la otra se obtiene
\[
  q_1^* = 40, \qquad q_2^* = 20, \qquad P^* = 240 - 2\cdot 60 = 120, \qquad \pi_1^* = 3200, \qquad \pi_2^* = 800 .
\]

\end{solution}

\part Suponga que las dos empresas de la industria se comportan à la Stackelberg (en cantidad), de modo que la empresa 1 actúa como líder y la 2 como seguidora. ¿Cuál sería la producción y el precio de equilibrio de mercado? ¿Cuánto producirá la empresa 1? ¿Cuánto producirá la empresa 2? ¿Cuáles serán los beneficios de ambas empresas?
\needspace{5\baselineskip}
\begin{solution}
Es un juego secuencial en dos etapas, que se resuelve por inducción hacia atrás. En $t=2$ la seguidora maximiza
\[
\begin{gathered}
  \max_{q_2}\ \pi_2 = (240 - 2(q_1+q_2))q_2 - 80q_2 \\
  \;\Rightarrow\; 240 - 2q_1 - 4q_2 - 80 = 0 \\
  \;\Rightarrow\; q_2(q_1) = \frac{160 - 2q_1}{4}.
\end{gathered}
\]
La firma 1 conoce la estrategia de la firma 2, por lo que escoge su cantidad óptima usando esta información (reemplaza la restricción $q_2 = q_2(q_1)$):
\[
  \max_{q_1}\ \pi_1 = \Bigl(240 - 2\bigl(q_1 + \tfrac{160 - 2q_1}{4}\bigr)\Bigr)q_1 - 40q_1
  \;\Rightarrow\; 160 - 2q_1 - 40 = 0 \;\Rightarrow\; q_1^* = 60 .
\]
La mejor respuesta de la firma 2 es producir $q_2^* = \frac{160-120}{4} = 10$. La producción total es $Q^* = 70$ y el precio $P^* = 240 - 140 = 100$. Las utilidades son
\[
  \pi_1^* = (100-40)\cdot 60 = 3600, \qquad \pi_2^* = (100-80)\cdot 10 = 200 .
\]
Comparando con Cournot, la líder produce y gana más, la seguidora produce y gana menos, y el precio es menor.

\end{solution}
\end{parts}

%-------------------------------------------------------------------- Comente: timing
\needspace{8\baselineskip}
\tema{Comente: timing de las decisiones en cantidades}
\question \textbf{Comente.} En una competencia de cantidades, y asumiendo simetría en la estructura de costos, da lo mismo si las firmas deciden el tamaño de producción en el mismo período o en períodos distintos, pues se llega al mismo equilibrio.
\needspace{5\baselineskip}
\begin{solution}
Falso. El \emph{timing} de la toma de decisiones de las firmas sí influye en el resultado final de equilibrio. Cuando las firmas deciden sus producciones de forma simultánea, nos encontramos en el típico caso de competencia à la Cournot, donde la producción y los beneficios son iguales entre ambas firmas. Cuando una firma decide su producción antes que otra, esto corresponde a un modelo Stackelberg, donde la firma que decide primero su producción tendrá una mayor participación de mercado y mayores beneficios.

\end{solution}

%-------------------------------------------------------------------- Comente: credibilidad
\needspace{8\baselineskip}
\tema{Comente: credibilidad en Stackelberg}
\question \textbf{Comente.} El elemento clave en el modelo de Stackelberg es la credibilidad. En concreto, si no es creíble que alguna firma pueda ajustar su plan de producción, entonces los resultados de este juego secuencial serán los mismos que los del juego simultáneo.
\needspace{5\baselineskip}
\begin{solution}
Falso. La credibilidad juega un rol clave en el modelo secuencial al interpretar el papel de una amenaza creíble. Para que una firma sea líder, necesitamos que, por ejemplo, no pueda ajustar su producción, es decir, que el compromiso de no ajustar su producción sea creíble. Si la firma seguidora sabe que tal compromiso es creíble (o, equivalentemente, que no es creíble que la firma líder pueda ajustar su producción), entonces efectivamente tendremos un juego secuencial, cuyos resultados serán diferentes al juego simultáneo.

\end{solution}

%-------------------------------------------------------------------- Comente: líder produce menos
\needspace{8\baselineskip}
\tema{Comente: cantidad del líder en Stackelberg}
\question \textbf{Comente.} En la competencia à la Stackelberg, el \emph{timing} indica que la firma que mueve primero producirá una cantidad menor para cobrar un precio alto del bien que produce y, por ende, maximizar sus beneficios.
\needspace{5\baselineskip}
\begin{solution}
Falso. La firma que mueve primero escogerá una cantidad \emph{mayor}. Luego, la firma seguidora escogerá una cantidad menor para no deprimir el precio y obtener beneficios positivos.

\end{solution}

%-------------------------------------------------------------------- Comente: Cournot vs Stackelberg beneficios
\needspace{8\baselineskip}
\tema{Comente: producción y beneficios del líder y seguidor}
\question \textbf{Comente.} En competencia à la Cournot con dos firmas, la cantidad producida por la firma que escoge primero en el modelo de Stackelberg es mayor que en el modelo de decisiones simultáneas, y los beneficios de las firmas también son mayores.
\needspace{5\baselineskip}
\begin{solution}
La primera afirmación es verdadera, ya que la firma líder sí producirá más que en el caso de decisiones simultáneas. Sin embargo, la segunda parte del comente es falsa, ya que solo aumentarán los beneficios de la firma que escoge primero, mientras que en el caso de la firma seguidora sus beneficios serán menores que si las decisiones fuesen simultáneas.

\end{solution}

%-------------------------------------------------------------------- Comente: ventaja primer movimiento
\needspace{8\baselineskip}
\tema{Comente: ventaja de mover primero}
\question \textbf{Comente.} En Stackelberg, siempre será ventajoso mover primero.
\needspace{5\baselineskip}
\begin{solution}
Falso. Cuando la competencia es en precios, resulta ventajoso mover después, dado que puedo reducir el precio ofrecido por la primera firma y quedarme con una mayor demanda.

\end{solution}

%-------------------------------------------------------------------- Stackelberg precios distintos
\needspace{8\baselineskip}
\tema{Stackelberg con funciones de precio distintas}
\question En una economía solo existen dos firmas que abastecen todo el mercado. Sus funciones de precio son
\begin{align*}
  P_A &= 80 - \tfrac{3}{2}q_A - q_B \\
  P_B &= 80 - \tfrac{3}{2}q_B - q_A
\end{align*}
Asuma que ninguna posee costos de producción y que la firma $A$ es líder y $B$ seguidora.
\begin{parts}
\part Suponga que las dos empresas de la industria se comportan à la Stackelberg (en cantidad). ¿Cuál sería la producción y el precio de equilibrio de mercado? ¿Cuánto producirá la empresa $A$? ¿Cuánto producirá la empresa $B$? ¿Cuáles serán los beneficios de ambas empresas?
\needspace{5\baselineskip}
\begin{solution}
Es un juego secuencial en dos etapas, que se resuelve por inducción hacia atrás. En $t=2$ la seguidora maximiza
\[
  \max_{q_B}\ \pi_B = \bigl(80 - \tfrac{3}{2}q_B - q_A\bigr)q_B
  \;\Rightarrow\; 80 - 3q_B - q_A = 0
  \;\Rightarrow\; q_B(q_A) = \frac{80}{3} - \frac{1}{3}q_A .
\]
La firma $A$ conoce la estrategia de la firma $B$, por lo que escoge su cantidad óptima usando esta información:
\[
  \max_{q_A}\ \pi_A = \bigl(80 - \tfrac{3}{2}q_A - q_B\bigr)q_A
  \quad \text{s.a.} \quad q_B = \frac{80}{3} - \frac{1}{3}q_A .
\]
Reemplazando la restricción,
\[
  \max_{q_A}\ \pi_A = \Bigl(80 - \tfrac{3}{2}q_A - \tfrac{80}{3} + \tfrac{1}{3}q_A\Bigr)q_A
  \;\Rightarrow\; \frac{160}{3} - \frac{7}{3}q_A = 0
  \;\Rightarrow\; q_A^* = \frac{160}{7} \approx 22{,}86 .
\]
Por lo tanto, la mejor respuesta de la firma $B$ es producir $q_B^* = \frac{400}{21} \approx 19{,}05$, y la producción total es $Q^* \approx 41{,}90$. Los precios (distintos, pues las funciones de precio lo son) y los beneficios son
\[
  P_A^* = \tfrac{80}{3} \approx 26{,}67, \quad P_B^* = \tfrac{200}{7} \approx 28{,}57, \qquad
  \pi_A^* \approx 609{,}5, \quad \pi_B^* \approx 544{,}2 .
\]

\end{solution}

\part \emph{(Propuesto)} Suponga ahora que las empresas enfrentan costos $C_{mg,A} = 2$ y $C_{mg,B} = 3$. Encuentre lo solicitado en (a) con la nueva información.
\needspace{5\baselineskip}
\begin{solution}
Se procede igual, con los nuevos costos. La seguidora resuelve $80 - 3q_B - q_A - 3 = 0$, es decir $q_B(q_A) = \frac{77}{3} - \frac{1}{3}q_A$. La líder maximiza
\[
  \max_{q_A}\ \pi_A = \Bigl(80 - \tfrac{3}{2}q_A - \tfrac{77}{3} + \tfrac{1}{3}q_A - 2\Bigr)q_A
  \;\Rightarrow\; \frac{157}{3} - \frac{7}{3}q_A = 0
  \;\Rightarrow\; q_A^* = \frac{157}{7} \approx 22{,}43 ,
\]
y $q_B^* = \frac{382}{21} \approx 18{,}19$. Los precios son $P_A^* \approx 28{,}17$ y $P_B^* \approx 30{,}29$, y los beneficios $\pi_A^* \approx 586{,}9$ y $\pi_B^* \approx 496{,}3$.

\end{solution}
\end{parts}

%-------------------------------------------------------------------- Comente: Stackelberg eficiente
\needspace{8\baselineskip}
\tema{Comente: eficiencia de Stackelberg vs Cournot}
\question \textbf{Comente.} ¿Cuándo la competencia tipo Stackelberg puede ser más eficiente que Cournot? ¿Cuándo podría ser más ineficiente que Cournot?
\needspace{5\baselineskip}
\begin{solution}
En caso de que tanto la firma líder como la seguidora sean igual de eficientes, el resultado Stackelberg será más eficiente que Cournot: el precio es menor y la cantidad producida es mayor.

En caso de que la líder sea más ineficiente, puede ocurrir que Stackelberg no sea más eficiente que Cournot, pues parte de la producción se traspasa a la firma más ineficiente (las firmas más ineficientes en Cournot producen menos).

\end{solution}

%-------------------------------------------------------------------- Comente: EN vs EPS
\needspace{8\baselineskip}
\tema{Comente: equilibrio de Nash y equilibrio perfecto en subjuegos}
\question \textbf{Comente.} Defina equilibrio de Nash y equilibrio perfecto en subjuegos. ¿Cómo se relacionan?
\needspace{5\baselineskip}
\begin{solution}
Un equilibrio de Nash (EN) es un resultado de un juego (conjunto de estrategias) tal que ningún individuo tiene incentivo a cambiar su estrategia. Un equilibrio perfecto en subjuegos (EPS) es un equilibrio de Nash creíble en un juego secuencial: es un EN en todos los subjuegos. Todo EPS es un EN, pero no todo EN es un EPS.

\end{solution}

%-------------------------------------------------------------------- Juego extensivo
\needspace{8\baselineskip}
\tema{Juego en forma extensiva: Nash y perfección en subjuegos}
\question Considere el siguiente juego secuencial. El jugador 1 elige entre $A$, $B$ y $C$. Si elige $C$, el juego termina con pagos $(2,2)$. Si elige $A$, el jugador 2 elige entre $a$ y $b$; si elige $b$ el juego termina con pagos $(1,-1)$, y si elige $a$ el jugador 1 elige entre $X$ (pagos $(3,0)$) e $Y$ (pagos $(0,3)$). Si elige $B$, el jugador 1 elige entre $\alpha$ y $\beta$; si elige $\beta$ el juego termina con pagos $(1,-1)$, y si elige $\alpha$ el jugador 2 elige entre $x$ (pagos $(3,0)$) e $y$ (pagos $(0,3)$). (Los pagos se escriben como $(\text{jugador 1},\text{jugador 2})$.)
\begin{center}
\begin{tikzpicture}[x=1cm,y=1cm,font=\small,every node/.style={inner sep=1pt},
  lev/.style={sloped,midway,above,font=\scriptsize}]
  \node (r) at (0,0) {1};
  \node (n2a) at (-3,-1.2) {2};
  \node (n1a) at (-4.5,-2.4) {1};
  \node (n1b) at (3,-1.2) {1};
  \node (n2b) at (1.8,-2.4) {2};
  \draw (r) -- (n2a) node[lev] {$A$};
  \draw (r) -- (0,-1.6) node[below,align=center] {$(2,2)$} node[midway,right] {$C$};
  \draw (r) -- (n1b) node[lev] {$B$};
  \draw (n2a) -- (n1a) node[lev] {$a$};
  \draw (n2a) -- (-2,-2.4) node[below,align=center] {$(1,-1)$} node[midway,right] {$b$};
  \draw (n1a) -- (-5.4,-3.6) node[below] {$(3,0)$} node[midway,left] {$X$};
  \draw (n1a) -- (-3.6,-3.6) node[below] {$(0,3)$} node[midway,right] {$Y$};
  \draw (n1b) -- (n2b) node[lev] {$\alpha$};
  \draw (n1b) -- (4.2,-2.4) node[below,align=center] {$(1,-1)$} node[midway,right] {$\beta$};
  \draw (n2b) -- (1,-3.6) node[below] {$(3,0)$} node[midway,left] {$x$};
  \draw (n2b) -- (2.6,-3.6) node[below] {$(0,3)$} node[midway,right] {$y$};
\end{tikzpicture}
\end{center}
\begin{parts}
\part Encuentre el equilibrio de Nash (perfecto en subjuegos).
\needspace{5\baselineskip}
\begin{solution}
Resolvemos por inducción hacia atrás, empezando desde abajo. En el nodo inferior izquierdo, el jugador 1 escoge $X$ (3 > 0). En el nodo inferior derecho, el jugador 2 escoge $y$ (3 > 0). En el nodo del medio izquierdo, el jugador 2 escoge $a$ (si elige $b$ obtiene $-1$, y con $a$ obtiene $0$). En el nodo del medio derecho, el jugador 1 escoge $\beta$ (obtiene 1, versus 0 si elige $\alpha$ y el jugador 2 responde $y$). Finalmente, en el nodo inicial, el jugador 1 compara: $A \to 3$, $B \to 1$, $C \to 2$, y escoge $A$.

El equilibrio de Nash perfecto en subjuegos es $(AX\beta,\ ay)$, con pagos $(3,0)$.

\end{solution}

\part Escriba el juego en forma normal.
\needspace{5\baselineskip}
\begin{solution}
Las estrategias del jugador 1 son combinaciones de su acción inicial, su acción tras $(A,a)$ y su acción tras $B$; las del jugador 2 son (acción tras $A$, acción tras $B\alpha$):
\begin{center}
\renewcommand{\arraystretch}{1.05}
\begin{tabular}{c|cccc}
  & $ax$ & $ay$ & $bx$ & $by$ \\ \hline
  $AX\alpha$ & $3,0$ & $3,0$ & $1,-1$ & $1,-1$ \\
  $AX\beta$ & $3,0$ & $3,0$ & $1,-1$ & $1,-1$ \\
  $AY\alpha$ & $0,3$ & $0,3$ & $1,-1$ & $1,-1$ \\
  $AY\beta$ & $0,3$ & $0,3$ & $1,-1$ & $1,-1$ \\
  $BX\alpha$ & $3,0$ & $0,3$ & $3,0$ & $0,3$ \\
  $BX\beta$ & $1,-1$ & $1,-1$ & $1,-1$ & $1,-1$ \\
  $BY\alpha$ & $3,0$ & $0,3$ & $3,0$ & $0,3$ \\
  $BY\beta$ & $1,-1$ & $1,-1$ & $1,-1$ & $1,-1$ \\
  $C\cdot\cdot$ (4 estrategias) & $2,2$ & $2,2$ & $2,2$ & $2,2$
\end{tabular}
\end{center}

\end{solution}

\part Encuentre un equilibrio de Nash en forma normal que lleve a un resultado diferente al de la solución en (a).
\needspace{5\baselineskip}
\begin{solution}
$(CX\alpha,\ by)$ es otro equilibrio de Nash, con pagos $(2,2)$. El jugador 2 no tiene incentivos a desviarse si el jugador 1 juega $C$ (el juego acaba ahí). Y el jugador 1 tampoco tiene incentivo para desviarse: si juega $A$, el otro le responde $b$ y gana 1; si escoge $B$, gana 0 o 1. Básicamente, en forma normal el jugador 1 escoge $C$ y gana 2, pero en forma secuencial escoge $A$, puesto que sabe que, dadas las estrategias creíbles, acabará ganando 3. Este equilibrio no es perfecto en subjuegos, ya que la respuesta $b$ del jugador 2 después de $A$ no es creíble.

\end{solution}
\end{parts}

%-------------------------------------------------------------------- Isobeneficio
\needspace{8\baselineskip}
\tema{Curvas de isobeneficio y Stackelberg}
\question La curva de demanda inversa de un bien está dada por la función $P(Q) = 240 - 2Q$. Existen dos firmas con funciones de costos $C_1(q_1) = 40q_1$ y $C_2(q_2) = 80q_2$.
\begin{parts}
\part Encuentre la expresión de la curva de isobeneficio de la firma 1. ¿Qué representa? Grafique.
\needspace{5\baselineskip}
\begin{solution}
La curva de isobeneficio de la firma 1 es el conjunto de combinaciones $(q_1, q_2)$ que le entregan un mismo nivel de beneficios $\bar\pi_1$. De manera genérica, con $P = A - \eta Q$:
\[
  \bar\pi_1 = \bigl(A - \eta(q_1 + q_2) - c_1\bigr)q_1
  \;\Longrightarrow\;
  q_2 = \frac{A - c_1}{\eta} - q_1 - \frac{\bar\pi_1}{\eta q_1}.
\]
Con los datos del problema ($A=240$, $\eta=2$, $c_1=40$): $q_2 = 100 - q_1 - \frac{\bar\pi_1}{2q_1}$. Es una curva con forma de ``joroba'': para cada $q_2$ la firma 1 prefiere estar en la curva más baja posible (menor $q_2$ implica mayor beneficio), y el máximo de cada curva se encuentra sobre su función de reacción $q_1^*(q_2)$. En el equilibrio de Cournot la curva de isobeneficio que pasa por el punto $(40,20)$ corresponde a $\bar\pi_1 = 3200$. En Stackelberg, la líder elige el punto de la función de reacción de la seguidora, $q_2 = 40 - \frac{q_1}{2}$, que se encuentra en la curva de isobeneficio más baja posible, tangente a ella: $(60,10)$, con $\bar\pi_1 = 3600$.
\begin{center}
\begin{tikzpicture}[x=0.075cm,y=0.1cm,font=\small,>=stealth]
  \draw[->] (0,0) -- (105,0) node[right] {$q_1$};
  \draw[->] (0,0) -- (0,52) node[above] {$q_2$};
  \draw[thick,blue!70,domain=20:80,samples=80,smooth] plot (\x,{100-\x-1600/\x});
  \draw[thick,violet,domain=23.6:76.4,samples=80,smooth] plot (\x,{100-\x-1800/\x});
  \draw[thick,gray] (24,52) -- (50,0);
  \draw[thick,orange!80!black] (0,40) -- (80,0);
  \fill (40,20) circle (2pt) node[above right,xshift=-3pt] {Cournot $(40,20)$};
  \fill[red] (60,10) circle (2pt) node[below left,red] {Stackelberg $(60,10)$};
  \node[blue!70] at (18,42) {$\bar\pi_1=3200$};
  \node[violet] at (90,7) {$\bar\pi_1=3600$};
  \node[orange!80!black,font=\scriptsize,anchor=west] at (62,31) {$q_2^*(q_1)$};
  \node[gray,font=\scriptsize,anchor=west] at (26,46) {$q_1^*(q_2)$};
\end{tikzpicture}
\end{center}
(La recta decreciente naranja es la función de reacción de la firma 2, $q_2 = 40 - \frac{q_1}{2}$; la recta gris, $q_2 = 100 - 2q_1$, es la función de reacción de la firma 1, $q_1 = 50 - \frac{q_2}{2}$.)

\end{solution}

\part Suponga que las dos empresas de la industria se comportan à la Stackelberg (en cantidad), de modo que la empresa 1 actúa como líder y la 2 como seguidora. ¿Cuál sería la producción y el precio de equilibrio de mercado? ¿Cuánto producirá la empresa 1? ¿Y la empresa 2? ¿Cuáles serán los beneficios de ambas empresas?
\needspace{5\baselineskip}
\begin{solution}
Se resuelve por inducción hacia atrás (ver el ejercicio anterior de Stackelberg con estos mismos datos): $q_1^* = 60$, $q_2^* = 10$, $Q^* = 70$, $P^* = 100$, $\pi_1 = 3600$ y $\pi_2 = 200$.

\end{solution}
\end{parts}

%-------------------------------------------------------------------- Stackelberg con inversión
\needspace{8\baselineskip}
\tema{Stackelberg con inversión para convertirse en líder}
\question Considere una industria con dos empresas, $L$ y $S$, que producen un producto homogéneo. La empresa $L$ es la líder y la $S$ la seguidora. La función inversa de demanda es $p(q_L,q_S) = 1 - q_L - q_S$, donde $q_L$ y $q_S$ son las cantidades producidas por la líder y la seguidora, respectivamente, y $p$ es el precio de mercado. Suponga además que la función de costos es $C_i(q_i) = \frac{1}{2}q_i$, para $i = L,S$.
\begin{parts}
\part Calcule las producciones en equilibrio, el precio de mercado, las cuotas de mercado y el beneficio de cada empresa.
\needspace{5\baselineskip}
\begin{solution}
Por inducción hacia atrás, la seguidora resuelve $\max_{q_S}\ (\frac{1}{2} - q_L - q_S)q_S$, de donde $q_S(q_L) = \frac{1}{4} - \frac{q_L}{2}$. La líder resuelve
\[
  \max_{q_L}\ \Bigl(\tfrac{1}{2} - q_L - \tfrac{1}{4} + \tfrac{q_L}{2}\Bigr)q_L = \Bigl(\tfrac{1}{4} - \tfrac{q_L}{2}\Bigr)q_L
  \;\Rightarrow\; q_L^* = \frac{1}{4}, \qquad q_S^* = \frac{1}{8}.
\]
Así, $Q = \frac{3}{8}$ y $p = \frac{5}{8}$. Las cuotas de mercado son $\frac{2}{3}$ para $L$ y $\frac{1}{3}$ para $S$, y los beneficios son $\pi_L = \frac{1}{8}\cdot\frac{1}{4} = \frac{1}{32}$ y $\pi_S = \frac{1}{8}\cdot\frac{1}{8} = \frac{1}{64}$.

\end{solution}

\part Suponga que las empresas compiten en cantidades simultáneamente (Cournot). Calcule las cantidades de equilibrio, el precio, las cuotas y los beneficios, y compare su respuesta con la anterior.
\needspace{5\baselineskip}
\begin{solution}
Cada firma resuelve $\max_{q_i}\ (\frac{1}{2} - q_i - q_j)q_i$, de modo que $q_i = \frac{1/2 - q_j}{2}$ y $q_L^* = q_S^* = \frac{1}{6}$. Entonces $Q = \frac{1}{3}$, $p = \frac{2}{3}$, las cuotas son $\frac{1}{2}$ cada una y $\pi_L = \pi_S = \frac{1}{6}\cdot\frac{1}{6} = \frac{1}{36}$.

Comparando con Stackelberg: la cantidad total es mayor y el precio menor ($\frac{5}{8} < \frac{2}{3}$); la líder gana más ($\frac{1}{32} > \frac{1}{36}$) y la seguidora gana menos ($\frac{1}{64} < \frac{1}{36}$) que en Cournot.

\end{solution}

\part Suponga que las empresas compiten en cantidades secuencialmente, pero ahora la empresa $S$ puede invertir $\bar S$ y convertirse en líder. ¿Debería la seguidora invertir? ¿Cómo depende su decisión del valor de $\bar S$?
\needspace{5\baselineskip}
\begin{solution}
Si $S$ invierte, se convierte en líder y obtiene el beneficio de la líder menos la inversión, $\frac{1}{32} - \bar S$. Si no invierte, queda como seguidora y obtiene $\frac{1}{64}$. Por lo tanto, $S$ debe invertir si y solo si
\[
  \frac{1}{32} - \bar S > \frac{1}{64}
  \quad\Longleftrightarrow\quad
  \bar S < \frac{1}{64}.
\]
Es decir, la seguidora invierte solo si el costo de la inversión es menor que la ganancia de ser líder en lugar de seguidora.

\end{solution}
\end{parts}

%-------------------------------------------------------------------- Juego de inversión del monopolista
\needspace{8\baselineskip}
\tema{Juego secuencial: inversión del monopolista y entrada}
\question Considere un juego secuencial donde un monopolista puede invertir en tecnología antes de la potencial entrada de un competidor. Si el monopolista invierte, la estructura de pagos favorece la estrategia de guerra en caso de entrada del competidor. El monopolista elige primero entre invertir ($I$) o no invertir ($NI$). Luego el competidor debe decidir entre entrar ($E$) o no entrar ($NE$) al mercado, y si entra el monopolista elige entre acomodar ($A$) o iniciar una guerra de precios ($G$). Los pagos (monopolista, entrante) son:
\begin{center}
\renewcommand{\arraystretch}{1.1}
\begin{tabular}{llc}
  Invierte & Entra y guerra ($I,E,G$) & $(10,-10)$ \\
  Invierte & Entra y acomoda ($I,E,A$) & $(0,20)$ \\
  Invierte & No entra ($I,NE$) & $(30,0)$ \\
  No invierte & No entra ($NI,NE$) & $(50,0)$ \\
  No invierte & Entra y guerra ($NI,E,G$) & $(-10,-10)$ \\
  No invierte & Entra y acomoda ($NI,E,A$) & $(20,20)$
\end{tabular}
\end{center}
\begin{parts}
\part Obtenga el equilibrio perfecto en subgrupos (subjuegos) del anterior juego secuencial.
\needspace{5\baselineskip}
\begin{solution}
Resolvemos desde abajo hacia arriba. El último en elegir es el monopolista, que decide entre guerra ($G$) o acomodar ($A$) en el caso de que el competidor decida entrar. En el lado izquierdo, cuando el monopolista invirtió, optará por una guerra de precios, pues obtiene mayores beneficios ($10 > 0$). De la misma forma, en el lado derecho, cuando el monopolista no invierte, preferirá acomodar al competidor ($20 > -10$).

Luego, el competidor decide entre entrar ($E$) o no entrar ($NE$), sabiendo lo que hará el monopolista. Si el monopolista invirtió, el competidor decidirá no entrar ($0 > -10$). Si el monopolista decidió no invertir, el competidor elegirá entrar ($20 > 0$).

Finalmente, el monopolista, sabiendo todo lo anterior, sabe que si invierte tendrá un pago de $30$, mientras que al no invertir tendrá un pago de $20$; por lo tanto, decide realizar la inversión, ya que esta (en el sentido económico) amenaza al entrante, persuadiéndolo de no entrar, aunque con un costo asociado: hubiera obtenido un mejor pago ($50$) si no hubiera invertido y el competidor no entrara. El equilibrio perfecto en subjuegos es $\{(I,G,A),\ (NE,E)\}$, con pagos $(30,0)$.

\end{solution}

\part Escriba el juego secuencial en su forma normal.
\needspace{5\baselineskip}
\begin{solution}
Las estrategias del monopolista son (decisión de invertir, acción si hay entrada tras $I$, acción si hay entrada tras $NI$); las del entrante son (acción tras $I$, acción tras $NI$):
\begin{center}
\renewcommand{\arraystretch}{1.1}
\begin{tabular}{l|cccc}
  & $E,NE$ & $E,E$ & $NE,NE$ & $NE,E$ \\ \hline
  $I,G,G$ & $10,-10$ & $10,-10$ & $30,0$ & $30,0$ \\
  $I,G,A$ & $10,-10$ & $10,-10$ & $30,0$ & $30,0$ \\
  $I,A,G$ & $0,20$ & $0,20$ & $30,0$ & $30,0$ \\
  $I,A,A$ & $0,20$ & $0,20$ & $30,0$ & $30,0$ \\
  $NI,G,G$ & $50,0$ & $-10,-10$ & $50,0$ & $-10,-10$ \\
  $NI,G,A$ & $50,0$ & $20,20$ & $50,0$ & $20,20$ \\
  $NI,A,G$ & $50,0$ & $-10,-10$ & $50,0$ & $-10,-10$ \\
  $NI,A,A$ & $50,0$ & $20,20$ & $50,0$ & $20,20$
\end{tabular}
\end{center}

\end{solution}

\part Encuentre un equilibrio de Nash en forma normal que lleve a un resultado diferente al de la solución en (a).
\needspace{5\baselineskip}
\begin{solution}
En este juego simultáneo aparecen distintos equilibrios de Nash, como el par de estrategias $(NI,G,G)$ y $(E,NE)$, con pagos $(50,0)$: el monopolista no invierte y amenaza con guerra si hay entrada, y el entrante no entra tras $NI$ porque anticipa la guerra (el monopolista no tiene incentivos a cambiar, pues obtiene el máximo $50$, y el entrante tampoco). Esto se debe a que en la forma normal las decisiones se toman al mismo tiempo y se pierde la información secuencial crucial del juego: la amenaza de guerra tras $NI$ no es creíble en el juego secuencial (acomodar da $20 > -10$), pero resulta ser un equilibrio de Nash válido en la forma normal. Esto ilustra por qué fue desarrollado el concepto de equilibrio perfecto en subjuegos: para eliminar equilibrios basados en amenazas no creíbles que, aunque matemáticamente consistentes en la forma normal, no reflejan el comportamiento racional esperado cuando las decisiones se toman secuencialmente.

\end{solution}
\end{parts}

%-------------------------------------------------------------------- Stackelberg con costo fijo
\needspace{8\baselineskip}
\tema{Cournot y Stackelberg con costos asimétricos y costo fijo (líder 1 o 2)}
\question Asuma que existen 2 firmas que producen un bien homogéneo. Sus funciones de costos vienen dadas por $C_1(q_1) = 20q_1 + 1$ y $C_2(q_2) = 30q_2$. La curva de demanda inversa en el mercado es $P = 60 - 5Q$.
\begin{parts}
\part Si ambas firmas compiten à la Cournot, ¿cuál es el precio de equilibrio? ¿Cuál es el nivel de producción y beneficios de cada firma?
\needspace{5\baselineskip}
\begin{solution}
Firma 1: $\max_{q_1}\ (60 - 5q_1 - 5q_2 - 20)q_1 - 1$, con CPO $40 - 10q_1 - 5q_2 = 0$, es decir, $q_1 = 4 - \frac{q_2}{2}$. Firma 2: $\max_{q_2}\ (60 - 5q_1 - 5q_2 - 30)q_2$, con CPO $30 - 5q_1 - 10q_2 = 0$, es decir, $q_2 = 3 - \frac{q_1}{2}$. Reemplazando: $q_1 = 4 - \frac{3 - q_1/2}{2}$, de donde $\frac{3}{4}q_1 = \frac{5}{2}$ y
\[
  q_1^C = \frac{10}{3}, \qquad q_2^C = \frac{4}{3}, \qquad Q = \frac{14}{3}, \qquad P^C = 60 - 5\cdot\frac{14}{3} = \frac{110}{3},
\]
\[
  \pi_1^C = \Bigl(\frac{110}{3} - 20\Bigr)\frac{10}{3} - 1 = \frac{491}{9} \approx 54{,}56, \qquad
  \pi_2^C = \Bigl(\frac{110}{3} - 30\Bigr)\frac{4}{3} = \frac{80}{9} \approx 8{,}89 .
\]

\end{solution}

\part Suponga que ambas firmas se comportan à la Stackelberg (en cantidad), de modo que la empresa 1 actúa como líder y la 2 como seguidora. ¿Cuánto producirán las firmas 1 y 2? Calcule el nuevo precio de equilibrio y los beneficios de cada firma.
\needspace{5\baselineskip}
\begin{solution}
Por inducción hacia atrás, en $t=2$ la seguidora tiene $q_2(q_1) = 3 - \frac{q_1}{2}$. La líder maximiza
\[
\begin{gathered}
  \max_{q_1}\ \Bigl(60 - 5q_1 - 5\bigl(3 - \tfrac{q_1}{2}\bigr) - 20\Bigr)q_1 - 1 = (25 - 2{,}5q_1)q_1 - 1 \\
  \Longrightarrow\; 25 - 5q_1 = 0 \;\Longrightarrow\; q_1 = 5 .
\end{gathered}
\]
Por lo tanto $q_2 = 3 - \frac{5}{2} = \frac{1}{2}$, $Q = \frac{11}{2}$ y $P = 60 - 5\cdot\frac{11}{2} = 32{,}5$. Los beneficios son
\[
  \pi_1 = (32{,}5 - 20)\cdot 5 - 1 = 61{,}5, \qquad \pi_2 = (32{,}5 - 30)\cdot\frac{1}{2} = 1{,}25 .
\]

\end{solution}

\part ¿Cómo cambia su respuesta en (b) si la firma 2 es la líder?
\needspace{5\baselineskip}
\begin{solution}
Ahora la firma 1 es seguidora, con $q_1(q_2) = 4 - \frac{q_2}{2}$. La líder 2 maximiza
\[
  \max_{q_2}\ \Bigl(60 - 5\bigl(4 - \tfrac{q_2}{2}\bigr) - 5q_2 - 30\Bigr)q_2 = (10 - 2{,}5q_2)q_2
  \;\Longrightarrow\; q_2 = 2 .
\]
Entonces $q_1 = 3$, $Q = 5$ y $P = 60 - 25 = 35$. Los beneficios son $\pi_1 = (35 - 20)\cdot 3 - 1 = 44$ (seguidora) y $\pi_2 = (35 - 30)\cdot 2 = 10$ (líder). Dada la asimetría de costos, cuando la firma 2 (de mayor costo) actúa como líder se observa una menor oferta en el mercado y un precio más alto; la firma 2 produce menos y obtiene menores beneficios en comparación con la firma 1, pero estos beneficios son superiores a los que obtenía cuando actuaba como seguidora.

\end{solution}

\part ¿Qué tipo de competencia le conviene a las firmas 1 y 2?
\needspace{5\baselineskip}
\begin{solution}
Comparamos los beneficios en cada caso. Para la firma 1: $\pi_1^{\text{líder}} = 61{,}5 > \pi_1^{C} = 54{,}56 > \pi_1^{\text{seguidora}} = 44$. Para la firma 2: $\pi_2^{\text{líder}} = 10 > \pi_2^{C} = 8{,}89 > \pi_2^{\text{seguidora}} = 1{,}25$. Por lo tanto, a ambas firmas les conviene competir à la Stackelberg actuando como líderes; si actúan como seguidoras, les resulta más beneficioso competir à la Cournot.

\end{solution}
\end{parts}

\end{unidad}

\unidadtitulo{Colusión y juegos repetidos}
\begin{unidad}
%-------------------------------------------------------------------- Comente: transparencia
\needspace{8\baselineskip}
\tema{Comente: transparencia y colusión}
\question \textbf{Comente.} La Fiscalía Nacional Económica quiere aumentar al máximo la transparencia entre empresas dentro del mercado para así facilitar su tarea de fiscalización y evitar posibles colusiones.
\needspace{5\baselineskip}
\begin{solution}
Falso. Ante un mayor nivel de transparencia, la colusión se vuelve más fácil entre firmas, ya que cualquier desvío del acuerdo colusivo será detectado más fácilmente, haciendo que el castigo sea instantáneo. Así disminuyen los incentivos a desviarse (es mayor el costo esperado del desvío). Si, por el contrario, existe información que posee solo la propia firma, el desvío es mucho más difícil de detectar y por ende también lo es la aplicación del castigo por parte del resto de los participantes del acuerdo colusivo, de modo que hay mayor incentivo a dejar de cooperar. Por lo tanto, ante menor transparencia en el mercado y menos acceso a las decisiones del rival, más difícil es la colusión.

\end{solution}

%-------------------------------------------------------------------- Colusión Bertrand y Cournot
\needspace{8\baselineskip}
\tema{Colusión en Bertrand y Cournot (estrategia gatillo)}
\question Suponga que en China, debido al coronavirus, solo han quedado dos empresas que comercian animales exóticos para consumo. Ambas empresas tienen los mismos costos marginales, iguales a $c$, y venden un producto homogéneo. La demanda inversa de mercado que enfrentan está dada por $P = A - Q$. Las firmas están estudiando la posibilidad de coludirse en diferentes escenarios; consideremos que las firmas descuentan los beneficios futuros a un factor $\delta$ y que ante un desvío aplican la estrategia gatillo.
\begin{parts}
\part Derive la condición que debe cumplir $\delta$ para que la colusión sea sostenible y encuentre el valor del factor $\delta_C$ que hace posible la colusión si estas firmas compiten en cantidades, y el factor $\delta_B$ que hace posible la colusión si estas firmas compiten en precios. Considere que al coludirse se reparten los beneficios equitativamente.
\needspace{5\baselineskip}
\begin{solution}
Para que una estrategia colusiva sea sostenible, el valor presente de colusionar debe ser mayor o igual al de desviarse y competir:
\[
  VP(\text{colusión}) \geq VP(\text{desvío})
  \quad\Longleftrightarrow\quad
  \sum_{t=0}^{\infty}\delta^t \pi^C \geq \pi^D + \sum_{t=1}^{\infty}\delta^t \pi^N .
\]
Aplicando la suma geométrica:
\[
  \frac{\pi^C}{1-\delta} \geq \pi^D + \frac{\delta\,\pi^N}{1-\delta}
  \;\Longrightarrow\;
  \pi^C \geq \pi^D(1-\delta) + \delta\pi^N
  \;\Longrightarrow\;
  \delta \geq \frac{\pi^D - \pi^C}{\pi^D - \pi^N}.
\]
Falta calcular los pagos en cada caso.

\textbf{Bertrand.} Los beneficios de coludirse son los beneficios monopólicos repartidos entre ambas firmas, $\pi^C = \frac{\pi^M}{2} = \frac{(A-c)^2}{8}$. Al desviarse (bajando levemente el precio) la firma obtiene los beneficios monopólicos, $\pi^D = \pi^M = \frac{(A-c)^2}{4}$, y los beneficios de no cooperar son $\pi^N = 0$. Luego
\[
  \delta_B \geq \frac{\frac{(A-c)^2}{4} - \frac{(A-c)^2}{8}}{\frac{(A-c)^2}{4} - 0} = \frac{1}{2}.
\]

\textbf{Cournot.} Los beneficios colusivos son iguales al caso de precios, $\pi^C = \frac{(A-c)^2}{8}$ (cada firma produce $q^C = \frac{A-c}{4}$). En la competencia, la función de reacción es $q_i = \frac{A - q_j - c}{2}$, de modo que cada firma produce $q^N = \frac{A-c}{3}$ y
\[
  \pi^N = \frac{(A-c)^2}{9}.
\]
Para calcular el beneficio de desvío, reemplazamos la cantidad colusiva de la firma $i$ en la función de reacción de la firma $j$:
\[
\begin{gathered}
  q^D = \frac{A - \frac{A-c}{4} - c}{2} = \frac{3(A-c)}{8} \\
  \Longrightarrow\; \pi^D = \frac{9(A-c)^2}{64}.
\end{gathered}
\]
Finalmente,
\[
  \delta_C \geq \frac{\frac{9}{64} - \frac{1}{8}}{\frac{9}{64} - \frac{1}{9}} = \frac{9}{17} \approx 0{,}53 .
\]

\end{solution}

\part ¿Bajo qué tipo de competencia es más factible la colusión?
\needspace{5\baselineskip}
\begin{solution}
De la parte anterior se tiene que $\delta_C > \delta_B$, por lo tanto es más fácil coludirse en Bertrand, debido a que el castigo (competir à la Bertrand, con beneficios nulos) es más severo que en Cournot.

\end{solution}
\end{parts}

%-------------------------------------------------------------------- Colusión asimétrica
\needspace{8\baselineskip}
\tema{Colusión con firmas asimétricas}
\question Un mercado posee una demanda $Q(P) = 36 - P$. Existen dos empresas que compiten en él mediante precios. La primera tiene costo marginal $c_1 = 0$, mientras que la segunda tiene costo $c_2 = 4$.
\begin{parts}
\part Suponga que las empresas desean coludirse. ¿Cuál será el precio de colusión que escogerían y por qué?
\needspace{5\baselineskip}
\begin{solution}
Sin restricciones de capacidad, todo se produce al costo de la más eficiente. Por ende, maximizamos el beneficio conjunto:
\[
  \max_p\ \pi = p(36 - p)
  \;\Rightarrow\; 36 - 2p = 0
  \;\Rightarrow\; p^C = 18, \qquad \pi = 324 .
\]

\end{solution}

\part ¿Cuál es el máximo reparto del mercado $S_2$ (\%) que podría llevarse la firma 2 para que el acuerdo sea factible, si el factor de descuento intertemporal es $\delta = 0{,}75$?
\needspace{5\baselineskip}
\begin{solution}
En caso de competencia, $P = 4 - \epsilon \approx 4$, $Q = 32$, $\pi_1^N = 128$ y $\pi_2^N = 0$. Si la firma 1 se desvía, gana beneficios monopólicos $\pi^M = 324$. Para que el acuerdo sea sostenible, ambas firmas deben aceptarlo; por lo tanto, realizamos el proceso para cada una.

\textbf{Firma 1} (con participación $S_1$):
\[
\begin{gathered}
  \delta \geq \frac{\pi^D - \pi^C}{\pi^D - \pi^N} \\
  \;\Longrightarrow\; \\
  0{,}75 \geq \frac{324 - 324S_1}{324 - 128} \\
  \;\Longrightarrow\; \\
  147 \geq 324 - 324S_1 \\
  \;\Longrightarrow\; \\
  S_1 \geq 0{,}546 .
\end{gathered}
\]
Como $S_2 = 1 - S_1$, se obtiene $S_2 \leq 0{,}454$, es decir, la firma 2 puede llevarse como máximo $\approx 45{,}4\%$ del mercado. (La firma 2, por su parte, solo exige $S_2 \geq 0{,}194$, pues si se desvía vendiendo al precio $18-\epsilon$ obtiene $252$ y luego beneficios nulos.)

\end{solution}

\part Obtenga las condiciones para que el acuerdo colusivo sea sostenible si las empresas deciden turnarse la producción. Es decir, un período solo produce una de ellas y en el siguiente período produce la otra, y así sucesivamente. \emph{(Asuma que en $t=0$ la firma 1 parte produciendo, por lo que la firma 2 se puede desviar desde un inicio y vender en el próximo período.)}
\needspace{5\baselineskip}
\begin{solution}
Bajo este acuerdo, cuando le toca fijar el precio a cada firma, fija su propio precio monopólico. Para la firma 1: $p^M = 18$, $Q=18$, $\pi = 324$. Para la firma 2: $p^M = 20$, $Q = 16$, $\pi = 256$. El equilibrio no cooperativo es $p = 4-\epsilon$, $Q = 32$, $\pi_1 = 128$, $\pi_2 = 0$.

El desvío consiste en producir cuando no le toca. Para la firma 1, fijar $p = 18 < p_2^M = 20$ le da $\pi = 324$. Para la firma 2, fijar $p = 18 - \epsilon < p_1^M$ le da $Q = 18$ y $\pi_2 = 252$.

\textbf{Firma 1:} $VP_1(\text{cooperar}) \geq VP_1(\text{desvío})$,
\[
\begin{gathered}
  \sum_{t=0}^{\infty}\delta_1^{2t}\,324 \geq 324(1+\delta_1) + \sum_{t=2}^{\infty}\delta_1^{t}\,128 \\
  \;\Longrightarrow\; \\
  \frac{324}{1-\delta_1^2} \geq 324(1+\delta_1) + \frac{128\,\delta_1^2}{1-\delta_1} \\
  \;\Longrightarrow\; \\
  196\delta_1^2 + 196\delta_1 - 324 \geq 0,
\end{gathered}
\]
de donde $\delta_1 \geq 0{,}8795$.

\textbf{Firma 2:} $VP_2(\text{cooperar}) \geq VP_2(\text{desvío})$,
\[
  \sum_{t=0}^{\infty}\delta_2^{2t+1}\,256 \geq 252
  \;\Longrightarrow\;
  \frac{256\,\delta_2}{1-\delta_2^2} \geq 252
  \;\Longrightarrow\;
  252\delta_2^2 + 256\delta_2 - 252 \geq 0,
\]
de donde $\delta_2 \geq 0{,}6137$. (A la firma 2 le conviene desviarse solo en el período en que produce la otra firma.)

Con $\delta = 0{,}75$ para ambas firmas, a la firma 2 le conviene cooperar pero a la firma 1 no, por lo que el acuerdo de turnos no es sostenible (se requiere $\delta \geq 0{,}8795$).

\end{solution}
\end{parts}

%-------------------------------------------------------------------- Comente: frecuencia
\needspace{8\baselineskip}
\tema{Comente: frecuencia de interacción y colusión}
\question \textbf{Comente.} Mantener una colusión estable será más fácil a menor frecuencia de operaciones (frecuencia con que las firmas interactúan y fijan los precios).
\needspace{5\baselineskip}
\begin{solution}
Falso. Mientras menor sea la frecuencia con que las firmas se reúnen y fijan los precios, menor será la estabilidad del acuerdo colusivo. Si la frecuencia con la que las firmas se reúnen es menor, cualquier desvío del acuerdo tarda más en ser descubierto y, por tanto, los castigos demoran más en ser aplicados, lo que favorece el beneficio de desviarse.

\end{solution}

%-------------------------------------------------------------------- Comente: asimetría de costos
\needspace{8\baselineskip}
\tema{Comente: asimetría de costos y estabilidad de la colusión}
\question \textbf{Comente.} Asuma que las firmas compiten en precios. Se puede afirmar que, a mayor asimetría entre firmas (en términos de costos), menor estabilidad del acuerdo colusivo.
\needspace{5\baselineskip}
\begin{solution}
Verdadero. Si las firmas son muy distintas, es poco probable que logren mantener un acuerdo colusivo. La razón es que la firma más eficiente tiene beneficios positivos aun cuando no se coluda con sus rivales, por lo que estará dispuesta a ser parte del acuerdo solo si se le concede una participación de mercado relativamente alta (mayor a $0{,}5$). El problema es que, si se le concede una porción del mercado demasiado alta a la firma más eficiente, entonces será la firma más ineficiente la que no tenga incentivos a formar parte del acuerdo, ya que tendrá una participación de mercado tan baja que le convendrá más desviarse y asumir el posterior castigo que cooperar.

\end{solution}

%-------------------------------------------------------------------- Comente: garantía de precios bajos
\needspace{8\baselineskip}
\tema{Comente: garantía de precios bajos}
\question \textbf{Comente.} Lea el siguiente texto y comente.

\emph{«Garantía de precios bajos: Precios bajos todos los días. ¡Si encuentras un precio más bajo, lo igualamos y te damos un 20\% de descuento sobre el precio igualado! ¿Precio más bajo en otro lugar? ¡Imposible! *Recuerda: la cotización que entregues para la garantía de precios debe ser de un competidor de la misma localidad».}

Lo anterior corresponde a una estrategia de precios de una empresa en Chile. Analice si dicha práctica podría afectar la sostenibilidad de un eventual acuerdo colusivo entre dicha firma y sus competidores. Justifique su respuesta.
\needspace{5\baselineskip}
\begin{solution}
Esta práctica efectivamente afecta la sostenibilidad de un acuerdo colusivo, y lo hace robusteciéndola por dos vías:
\begin{itemize}
  \item Hace partícipes a los consumidores del monitoreo de los precios de la firma competidora, facilitando su revisión y aumentando con ello la transparencia del mercado.
  \item Fortalece la estrategia de castigo y hace menos beneficiosa la estrategia de desvío, ya que ante cualquier desvío del precio acordado la otra firma reaccionaría inmediatamente y aplicaría una disminución del precio de un 20\% por debajo del precio de desvío.
\end{itemize}

\end{solution}

%-------------------------------------------------------------------- Colusión en mercados en expansión y declive
\needspace{8\baselineskip}
\tema{Colusión en mercados en expansión y declive}
\question Suponga un mercado donde $n$ empresas simétricas compiten en precios con productos homogéneos. Suponga que las ganancias monopólicas crecen cada período a una tasa $g$, y que la tasa de descuento de cada firma es $\rho$. Derive la condición para $\rho$ para que la colusión sea sostenible en esta industria. Comente el efecto que tiene que sea una industria en expansión ($g>0$) o en declive ($g<0$) para la sostenibilidad de la colusión.
\needspace{5\baselineskip}
\begin{solution}
Bajo la estrategia de colusión, los pagos se incrementan en cada período en un factor $1+g$, por lo tanto el valor presente de seguir colusionando es
\[
  VP(\text{colusión}) = \frac{\pi^M}{n} + \delta\frac{(1+g)\pi^M}{n} + \delta^2\frac{(1+g)^2\pi^M}{n} + \dots = \frac{\pi^M}{n\bigl(1 - \delta(1+g)\bigr)} .
\]
El pago por desviarse sigue siendo el pago monopólico por una vez (no depende de $g$) y el pago en la etapa de castigo es cero. Por lo tanto, la condición de sostenibilidad es
\[
  \frac{\pi^M}{n\bigl(1-\delta(1+g)\bigr)} \geq \pi^M
  \;\Longrightarrow\;
  1 \geq n\bigl(1-\delta(1+g)\bigr)
  \;\Longrightarrow\;
  \delta(1+g) \geq 1 - \frac{1}{n}.
\]
Escribiendo $\delta = \frac{1}{1+\rho}$, con $\rho \in [0,+\infty)$ una tasa de impaciencia, la condición queda
\[
  \frac{1+g}{1+\rho} \geq 1 - \frac{1}{n}
  \quad\Longleftrightarrow\quad
  \rho \leq \frac{1+g}{1-\frac{1}{n}} - 1,
  \qquad \text{o bien} \qquad
  \delta \geq \frac{1-\frac{1}{n}}{1+g}.
\]
El $\delta$ mínimo necesario, $\frac{1-1/n}{1+g}$, es decreciente en $g$. Por lo tanto, si $g<0$ (industria en declive) se hace más difícil sostener la colusión: la tasa de impaciencia máxima con la que se sostiene es menor. Esto refleja que, si el crecimiento es negativo, los pagos futuros de seguir coludidos son menos atractivos y la tentación del desvío es mayor. Lo opuesto ocurre si $g>0$.

\end{solution}

%-------------------------------------------------------------------- Comente: horizonte finito
\needspace{8\baselineskip}
\tema{Comente: colusión en Cournot con horizonte finito}
\question \textbf{Comente.} En un duopolio con competencia à la Cournot con un período finito, las empresas producen una cantidad equitativa para repartirse así las utilidades de monopolio. Dado que las utilidades de monopolio son las máximas posibles, las empresas decidirán mantener el acuerdo incluso si no se fuerza legalmente o de alguna otra manera.
\needspace{5\baselineskip}
\begin{solution}
Falso. Sujeto a que la otra firma produce la mitad de lo que haría un monopolio, la estrategia óptima desde el punto de vista individual de la otra será producir más que la mitad de un monopolio. De esta manera, el acuerdo no se mantiene sin un mecanismo de coacción. Esto es válido para cualquier horizonte finito (por inducción hacia atrás, en el último período se juega el equilibrio de Cournot).

\end{solution}

%-------------------------------------------------------------------- Comente: Folk
\needspace{8\baselineskip}
\tema{Comente: teorema de Folk}
\question \textbf{Comente.} ¿Hay alguna manera de que estas empresas decidan cooperar? Explique en qué consiste el teorema de Folk.
\needspace{5\baselineskip}
\begin{solution}
El teorema de Folk nos indica que existe algún nivel de impaciencia (descuento) tal que un equilibrio cooperativo es factible en un juego repetido de horizonte infinito. Es decir, existe $0 < \delta < 1$ tal que la utilidad de la cooperación es mayor a la utilidad de no cooperar, donde la utilidad es
\[
  U_i = \sum_{t \geq 0}\delta^t u_i(x_t).
\]

\end{solution}

%-------------------------------------------------------------------- Comente: colusión mejor estado
\needspace{8\baselineskip}
\tema{Comente: la colusión como mejor estado de la industria}
\question \textbf{Comente.} La colusión, sin lugar a dudas, es el mejor estado que puede tener una industria, dado que todos sus miembros obtienen el ingreso monopólico.
\needspace{5\baselineskip}
\begin{solution}
Falso. Cuando los miembros de una industria se coluden, cada uno obtiene solo una fracción del ingreso monopólico, en función de la cantidad monopólica que producen o del poder de negociación de cada una de las firmas (y con asimetrías, la repartición puede dejar a alguna firma peor que compitiendo).

\end{solution}

%-------------------------------------------------------------------- Comente: licitaciones
\needspace{8\baselineskip}
\tema{Comente: colusión en licitaciones de obras públicas}
\question \textbf{Comente.} En las licitaciones de obras públicas, es muy probable que las firmas logren sostener un acuerdo colusivo debido a que para ambas es muy conveniente dejar de competir por quién ofrece el precio más bajo.
\needspace{5\baselineskip}
\begin{solution}
Falso. El hecho de que las firmas no se encuentren todos los períodos en el mercado induce a que tengan más incentivos a desviarse de la colusión, debido a que el castigo demorará más en llegar y por lo tanto podrán percibir durante todo el período que dure la licitación las ganancias del desvío. Se debe considerar además el papel de la fiscalización ante posibles colusiones, el tipo de competencia (en este caso, en precios) y el posible contacto multimercado.

\end{solution}

%-------------------------------------------------------------------- Comente: precios vs cantidades colusión
\needspace{8\baselineskip}
\tema{Comente: colusión en precios o cantidades con $N$ firmas}
\question \textbf{Comente.} En un mercado con 2 firmas simétricas, ¿es más fácil mantener un acuerdo colusivo si compiten en precios o en cantidades? ¿Y en un mercado con $N$ firmas simétricas? Explique los motivos.
\needspace{5\baselineskip}
\begin{solution}
En un mercado de 2 firmas es más fácil mantener un acuerdo colusorio si compiten en precios. Esto se debe a que el castigo por desviarse es mucho mayor en la competencia por precios que en la competencia por cantidades: cuando compiten à la Bertrand las firmas obtienen beneficios iguales a cero. En Bertrand los efectos pro-colusivos son superiores a los anti-colusivos.

Cuando crece el número de firmas (hasta $N$), tanto en Cournot como en Bertrand se reparten los beneficios de la colusión entre más firmas, por lo que el efecto pro-colusivo se debilita, pues la diferencia entre coludirse y competir disminuye. El efecto anti-colusivo de Bertrand aumenta más que el de Cournot, pues en caso de desviarse del acuerdo una firma puede quedarse con todo el mercado y obtener $N$ veces el beneficio de coludirse. Esto hace que para $N$ firmas la colusión sea más sostenible en competencia en cantidades.

\end{solution}

%-------------------------------------------------------------------- Colusión asimétrica 24-P
\needspace{8\baselineskip}
\tema{Colusión asimétrica en precios y turnos (variante numérica)}
\question Un mercado posee una demanda $Q(P) = 24 - P$. Existen dos empresas que compiten en él mediante precios. La primera tiene costo marginal $c_1 = 0$, mientras que la segunda tiene costo $c_2 = 3$.
\begin{parts}
\part Suponga que las empresas desean coludirse. ¿Cuál será el precio de colusión que escogerían y por qué?
\needspace{5\baselineskip}
\begin{solution}
Sin restricciones de capacidad, todo se produce al costo de la firma más eficiente, por ende el precio será el precio monopólico de la firma 1. Maximizando el beneficio conjunto:
\[
  \max_P\ \pi^C = P(24 - P) \;\Rightarrow\; 24 - 2P = 0 \;\Rightarrow\; P = 12, \qquad \pi^C = 144 .
\]

\end{solution}

\part ¿Cuál es el máximo reparto del mercado $S_2$ (\%) que podría llevarse la firma 2 para que el acuerdo sea factible, si el factor de descuento es $\delta = 0{,}7$?
\needspace{5\baselineskip}
\begin{solution}
Se debe cumplir $\delta \geq \frac{\pi^D - \pi^C}{\pi^D - \pi^N}$. Los pagos son:
\begin{itemize}
  \item \emph{Competencia:} $P = 3 - \epsilon$, $Q = 21$, $\pi_1^N = 63$, $\pi_2^N = 0$.
  \item \emph{Desvío de la firma 1:} $P = 12 - \epsilon$, $Q = 12$, $\pi_1^D = 144$.
  \item \emph{Colusión:} $P^M = 12$, $Q^M = 12$, $\pi^C = 144$ (la firma 1 recibe $144S_1$).
\end{itemize}
Para la firma 1:
\[
  0{,}7 \geq \frac{144 - 144S_1}{144 - 63}
  \;\Longrightarrow\;
  144S_1 \geq 87{,}3
  \;\Longrightarrow\;
  S_1 \geq 0{,}606 .
\]
Como $S_1 + S_2 = 1$, el máximo reparto que se puede llevar la firma 2 es $S_2 \leq 0{,}394$, es decir, $\approx 39\%$. (La firma 2, que si se desvía obtiene $108$, solo exige $S_2 \geq 0{,}225$.)

\end{solution}

\part Obtenga las condiciones para que el acuerdo colusivo sea sostenible si las empresas deciden turnarse la producción. Es decir, un período solo produce una de ellas y en el siguiente período produce la otra, y así sucesivamente.
\needspace{5\baselineskip}
\begin{solution}
Primero, calculamos los beneficios de competir, coludirse y desviarse.

\emph{Competencia:} $P = 3 - \epsilon$, $Q = 21$, $\pi_1^N = 63$, $\pi_2^N = 0$.

\emph{Colusión (beneficio monopólico de cada firma en su turno):} firma 1: $P^M = 12$, $Q^M = 12$, $\pi_1^C = 144$. Firma 2: $\max (P-3)(24-P) \Rightarrow P^M = 13{,}5$, $Q^M = 10{,}5$, $\pi_2^C = 110{,}25$.

\emph{Desvío:} la firma 1 (que en el turno de la 2 fija $P=12 < 13{,}5$) obtiene $\pi_1^D = 144$. La firma 2 (que en el turno de la 1 fija $P = 12 - \epsilon$) obtiene $\pi_2^D = 108$.

\textbf{Si comienza la firma 1:}
\begin{align*}
  \text{Firma 1:}\quad & \sum_{i=0}^{\infty}144\,\delta^{2i} \geq 144 + 144\delta + \sum_{i=2}^{\infty}63\,\delta^{i}, \\
  \text{Firma 2:}\quad & \sum_{i=0}^{\infty}110{,}25\,\delta^{2i+1} \geq 108 .
\end{align*}
\textbf{Si comienza la firma 2:}
\begin{align*}
  \text{Firma 2:}\quad & \sum_{i=0}^{\infty}110{,}25\,\delta^{2i} \geq 110{,}25 + 108\delta + 0, \\
  \text{Firma 1:}\quad & \sum_{i=0}^{\infty}144\,\delta^{2i+1} \geq 144 + \sum_{i=1}^{\infty}63\,\delta^{i}.
\end{align*}
Resolviendo, en ambos casos se obtiene $81\delta^2 + 81\delta - 144 \geq 0$ para la firma 1, es decir $\delta \geq 0{,}924$, y $108\delta^2 + 110{,}25\delta - 108 \geq 0$ para la firma 2, es decir $\delta \geq 0{,}612$. Con $\delta = 0{,}7$ la firma 1 no cumple, por lo que el acuerdo de turnos no es sostenible.

\end{solution}
\end{parts}

%-------------------------------------------------------------------- Galletas
\needspace{8\baselineskip}
\tema{Colusión y galletas (Bertrand repetido)}
\question Galletas y Cookies son dos empresas que venden productos homogéneos y compiten en precios. La demanda viene dada por $D(p) = 1000 - p$ y los consumidores comprarán a la firma que vende al menor precio. En el caso de que ambas empresas vendan al mismo precio, asumiremos que la mitad de los consumidores comprará a cada firma. Suponga que el costo marginal de cada firma es constante e igual a \$20.
\begin{parts}
\part Encuentre el equilibrio de Bertrand en precios si las firmas eligen precios en forma simultánea y juegan solo una vez.
\needspace{5\baselineskip}
\begin{solution}
$p_G = p_C = 20$. Para cualquier otra combinación de precios, al menos una firma tiene una desviación unilateral beneficiosa.

\end{solution}

\part Suponga que el juego se repite y se juega dos veces. Sea $0 < \delta < 1$ el factor de descuento de las firmas. Encuentre el equilibrio del juego repetido.
\needspace{5\baselineskip}
\begin{solution}
Resolvemos el equilibrio perfecto en subjuegos por inducción hacia atrás. En el período 2, el juego es idéntico a un juego de un período, por lo tanto el equilibrio es $p_{G,2} = p_{C,2} = 20$. En el período 1, las firmas reconocen que en el período 2 no se puede sostener ningún equilibrio que no sea precios iguales a costos marginales. Por lo tanto, no tienen incentivos para cooperar y el único equilibrio posible es el del juego estático: $p_{G,1} = p_{C,1} = 20$.

\end{solution}

\part Suponga que el juego se repite infinitas veces. Proponga una estrategia de gatillo para cada firma de forma de mantener el precio colusivo (el precio de un monopolista) en cada período, asumiendo que cada firma vuelve al equilibrio de Bertrand si alguna firma se desvía del precio colusivo. Calcule el valor mínimo de $\delta$ tal que el precio colusivo es el equilibrio de este juego.
\needspace{5\baselineskip}
\begin{solution}
El precio de monopolio $p^M = 510$ (con la cantidad $q^M = 490$) maximiza los beneficios conjuntos de las firmas, $\pi^M = 240\,100$.

\emph{Estrategia gatillo:} la firma $i$ fija $p_{i,t} = 510$ si en todos los períodos anteriores ambas fijaron $510$, y fija $p_{i,t} = 20$ en caso contrario.

Beneficio de cooperar:
\[
  \pi_i^{COOP} = \frac{240\,100}{2} + \delta\frac{240\,100}{2} + \delta^2\frac{240\,100}{2} + \dots = \frac{240\,100}{2}\cdot\frac{1}{1-\delta}.
\]
Beneficio de desviarse en $t$: $\pi_i^{DESV} = 240\,100 + \delta\cdot 0 + \delta^2\cdot 0 + \dots = 240\,100$.

Para una historia del juego donde no han existido desvíos: $\pi_i^{COOP} \geq \pi_i^{DESV}$ implica $\frac{120\,050}{1-\delta} \geq 240\,100$, es decir, $\delta \geq \frac{1}{2}$.

Para una historia del juego donde sí han existido desvíos, el rival juega $p_2 = 20$ hasta el final; para la firma 1, jugar $p_1 = 20$ (como lo dice su estrategia de gatillo) es mejor respuesta para cualquier $\delta$.

\end{solution}

\part Suponga ahora que Galletas y Cookies producen galletas de navidad, un producto que se vende solo durante los últimos 3 meses del año. Si definimos un período como un trimestre, las empresas interactúan cada 4 períodos. ¿Es más difícil o más fácil coludir si las empresas interactúan solo en un trimestre del año? No necesita calcular, solo puede explicarlo.
\needspace{5\baselineskip}
\begin{solution}
La colusión es más difícil cuando las ventas son menos frecuentes. Al bajar la frecuencia de las ventas, las ganancias de desviarse en $t$ no cambian, pero los beneficios futuros de cooperar disminuyen (el castigo llega más tarde, y en términos de valor presente pesa menos), por lo que es más difícil sostener un acuerdo colusivo.

\end{solution}
\end{parts}

%-------------------------------------------------------------------- Colusión con 3 firmas
\needspace{8\baselineskip}
\tema{Colusión con tres firmas, horizonte infinito y finito}
\question Un mercado enfrenta una función de demanda equivalente a $Q(P) = 12 - P$. Actualmente existe un conjunto de firmas, las cuales compiten en precios.
\begin{parts}
\part Suponga que el mercado está constituido por 3 firmas. Determine el equilibrio del mercado si las firmas compiten en precios y presentan costos marginales $c_1 < c_2 = c_3$.
\needspace{5\baselineskip}
\begin{solution}
Dado que no existen restricciones de capacidad y la competencia es en precios, el equilibrio del mercado viene dado por $(P_1, P_2, P_3) = (c_2 - \epsilon, c_2, c_2)$, con $c_2 = c_3$: la firma 1 se queda con todo el mercado.

\end{solution}

\part Suponga que la firma 1 presenta un costo marginal de producción equivalente a 2, mientras que las firmas 2 y 3 tienen cada una un costo marginal de 4. Para un período de tiempo en particular, ¿una estrategia colusiva en precios sería beneficiosa para la firma 1? Justifique en base a los beneficios que obtendría dicha firma.
\needspace{5\baselineskip}
\begin{solution}
Si las firmas no cooperan, la firma 1 recorta el precio mínimo de las otras dos: $\pi_1^{NC} = (4 - \epsilon - 2)(12 - 4 - \epsilon) \approx 16$. Si deciden cooperar, lo razonable es que todo lo produzca la firma más eficiente, todas adoptan el precio monopólico para $c_1 = 2$ ($P = 7$, $Q = 5$, $\pi = 25$) y se reparten los beneficios de manera simétrica, es decir, $\pi_1^C = \frac{1}{3}\cdot 25 = \frac{25}{3}$. Por último, si la firma eficiente decide desviarse y recortar el precio acordado, obtendrá $\pi_1^D = (7 - \epsilon - 2)\cdot 5 = 25$.

Para que la colusión sea la mejor estrategia para la firma 1 debe suceder que $\pi_1^C \geq \pi_1^D$, es decir, $\frac{25}{3} \geq 25$, lo que no se cumple. Es claro que la firma 1 preferirá desviarse antes que respetar un acuerdo colusivo, e incluso preferirá competir ($16 > \frac{25}{3}$).

\end{solution}

\part Gracias a una innovación tecnológica, ahora todas las firmas presentan costos marginales equivalentes a 2. Debido a la fuerte competencia, las firmas acuerdan coludirse de manera indefinida. Obtenga el equilibrio y los beneficios de las firmas bajo esta estrategia de largo plazo si las firmas acuerdan llevar a cabo una colusión en precios.
\needspace{5\baselineskip}
\begin{solution}
Si las firmas compiten en precios teniendo costos marginales simétricos, el resultado es la paradoja de Bertrand, con $\pi_i^{NC} = 0$. Si deciden coludirse fijando el precio monopólico correspondiente a la demanda del mercado y a los costos marginales (y también en el caso de desvío), los resultados son idénticos a los del caso anterior: $\pi_i^C = \frac{25}{3}$ y $\pi_i^D = 25$. Usando el teorema de Folk:
\[
  \delta \geq \frac{\pi_i^D - \pi_i^C}{\pi_i^D - \pi_i^{NC}} = \frac{25 - \frac{25}{3}}{25 - 0} = \frac{2}{3}.
\]

\end{solution}

\part Suponga que las firmas saben que solamente estarán por dos períodos en el mercado. Recalcule los beneficios de realizar una estrategia colusiva en precios.
\needspace{5\baselineskip}
\begin{solution}
En este caso los pagos son los mismos, pero no es posible aplicar el teorema de Folk y se debe plantear la condición $VP(\text{colusión}) \geq VP(\text{desvío})$:
\[
  \pi_i^C + \delta\pi_i^C \geq \pi_i^D + \delta\pi_i^{NC}
  \;\Longrightarrow\;
  (1+\delta)\frac{25}{3} \geq 25
  \;\Longrightarrow\;
  \delta \geq 2,
\]
lo que evidentemente no es sostenible (pues $\delta < 1$).

\end{solution}

\part Compare y comente los resultados obtenidos en (c) y (d): ¿cómo afecta a la colusión el hecho de que las firmas no compitan por infinitos períodos?
\needspace{5\baselineskip}
\begin{solution}
Comparando (c) y (d): es la comprobación del teorema de Folk, que dice que cualquier resultado factible en el cual todos los jugadores obtienen un beneficio mayor que el que logran en un equilibrio de Nash de la versión estática de un juego puede sostenerse como un equilibrio de la versión infinitamente repetida del mismo, siempre que los jugadores valoren lo suficiente el futuro. La colusión que no era posible en el caso de un juego estático, ahora sí lo es en un juego infinito. Con solo dos períodos, en cambio, un resultado alcanzable mediante un juego infinitamente repetido resulta imposible.

\end{solution}
\end{parts}

%-------------------------------------------------------------------- Comente: auge vs declive
\needspace{8\baselineskip}
\tema{Comente: colusión tácita en industrias en auge y en declive}
\question \textbf{Comente.} Comparando una industria en auge (demanda crece todos los años) con una industria en declive (demanda cae todos los años), ¿en cuál de las dos es más probable que surja un equilibrio de colusión tácita? (Todo lo demás constante; esto es, un mismo número de firmas, costos, etc.) Fundamente intuitivamente su respuesta.
\needspace{5\baselineskip}
\begin{solution}
Es más probable que surja en la industria en auge. La intuición es que un equilibrio de colusión tácita surge cuando los beneficios futuros de seguir coludido son mayores que los beneficios de desviarse hoy y ganar más participación de mercado en el presente. En una industria en declive, el futuro es poco rentable ya que la demanda va cayendo, haciendo más probable que las firmas se desvíen, intentando obtener más ganancias ahora ya que en el futuro no parece muy promisorio. En cambio, en una industria en auge el futuro es cada vez más rentable, por lo que es más probable que las firmas no quieran desviarse hoy para no perder los mayores beneficios futuros de una estrategia de precios altos cuando la demanda sea también alta.

\end{solution}

%-------------------------------------------------------------------- Comente: impaciencia siempre alcanza
\needspace{8\baselineskip}
\tema{Comente: ¿siempre hay un nivel de impaciencia que sostenga la cooperación?}
\question \textbf{Comente.} Dadas las condiciones, siempre habrá un nivel de impaciencia en que el equilibrio perfecto en subjuegos cooperativo sea posible.
\needspace{5\baselineskip}
\begin{solution}
Falso. Es condición necesaria tener un horizonte infinito (o incierto). Bajo horizontes finitos, al resolver por inducción hacia atrás siempre llegaremos a un equilibrio de Nash no cooperativo (el de competencia).

\end{solution}

%-------------------------------------------------------------------- Pollos
\needspace{8\baselineskip}
\tema{Colusión en Cournot con 3 y 2 firmas (los pollos)}
\question Las firmas SuperPollo, Ariztía y DonPollo son conocidas por sus carnes de pollo (un bien homogéneo). Las 3 firmas se enfrentan a una demanda de mercado $P = 12\,500 - 2Q$, donde $Q$ es la cantidad total producida en la industria. Todas las firmas tienen un costo marginal igual a 2500.
\begin{parts}
\part Suponga que las empresas compiten solo un período. Encuentre las cantidades y las utilidades si las empresas compiten en cantidades.
\needspace{5\baselineskip}
\begin{solution}
Las empresas compiten en un oligopolio de Cournot con 3 firmas. La función de beneficio de la firma 1 es $\pi_1 = (10\,000 - 2(q_1 + q_2 + q_3))q_1$, con CPO $10\,000 - 4q_1 - 2q_2 - 2q_3 = 0$. Por simetría, $q_1 = q_2 = q_3 = q^*$ y $10\,000 = 8q^*$, de donde
\[
\begin{gathered}
  q^* = 1250, \quad Q = 3750, \quad P = 12\,500 - 7500 = 5000, \\ \pi_i = (5000 - 2500)\cdot 1250 = 3\,125\,000 .
\end{gathered}
\]

\end{solution}

\part Considere ahora que las empresas compiten en un horizonte infinito. ¿Cuál es el factor de descuento que hace que las empresas cooperen, es decir, se repartan las utilidades monopólicas, en un equilibrio perfecto en subjuegos?
\needspace{5\baselineskip}
\begin{solution}
El monopolista maximiza $(10\,000 - 2Q)Q$: $Q^M = 2500$, $P^M = 7500$ y $\pi^M = 12\,500\,000$. Repartido entre 3 firmas, $\pi^{COL} = \frac{12\,500\,000}{3} \approx 4\,166\,667$ (cada una produce $\frac{2500}{3} \approx 833{,}33$).

Si una firma se desvía mientras las otras dos producen $\frac{5000}{3}$ en total, su mejor respuesta es $q^{DESV} = 2500 - \frac{5000/3}{2} = 1666{,}67$. Entonces $P = 12\,500 - 2(1666{,}67 + 1666{,}67) = 5833{,}33$ y
\[
  \pi^{DESV} = (5833{,}33 - 2500)\cdot 1666{,}67 \approx 5\,555\,556 .
\]
Con estrategia de gatillo (reversión al equilibrio de Cournot) la colusión es sostenible si
\[
  \delta \geq \frac{\pi^{DESV} - \pi^{COL}}{\pi^{DESV} - \pi^{COURNOT}} = \frac{5\,555\,556 - 4\,166\,667}{5\,555\,556 - 3\,125\,000} = \frac{1\,388\,889}{2\,430\,556} \approx 0{,}57 .
\]

\end{solution}

\part Por ciertas situaciones, Ariztía deja de competir en esta industria, dejando solo dos firmas. Encuentre el nuevo factor de descuento que hace que las empresas cooperen en un horizonte infinito.
\needspace{5\baselineskip}
\begin{solution}
Con dos firmas, la CPO de la firma $i$ es $10\,000 - 4q_i - 2q_j = 0$, es decir $q_i = 2500 - \frac{q_j}{2}$, por lo que $q^* = \frac{5000}{3} \approx 1666{,}67$, $Q = 3333{,}33$, $P = 5833{,}33$ y $\pi^{COURNOT} = 3333{,}33\cdot 1666{,}67 \approx 5\,555\,556$.

En colusión, cada una produce $1250$ y obtiene $\pi^{COL} = 6\,250\,000$. Si una firma se desvía mientras la otra coopera, produce $q^{DESV} = 2500 - \frac{1250}{2} = 1875$, el precio es $P = 12\,500 - 2(1875 + 1250) = 6250$ y $\pi^{DESV} = (6250 - 2500)\cdot 1875 = 7\,031\,250$. Entonces
\[
  \delta \geq \frac{7\,031\,250 - 6\,250\,000}{7\,031\,250 - 5\,555\,556} = \frac{781\,250}{1\,475\,694} \approx 0{,}53 .
\]

\end{solution}

\part Compare los factores de descuento obtenidos en los apartados anteriores. ¿Cuál de ellos es mayor y a qué se debe? ¿Cómo influye el número de firmas en este resultado?
\needspace{5\baselineskip}
\begin{solution}
El factor de descuento mínimo con tres firmas es $\delta_{min} \approx 0{,}57$, mientras que con dos firmas es $\delta_{min} \approx 0{,}53$. Se requiere mayor paciencia para sostener la colusión cuando hay más competidores. Esto se debe a varios factores: primero, con más firmas cada una recibe una porción menor de las utilidades monopólicas, lo que reduce el beneficio de la cooperación; segundo, el potencial de ganancias inmediatas por desviarse aumenta (se desvía contra rivales que producen poco cada una); por último, el castigo es menos severo cuando hay más firmas, pues las utilidades de Cournot son menores con más competidores pero no disminuyen tan rápidamente como la participación en las utilidades monopólicas. Todo esto hace más difícil sostener la colusión con un mayor número de firmas.

\end{solution}
\end{parts}

\end{unidad}

\unidadtitulo{Diferenciación de productos}
\begin{unidad}
\tema{Tres firmas con bienes diferenciados y fusión}
\question Suponga un mercado en que tres firmas 1, 2 y 3 ofrecen productos diferenciados y compiten en precios. La demanda que enfrenta el producto ofrecido por cada firma $i$ viene dada por $q_i = 100 - 3p_i + p_j + p_k$. Esta función de demanda es simétrica para los tres productos, de manera que
\begin{align*}
  q_1 &= 100 - 3p_1 + p_2 + p_3 \\
  q_2 &= 100 - 3p_2 + p_1 + p_3 \\
  q_3 &= 100 - 3p_3 + p_1 + p_2
\end{align*}
La función de costos de producción es idéntica para cada firma y viene dada por $C(q_i) = 20q_i$. Considerando esta información responda las siguientes preguntas:
\begin{parts}
\part Calcule las funciones de reacción de las firmas 1, 2 y 3. Grafique las funciones de reacción de las firmas 1 y 2 (asumiendo $\bar p_3$ como constante). Identifique el precio de equilibrio en su gráfico.
\needspace{5\baselineskip}
\begin{solution}
Planteamos el problema de la firma 1:
\[
  \max_{p_1}\ \pi_1 = (p_1 - 20)(100 - 3p_1 + p_2 + p_3)
\]
CPO:
\[
\begin{gathered}  \frac{\partial \pi_1}{\partial p_1} = 100 - 3p_1 + p_2 + p_3 - 3(p_1 - 20) = 160 - 6p_1 + p_2 + p_3 = 0
  \quad\Longrightarrow\quad\\
  p_1(p_2,p_3) = \frac{160 + p_2 + p_3}{6}.\end{gathered}
\]
Por simetría, para cualquier empresa,
\[
  p_i(p_j,p_k) = \frac{160 + p_j + p_k}{6}.
\]
Para graficar las funciones de reacción dejamos $\bar p_3$ como constante, de modo que $p_i(p_j) = \frac{p_j}{6} + \text{cte.}$ Las funciones de reacción tienen pendiente positiva (los precios son \emph{complementos estratégicos}) y las curvas de las firmas 1 y 2 se cruzan en el equilibrio de Nash.

\end{solution}

\part Obtenga el vector de precios $(p_1^*, p_2^*, p_3^*)$ que corresponde al equilibrio de Nash de este juego.
\needspace{5\baselineskip}
\begin{solution}
Aplicamos que $p_1^* = p_2^* = p_3^* = p^*$, y reemplazamos en alguna de las funciones de reacción:
\[
\begin{gathered}
  p^* = \frac{160 + p^* + p^*}{6} \\
  \Longrightarrow\; 4p^* = 160 \\
  \Longrightarrow\; p^* = 40 .
\end{gathered}
\]
El vector de equilibrio es $(40, 40, 40)$.

\end{solution}

\part Asuma que las firmas 2 y 3 deciden salirse del mercado por motivos internos. Ante esta situación la firma 1 logra adquirir la firma 2, por lo que las demandas pasan a ser las siguientes:
\begin{align*}
  q_1 &= 100 - 2p_1 + p_2 \\
  q_2 &= 100 - 2p_2 + p_1
\end{align*}
Calcule el precio, la cantidad y los beneficios de equilibrio bajo este nuevo escenario.
\needspace{5\baselineskip}
\begin{solution}
Planteamos ahora el problema de maximización de un monopolio multiproducto, $\pi = \sum_{i=1}^{2}\pi_i$:
\[
  \max_{p_1,p_2}\ \pi = (100 - 2p_1 + p_2)(p_1 - 20) + (100 - 2p_2 + p_1)(p_2 - 20)
\]
CPO:
\[
\begin{gathered}  \frac{\partial \pi}{\partial p_1} = 100 - 2p_1 + p_2 - 2(p_1 - 20) + (p_2 - 20) = 120 - 4p_1 + 2p_2 = 0
  \quad\Longrightarrow\quad\\
  p_1^*(p_2) = \frac{60 + p_2}{2}.\end{gathered}
\]
Por simetría, $p_2^*(p_1) = \frac{60 + p_1}{2}$, y podemos decir que $p_1 = p_2 = p$, por tanto $p = 60$ y consecuentemente $q = 100 - 120 + 60 = 40$. El beneficio es
\[
  \pi = 40\cdot(60-20) + 40\cdot(60-20) = 3200 .
\]

\end{solution}

\part ¿Cómo cambia su respuesta del ítem anterior si ahora las demandas son
\begin{align*}
  q_1 &= 100 - 2p_1 - p_2 \\
  q_2 &= 100 - 2p_2 - p_1
\end{align*}
¿Qué significa una demanda de este tipo?
\needspace{5\baselineskip}
\begin{solution}
Una demanda de este tipo refiere a que los productos 1 y 2 son \emph{complementarios} (un mayor precio de un bien reduce la demanda del otro), por lo que los beneficios del monopolio multiproducto serán menores. Calculémoslo:
\[
  \max_{p_1,p_2}\ \pi = (100 - 2p_1 - p_2)(p_1 - 20) + (100 - 2p_2 - p_1)(p_2 - 20)
\]
CPO:
\[
\begin{gathered}  \frac{\partial \pi}{\partial p_1} = 100 - 2p_1 - p_2 - 2(p_1 - 20) - (p_2 - 20) = 160 - 4p_1 - 2p_2 = 0
  \quad\Longrightarrow\quad\\
  p_i^*(p_j) = \frac{80 - p_j}{2}.\end{gathered}
\]
Por simetría $p_1=p_2=p$, de modo que $3p = 80$, es decir $p = \frac{80}{3} \approx 26{,}7$ y $q = 100 - 3p = 20$. El beneficio es
\[
  \pi = 2\cdot 20\cdot\Bigl(\frac{80}{3} - 20\Bigr) = \frac{800}{3} \approx 267 .
\]
Fíjese que $\pi_{\text{Sustitutos}} = 3200 > \pi_{\text{Complementos}} \approx 267$.

\end{solution}
\end{parts}

%-------------------------------------------------------------------- Comente: diferenciación y precios
\needspace{8\baselineskip}
\tema{Comente: diferenciación, precios y utilidades}
\question \textbf{Comente.} A mayor diferenciación de los productos, mayores serán los precios y utilidades de las firmas.
\needspace{5\baselineskip}
\begin{solution}
Verdadero. Mientras más diferenciados son los productos, mayor será la demanda cautiva (aquella que le permite subir los precios) de las firmas, lo que se traduce en un equilibrio con mayores precios en el mercado. En conclusión, a mayor diferenciación de los bienes, menor es la intensidad de la competencia.

\end{solution}

%-------------------------------------------------------------------- Comente: complementarios monopolio vs Bertrand
\needspace{8\baselineskip}
\tema{Comente: bienes complementarios, monopolio vs Bertrand}
\question \textbf{Comente.} Imagine un mercado donde existen dos bienes diferenciados complementarios, que enfrentan demandas lineales. En términos de los precios percibidos por los consumidores, será preferible que dicho mercado sea monopólico a tener una competencia oligopólica à la Bertrand.
\needspace{5\baselineskip}
\begin{solution}
Verdadero. Esto es lo que motiva la existencia de fusiones verticales o la entrega de patentes. Recordemos que ante bienes complementarios la demanda del bien 1 aumenta en la medida que disminuye el precio del bien 2, mientras que en bienes sustitutos disminuye. Esto genera incentivos para que un monopolista maximice el beneficio en un equilibrio donde el precio es menor al de competencia, ya que el monopolio percibe la totalidad de la demanda. En cambio, en el escenario de competencia en precios la demanda se divide, lo que lleva a que el óptimo de cada firma se encuentre en un precio mayor (no internaliza el efecto de su precio sobre la demanda del complemento).

\end{solution}

%-------------------------------------------------------------------- Comente: diferenciación e intensidad competitiva
\needspace{8\baselineskip}
\tema{Comente: diferenciación e intensidad competitiva}
\question \textbf{Comente.} Explique el efecto de la diferenciación de producto en la intensidad competitiva.
\needspace{5\baselineskip}
\begin{solution}
Desde el punto de vista de la competencia, la diferencia con el caso de producto homogéneo radica en que existen más variables de decisión para los individuos que solo el precio. Así, puede pasar que un consumidor prefiera a un productor que cobra más caro porque su producto se alinea mejor con las preferencias propias.

El efecto de la diferenciación es atenuar la competencia en precios, dando pie a que sea factible cobrar sobre los costos marginales y, por ende, obtener rentas.

\end{solution}

%-------------------------------------------------------------------- Comente: horizontal vs vertical
\needspace{8\baselineskip}
\tema{Comente: diferenciación horizontal vs vertical}
\question \textbf{Comente.} ¿Cuál es la diferencia entre una diferenciación horizontal y una vertical?
\needspace{5\baselineskip}
\begin{solution}
Una diferenciación horizontal refiere a caracterizarse por medio de características de un mismo producto que no afectan a la calidad. Esto puede ser estética (autos), sabor (cereales), aroma (perfumes), y suele ser ejemplificado como la ubicación de venta del producto.

Por otro lado, una diferenciación vertical hace referencia a un espectro de calidad de un mismo producto. Incluso no teniendo diferentes costos de producción, las firmas decidirán hacer productos de distinta calidad para disminuir la competencia.

Al fin y al cabo, ambas suelen tener un mismo resultado: las empresas prefieren diferenciarse para relajar la competencia.

\end{solution}

%-------------------------------------------------------------------- Comente: horizontal producto no cambia
\needspace{8\baselineskip}
\tema{Comente: la diferenciación horizontal no cambia el producto}
\question \textbf{Comente.} Al diferenciarse horizontalmente el producto nunca cambia.
\needspace{5\baselineskip}
\begin{solution}
Falso. Hay diferenciaciones horizontales que cambian características del producto. Si alguien está en búsqueda de un auto coupé deportivo de alta gama, puede encontrarse con un Porsche 911 Carrera o bien con un Chevrolet Corvette Stingray. Mientras un Porsche tiene un diseño más clásico, el Chevrolet tiene un diseño más audaz. Ambos tienen un precio y calidad similar, pero se diferencian en su posicionamiento frente al potencial comprador.

\end{solution}

%-------------------------------------------------------------------- Comente: Salop vs Hotelling
\needspace{8\baselineskip}
\tema{Comente: Salop vs Hotelling}
\question \textbf{Comente.} El modelo de Salop es lo mismo que el modelo de Hotelling pero asumiendo una ciudad circular.
\needspace{5\baselineskip}
\begin{solution}
Falso. El modelo de Salop tiene como objetivo entender la entrada de firmas. Se supone una ciudad circular para dar por sentado que no hay ubicaciones privilegiadas sobre otras: todas las empresas son equidistantes entre sí.

\end{solution}

%-------------------------------------------------------------------- Comente: más o menos diferenciación
\needspace{8\baselineskip}
\tema{Comente: efecto demanda y efecto estratégico en Hotelling}
\question \textbf{Comente.} En un modelo tipo Hotelling con decisiones de localización y luego competencia en precios, ¿cuáles son las razones a favor y en contra de que las firmas produzcan bienes cada vez más diferenciados?
\needspace{5\baselineskip}
\begin{solution}
El \emph{efecto de demanda} crea incentivos para que las firmas produzcan bienes con poca diferenciación: cuando las empresas se acercan hacia el centro de la ciudad capturan mayor demanda. Sin embargo, el \emph{efecto estratégico} nos dice que la competencia en precios es más fuerte cuanto menos diferenciados sean los bienes: cuando las empresas se alejan del centro de la ciudad reducen la competencia por precios.

\end{solution}

%-------------------------------------------------------------------- Hotelling
\needspace{8\baselineskip}
\tema{Modelo de Hotelling con costos de transporte cuadráticos}
\question Considere una ciudad lineal que va de 0 a 1. Dos empresas $L$ y $R$ deciden en qué parte ubicarse, $\theta_L, \theta_R \in [0,1]$, con $\theta_L < \theta_R$. Ambas ofrecen productos homogéneos (son sustituibles) a un precio $p_L, p_R$ según cada firma, con costo marginal $c$. Los potenciales consumidores de estas firmas se distribuyen de forma uniforme y los caracteriza la siguiente función de utilidad:
\[
  U_{ij} = \bar u + (y - p_j) - \tau(\theta_j - v_i)^2 .
\]
\begin{parts}
\part Explique cada parte de la función de utilidad y su interpretación intuitiva.
\needspace{5\baselineskip}
\begin{solution}
La función de utilidad es del individuo $i$ comprando a la firma $j$. En primer lugar, $\bar u$ es una utilidad fija que les brinda el producto, e $y - p_j$ es la utilidad neta de comprar el producto a la firma $j$ (ingreso menos precio). Por último, $\tau(\theta_j - v_i)^2$ es el costo de transporte que incurre el individuo $i$ (ubicado en $v_i$) para comprarle a la firma $j$, que crece con el cuadrado de la distancia.

\end{solution}

\part ¿Cuántos individuos indiferentes hay en una ciudad lineal? ¿Qué caracteriza a estos individuos? Encuentre su ubicación.
\needspace{5\baselineskip}
\begin{solution}
En una ciudad lineal habrá un único individuo indiferente entre comprar a la firma $L$ o $R$. Este individuo es tal que $U_{iL} = U_{iR}$. Su ubicación $\bar v$ se encuentra planteando tal ecuación:
\begin{gather*}
  \bar u + (y - p_L) - \tau(\theta_L - \bar v)^2 = \bar u + (y - p_R) - \tau(\theta_R - \bar v)^2 \\
  (p_R - p_L) + \tau(\theta_R^2 - \theta_L^2) = 2\tau\bar v(\theta_R - \theta_L) \\
  \bar v = \frac{p_R - p_L}{2\tau(\theta_R - \theta_L)} + \frac{\theta_R + \theta_L}{2}.
\end{gather*}
Lo único diferente de ambos productos es el lugar de la ciudad en que se venden.

\end{solution}

\part Calcule las cuotas de mercado de ambas firmas.
\needspace{5\baselineskip}
\begin{solution}
La cuota de mercado de la firma $L$ son todos los individuos a la izquierda de $\bar v$, y $R$ se queda con el resto del mercado, $1 - \bar v$:
\[
\begin{gathered}
  D_L(p,\theta;\tau) = \bar v = \frac{p_R - p_L}{2\tau(\theta_R - \theta_L)} + \frac{\theta_R + \theta_L}{2}, \\
  D_R(p,\theta;\tau) = 1 - \bar v = 1 - \frac{p_R - p_L}{2\tau(\theta_R - \theta_L)} - \frac{\theta_R + \theta_L}{2}.
\end{gathered}
\]
En el modelo de Hotelling las firmas, en una primera etapa, eligen dónde ubicarse en la ciudad, para luego, en una segunda etapa, empezar a vender a un cierto precio.

\end{solution}

\part Suponga que bajo un nuevo marco legal el precio del producto está fijado. ¿Dónde les conviene ubicarse a las firmas?
\needspace{5\baselineskip}
\begin{solution}
Dado que las firmas están obligadas a cobrar un precio fijo, dada la ubicación de la competencia la mejor respuesta siempre será estar un $\epsilon$ más cerca del centro ($0{,}5$). El equilibrio será $\theta_R = \theta_L = 0{,}5$. Aquí estamos frente a un caso de \emph{mínima diferenciación}: como los precios están dados, las firmas intentan maximizar su cuota de mercado mediante su ubicación. Con $p_R = p_L$ y $\theta_R = \theta_L = 0{,}5$, la demanda de cada una es $0{,}5$ y se reparten el mercado en partes iguales.

\end{solution}

\part Ahora volvamos al escenario sin el marco regulatorio: las firmas eligen sus precios. Encuentre las funciones de reacción de ambas firmas. Encuentre el equilibrio de Nash.
\needspace{5\baselineskip}
\begin{solution}
Hacemos un cambio de variable para simplificar: $t = \tau(\theta_R - \theta_L)$ y $s = \frac{\theta_R + \theta_L}{2}$. Entonces $D_L = \frac{p_R - p_L}{2t} + s$ y $D_R = \frac{p_L - p_R}{2t} + 1 - s$.

Para la empresa $L$:
\[
\begin{gathered}
  \max_{p_L}\ \pi_L = (p_L - c)\Bigl(\frac{p_R - p_L}{2t} + s\Bigr) \\
  \;\Rightarrow\; \frac{p_R - 2p_L + 2st + c}{2t} = 0 \\
  \;\Rightarrow\; p_L(p_R) = \frac{p_R + c}{2} + st .
\end{gathered}
\]
Para la empresa $R$:
\[
  \max_{p_R}\ \pi_R = (p_R - c)\Bigl(\frac{p_L - p_R}{2t} + 1 - s\Bigr)
  \;\Rightarrow\; p_R(p_L) = \frac{p_L + c}{2} + t(1-s).
\]
Deshaciendo el cambio de variable, la función de reacción de $L$ es $p_L(p_R) = \frac{p_R + c}{2} + \frac{\tau(\theta_R^2 - \theta_L^2)}{2}$. Resolviendo el sistema se obtiene el equilibrio de Nash:
\[
  p_L^* = c + \frac{\tau(\theta_R - \theta_L)(2 + \theta_L + \theta_R)}{3}, \qquad
  p_R^* = c + \frac{\tau(\theta_R - \theta_L)(4 - \theta_L - \theta_R)}{3}.
\]
En particular, cuando las firmas están a una misma distancia del centro ($\theta_L + \theta_R = 1$) los precios son iguales, $p_L^* = p_R^*$, y
\[
  p^* = c + \tau(\theta_R - \theta_L).
\]

\end{solution}

\part Grafique las funciones de reacción. Muestre cómo se mueven las curvas al cambiar los costos de transporte.
\needspace{5\baselineskip}
\begin{solution}
Las funciones de reacción son rectas con pendiente positiva $\frac{1}{2}$ (los precios son complementos estratégicos), con intercepto $\frac{c}{2} + st$ para $L$ y $\frac{c}{2} + t(1-s)$ para $R$. Un aumento en los costos de transporte $\tau$ aumenta $t$ y, por lo tanto, desplaza ambas curvas hacia arriba (hacia afuera): el equilibrio se mueve a precios más altos, pues la diferenciación relaja la competencia en precios. En el límite $\tau \to 0$ los productos son homogéneos y se llega al equilibrio de Bertrand $p = c$.

\end{solution}

\part Suponga que usted es un planificador social omnipotente que busca maximizar el bienestar de los consumidores. ¿Cómo cree que deberían ubicarse las firmas?
\needspace{5\baselineskip}
\begin{solution}
La solución socialmente óptima es la que minimiza los costos de transporte, y sería $\theta_L = \frac{1}{4}$ y $\theta_R = \frac{3}{4}$. Por tanto, desde el punto de vista social hay demasiada diferenciación del producto cuando el mercado es privado (en el equilibrio privado con precios endógenos y costos cuadráticos las firmas se ubican en los extremos, más lejos que el óptimo social).

\end{solution}
\end{parts}

%-------------------------------------------------------------------- Apple y Dell
\needspace{8\baselineskip}
\tema{Competencia entre Apple y Dell (fusión con bienes diferenciados)}
\question Suponga que Apple y Dell son los únicos competidores en el mercado de los notebooks. Apple y Dell venden productos diferenciados horizontalmente (Macbook y XPS) y compiten en precios. Las funciones de demanda de Apple y Dell son
\begin{align*}
  q_A &= 2 - 2p_A + p_D \\
  q_D &= 2 + p_A - 2p_D
\end{align*}
Los costos marginales para ambas firmas son constantes e iguales a $c_A = c_D = \frac{1}{2}$.
\begin{parts}
\part Encuentre los precios, cantidades y beneficios de cada firma en el equilibrio de Nash en precios. Grafique las funciones de mejor respuesta y el equilibrio.
\needspace{5\baselineskip}
\begin{solution}
Apple maximiza
\[
  \max_{p_A}\ \pi_A = \Bigl(p_A - \frac{1}{2}\Bigr)(2 - 2p_A + p_D).
\]
Derivando:
\[
  \frac{\partial\pi_A}{\partial p_A} = (2 - 2p_A + p_D) - 2\Bigl(p_A - \frac{1}{2}\Bigr) = 3 - 4p_A + p_D = 0
  \;\Longrightarrow\;
  p_A(p_D) = \frac{3}{4} + \frac{1}{4}p_D .
\]
Por simetría, $p_D(p_A) = \frac{3}{4} + \frac{1}{4}p_A$. En el equilibrio simétrico $p_A = p_D = p$: $3 - 3p = 0$, es decir
\[
  p_A^* = p_D^* = 1, \qquad q_A^* = q_D^* = 1, \qquad \pi_A^* = \pi_D^* = \frac{1}{2}.
\]
Las funciones de mejor respuesta son rectas de pendiente $\frac14$ (precios como complementos estratégicos) que se intersectan en $(1,1)$.

\end{solution}

\part Suponga que Apple y Dell deciden fusionarse y formar Dellapple, que se comportará como un monopolista multiproducto (sin sinergias como consecuencia de la fusión). Encuentre los precios, cantidades y beneficios del monopolista multiproducto.
\needspace{5\baselineskip}
\begin{solution}
Dellapple maximiza
\[
  \max_{p_A,p_D}\ \Bigl(p_A - \frac{1}{2}\Bigr)(2 - 2p_A + p_D) + \Bigl(p_D - \frac{1}{2}\Bigr)(2 - 2p_D + p_A).
\]
Derivando respecto a $p_A$:
\[
  (2 - 2p_A + p_D) - 2\Bigl(p_A - \frac{1}{2}\Bigr) + \Bigl(p_D - \frac{1}{2}\Bigr) = \frac{5}{2} - 4p_A + 2p_D = 0 .
\]
Usando simetría, $\frac{5}{2} - 2p = 0$, de donde
\[
  p_A = p_D = \frac{5}{4}, \qquad q_A = q_D = \frac{3}{4}, \qquad \pi_A = \pi_D = \frac{9}{16}, \qquad \pi_{\text{total}} = \frac{9}{8}.
\]
(Cada producto genera $\bigl(\frac54 - \frac12\bigr)\cdot\frac34 = \frac{9}{16}$, y el beneficio conjunto de $\frac98 > 1$ es mayor que el de la competencia.)

\end{solution}

\part ¿Los consumidores prefieren que las firmas compitan (como en (a)) o que exista un monopolista (como en (b))? ¿Su respuesta sería distinta si los bienes fuesen complementarios en lugar de sustitutos? Explique.
\needspace{5\baselineskip}
\begin{solution}
Los consumidores prefieren que las firmas compitan porque los precios de ambos productos son menores ($1 < \frac{5}{4}$) y, por ende, el excedente del consumidor es mayor. Si los bienes fuesen complementarios, los precios fijados por el monopolista multiproducto serían menores, por lo que con bienes complementarios los consumidores prefieren un monopolio a que las firmas compitan.

La intuición es la siguiente: si los productos son complementarios, cuando Apple fija un precio no tiene en cuenta que esto tiene un efecto sobre la demanda de Dell. A medida que Apple disminuye su precio, genera una externalidad positiva sobre Dell, que vende un bien complementario. Dellapple internaliza dicho efecto y fija precios menores que las firmas que operan de manera independiente.

\end{solution}
\end{parts}

%-------------------------------------------------------------------- Comente: vertical pura
\needspace{8\baselineskip}
\tema{Comente: diferenciación vertical pura y costos marginales}
\question \textbf{Comente.} Según los modelos de diferenciación vertical pura, la diferencia de precios se debe a una diferencia en los costos marginales de producir bienes de mayor o menor calidad.
\needspace{5\baselineskip}
\begin{solution}
Falso. En los modelos de diferenciación vertical pura, aunque no haya diferencias en los costos de producir bienes de menor o mayor calidad, habrá una diferenciación. Esto se debe a que, ante máxima diferenciación, las firmas relajan la competencia por precios (y obtienen beneficios positivos).

\end{solution}

%-------------------------------------------------------------------- Comente: con y sin localización
\needspace{8\baselineskip}
\tema{Comente: modelos de diferenciación con y sin localización}
\question \textbf{Comente.} La principal diferencia entre los modelos de diferenciación con y sin localización es que en los modelos sin localización los consumidores tienen utilidad de la variedad de productos, mientras que en los modelos con localización el consumidor compra solamente una marca (1 ordenador, 1 casa, etc.).
\needspace{5\baselineskip}
\begin{solution}
Verdadero. Cuando vemos Hotelling con diferenciación horizontal estamos utilizando modelos de utilidad discreta: se compra una unidad de alguno de los bienes que ofrecen las empresas. En estos modelos, si existen $J$ alternativas en el mercado, indexadas por $j = 1,\dots,J$, el problema del consumidor $i$ se reduce a elegir la que más utilidad le brinde:
\[
  \max_{j,q}\ U_i(x_j, q) \quad \text{sujeto a} \quad p_j + p_0 q = y_i ,
\]
donde $x_j$ son las características de la marca $j$ y su precio es $p_j$, mientras que $q$ es la cantidad de la opción alternativa (\emph{outside good}) y $p_0$ su precio (generalmente se normaliza a 1). La opción alternativa (que se denota como $j = 0$) es la opción de no comprar (lo que implica que se gasta el ingreso disponible en otros bienes).

\end{solution}

%-------------------------------------------------------------------- Comente: vertical cambia producto
\needspace{8\baselineskip}
\tema{Comente: la diferenciación vertical cambia el producto}
\question \textbf{Comente.} Al diferenciarse verticalmente el producto cambia, por ejemplo añadiendo funciones o mejorando la calidad de los materiales.
\needspace{5\baselineskip}
\begin{solution}
Verdadero. La diferenciación vertical cambia características del producto que inciden directamente en su calidad. Estos modelos se diferencian muy poco en aspecto, apuntan al mismo mercado, pero difieren en calidad de materiales y funcionalidades. Una marca ofrece, por ejemplo, relojes de calidades distintas pero con aspectos similares porque en un mismo mercado habrá personas más dispuestas que otras a pagar por calidad; para obtener mayor excedente de estas personas se mejora la calidad del bien cobrando un mayor precio. (En cambio, en la diferenciación horizontal los cambios no son rankeables: una marca introduce relojes de calidades similares pero de aspectos distintos porque darle en el gusto a las personas aumenta su disposición a pagar y, en competencia, el precio que se puede cobrar sobre el costo marginal.)

\end{solution}

%-------------------------------------------------------------------- Diferenciación vertical
\needspace{8\baselineskip}
\tema{Diferenciación vertical: calidades y precios}
\question Considere que las empresas no se distinguen en características horizontales de los productos, pero pueden decidir en el primer turno la calidad de su producto y en el segundo competir por precios. Podríamos imaginar una distribución lineal de calidad del producto, donde 1 es la máxima calidad y 0 la mínima. La empresa de menor calidad se ubica en $a$ y la de mayor calidad en $b$, con $a < b$. Considere un individuo con utilidad $U_i^\theta(p_i)$, donde $\theta \in [0,1]$ es la posición del individuo (su disposición a pagar por calidad, distribuida uniformemente) e $i = a,b$ es la empresa a la que le compra:
\[
  U_i^\theta(p_i) =
  \begin{cases}
    \theta a - p_a & \text{si } i = a \\
    \theta b - p_b & \text{si } i = b
  \end{cases}
\]
Suponga que no hay costos marginales y que el mercado está cubierto (todos compran alguno de los dos bienes).
\begin{parts}
\part ¿Cuál es el procedimiento para resolver este tipo de juego?
\needspace{5\baselineskip}
\begin{solution}
Estos juegos por turnos se resuelven por inducción hacia atrás para encontrar un equilibrio de Nash consistente con individuos que eligen su mejor respuesta mirando hacia el futuro. Es decir, se debe pensar en el precio que se fijaría dadas las ubicaciones en calidad, para luego elegir las ubicaciones que se tomarían. De lo contrario, si se eligiera una ubicación sin pensar en la pelea por precios que habría posteriormente, puede que lleguemos a un equilibrio (dadas las funciones de reacción de precios) que no sea consistente con individuos racionales que miran hacia el futuro.

\end{solution}

\part Encuentre el punto $\hat\theta$ en que se encuentra el consumidor indiferente. Demuestre que comprar calidad hace más felices a los dispuestos a pagar por calidad que a los que no valoran tanto la calidad.
\needspace{5\baselineskip}
\begin{solution}
El consumidor indiferente cumple $U_a^{\hat\theta}(p_a) = U_b^{\hat\theta}(p_b)$:
\[
  \hat\theta a - p_a = \hat\theta b - p_b
  \;\Longrightarrow\;
  \hat\theta(a - b) = p_a - p_b
  \;\Longrightarrow\;
  \hat\theta = \frac{p_b - p_a}{b - a}.
\]
Para ver que comprar calidad hace más felices a los que están dispuestos a pagar por calidad, evaluamos la diferencia de utilidades para un individuo $\hat\theta + \delta$ con $\delta > 0$:
\[
  U_b^{\hat\theta+\delta} - U_a^{\hat\theta+\delta} = \underbrace{\hat\theta(b-a) - p_b + p_a}_{=0} + \delta(b - a) = \delta(b-a) > 0,
\]
y para $\hat\theta - \delta$ la diferencia es $-\delta(b-a) < 0$. Es decir, los de mayor $\theta$ prefieren la alta calidad y los de menor $\theta$ prefieren la baja calidad.

\end{solution}

\part Encuentre las funciones de reacción de las firmas dadas las ubicaciones. Suponga que no hay costos marginales.
\needspace{5\baselineskip}
\begin{solution}
La demanda por los bienes de la empresa de calidad $a$ está dada por los individuos a la izquierda de $\hat\theta$, por lo que esa empresa maximiza
\[
  \max_{p_a}\ \pi_a(p_b; b, a) = p_a\hat\theta = p_a\,\frac{p_b - p_a}{b - a}
  \;\Longrightarrow\;
  \frac{-2p_a + p_b}{b-a} = 0
  \;\Longrightarrow\;
  p_a(p_b) = \frac{1}{2}p_b .
\]
Es decir, dadas las diferencias en calidad, el productor del producto de menor calidad fijará la mitad del precio del otro. Para la firma de mayor calidad es análogo:
\[
\begin{gathered}
  \max_{p_b}\ \pi_b = p_b(1 - \hat\theta) = p_b\Bigl(1 - \frac{p_b - p_a}{b-a}\Bigr) \\
  \;\Longrightarrow\; \\
  1 - \frac{2p_b - p_a}{b-a} = 0 \\
  \;\Longrightarrow\; \\
  p_b(p_a) = \frac{(b-a) + p_a}{2}.
\end{gathered}
\]

\end{solution}

\part Encuentre el equilibrio de Nash con las funciones de reacción que encontró en el ítem anterior.
\needspace{5\baselineskip}
\begin{solution}
Reemplazando: $p_a = \frac{1}{2}\cdot\frac{(b-a) + p_a}{2} = \frac{(b-a) + p_a}{4}$, de donde $\frac{3}{4}p_a = \frac{b-a}{4}$ y
\[
  p_a^* = \frac{b - a}{3}, \qquad p_b^* = \frac{(b-a) + \frac{b-a}{3}}{2} = \frac{2(b-a)}{3}.
\]
La empresa con más alta calidad cobra un precio superior, pero ambas empresas colocan un precio por encima del costo marginal. Es decir, relajaron la competencia diferenciándose, por lo que obtendrán beneficios.

\end{solution}

\part Encuentre la demanda y los beneficios en el equilibrio.
\needspace{5\baselineskip}
\begin{solution}
La demanda depende del individuo indiferente, que a su vez depende de los precios dadas las ubicaciones:
\[
  \hat\theta = \frac{p_b^* - p_a^*}{b-a} = \frac{\frac{2}{3}(b-a) - \frac{1}{3}(b-a)}{b-a} = \frac{1}{3}.
\]
Por tanto, la demanda de $a$ es $\frac{1}{3}$ y la demanda de $b$ es $\frac{2}{3}$. Los beneficios son
\[
  \pi_a^* = p_a^*\cdot\frac{1}{3} = \frac{b-a}{9}, \qquad \pi_b^* = p_b^*\cdot\frac{2}{3} = \frac{4(b-a)}{9}.
\]

\end{solution}

\part Encuentre las ubicaciones que le conviene a las firmas para maximizar sus beneficios. ¿Le conviene centrarse o diferenciarse?
\needspace{5\baselineskip}
\begin{solution}
En el primer turno, dados los beneficios en función de las ubicaciones, se obtiene la ubicación de cada empresa que maximiza los beneficios. Para la firma de menor calidad, $\max_a \pi_a = \frac{b-a}{9}$: la derivada es negativa, por lo que conviene producir con la menor calidad posible, $a = 0$. Para la firma de mayor calidad, $\max_b \pi_b = \frac{4(b-a)}{9}$ es creciente en $b$, por lo que conviene $b = 1$.

Las firmas se diferenciarán lo máximo posible para maximizar sus beneficios: una firma producirá la mayor calidad posible y la otra la menor calidad posible, incluso aunque sean igual de baratas de producir ($c = 0$).

\end{solution}
\end{parts}

%-------------------------------------------------------------------- Coca-Cola y Pepsi
\needspace{8\baselineskip}
\tema{Competencia entre Coca-Cola y Pepsi (fusión, variante numérica)}
\question Suponga que Coca-Cola y Pepsi son los únicos competidores en el mercado de bebidas gaseosas. Ambas empresas venden productos diferenciados horizontalmente y compiten en precios. Las funciones de demanda de Coca-Cola y Pepsi son
\begin{align*}
  q_C &= 3 - 3p_C + 2p_P \\
  q_P &= 3 + 2p_C - 3p_P
\end{align*}
Los costos marginales para ambas firmas son constantes e iguales a $c_C = c_P = \frac{1}{3}$.
\begin{parts}
\part Encuentre los precios, cantidades y beneficios de cada firma en el equilibrio de Nash en precios. Grafique las funciones de mejor respuesta y el equilibrio.
\needspace{5\baselineskip}
\begin{solution}
Coca-Cola maximiza $\pi_C = \bigl(p_C - \frac{1}{3}\bigr)(3 - 3p_C + 2p_P)$, con
\[
  \frac{\partial\pi_C}{\partial p_C} = (3 - 3p_C + 2p_P) - 3\Bigl(p_C - \frac{1}{3}\Bigr) = 4 - 6p_C + 2p_P = 0
  \;\Longrightarrow\;
  p_C(p_P) = \frac{2}{3} + \frac{1}{3}p_P .
\]
De manera similar (y por simetría), $p_P(p_C) = \frac{2}{3} + \frac{1}{3}p_C$. En el equilibrio simétrico, $p = \frac{2}{3} + \frac{1}{3}p$, por lo que
\[
  p_C^* = p_P^* = 1, \qquad q_C^* = q_P^* = 2, \qquad \pi_C^* = \pi_P^* = \frac{4}{3}.
\]

\end{solution}

\part Suponga que Coca-Cola y Pepsi deciden fusionarse y formar Pepsi-Cola, que se comportará como un monopolista multiproducto (no hay sinergias como consecuencia de la fusión). Encuentre los precios, cantidades y beneficios del monopolista multiproducto.
\needspace{5\baselineskip}
\begin{solution}
Pepsi-Cola maximiza $\bigl(p_C - \frac{1}{3}\bigr)(3 - 3p_C + 2p_P) + \bigl(p_P - \frac{1}{3}\bigr)(3 - 3p_P + 2p_C)$. Derivando respecto de $p_C$:
\[
  (3 - 3p_C + 2p_P) - 3\Bigl(p_C - \frac{1}{3}\Bigr) + 2\Bigl(p_P - \frac{1}{3}\Bigr) = \frac{10}{3} - 6p_C + 4p_P = 0 .
\]
Usando simetría, $\frac{10}{3} - 2p = 0$, de donde
\[
  p_C^* = p_P^* = \frac{5}{3}, \qquad q_C^* = q_P^* = \frac{4}{3}, \qquad \pi_C^* = \pi_P^* = \frac{16}{9}, \qquad \pi_{C+P} = \frac{32}{9}.
\]

\end{solution}

\part Explique por qué las firmas, al actuar en forma independiente en (a), obtienen menos beneficios que un monopolista multiproducto.
\needspace{5\baselineskip}
\begin{solution}
Cuando las firmas toman decisiones en forma no cooperativa se encuentran en una situación parecida al dilema del prisionero: existen precios más altos que dejan a ambas firmas mejor, pero a esos precios una firma puede bajar precios unilateralmente y obtener mayores beneficios (desviación unilateral beneficiosa).

Otra forma de verlo es comparar las condiciones de primer orden del monopolista multiproducto (que por definición maximiza los beneficios conjuntos) y del duopolio. Cuando Coca-Cola fija su precio no tiene en cuenta el efecto de su precio sobre la demanda de Pepsi; por lo tanto, no considera que si rebaja su precio está afectando los beneficios de Pepsi (externalidad negativa). Pepsi-Cola sí tiene en cuenta este efecto y, por lo tanto, fija precios mayores.

\end{solution}

\part ¿Los consumidores prefieren que las firmas compitan (como en (a)) o que exista un monopolista (como en (b))? ¿Su respuesta sería distinta si los bienes fuesen complementarios en lugar de sustitutos? Explique.
\needspace{5\baselineskip}
\begin{solution}
Los consumidores prefieren que las firmas compitan porque los precios de ambos productos son menores ($1 < \frac{5}{3}$) y por ende el excedente del consumidor es mayor. Si los bienes fuesen complementarios, los precios fijados por el monopolista multiproducto serían menores, y con bienes complementarios los consumidores prefieren un monopolio a que las firmas compitan: cuando Coca-Cola disminuye su precio genera una externalidad positiva sobre Pepsi, que vende un bien complementario, y Pepsi-Cola internaliza dicho efecto y fija precios menores que las firmas en forma independiente.

\end{solution}
\end{parts}

%-------------------------------------------------------------------- Comente: logit vs nested logit
\needspace{8\baselineskip}
\tema{Comente: logit vs nested logit}
\question \textbf{Comente.} Explique la principal diferencia de los modelos de estimación de demanda logit y nested logit.
\needspace{5\baselineskip}
\begin{solution}
La principal diferencia es la elasticidad propia y cruzada que se infieren de los modelos. El modelo logit asume que los consumidores eligen productos y tipos o segmentos de productos de manera independiente, es decir, que las preferencias de los consumidores son i.i.d. Esto hace que el modelo logit entregue elasticidades y patrones de sustitución poco realistas (marcas con la misma participación de mercado tienen el mismo efecto, lo que no es cierto en la realidad, ya que no todos los productos son sustitutos para el consumidor). El modelo \emph{nested logit} mejora las elasticidades estimadas al incorporar en la función de utilidad grupos de marcas, donde las marcas son sustitutos dentro del grupo, pero no entre grupos. Por lo tanto, las elasticidades y los patrones de sustitución del modelo nested logit son más realistas.

\end{solution}

%-------------------------------------------------------------------- Comente: logit
\needspace{8\baselineskip}
\tema{Comente: patrones de sustitución en el modelo logit}
\question \textbf{Comente.} En el modelo logit todos los productos son sustitutos perfectos entre sí, lo que genera patrones de sustitución poco realistas cuando hay muchos productos en el mercado.
\needspace{5\baselineskip}
\begin{solution}
Verdadero. El modelo logit estándar asume que las preferencias de los consumidores siguen una distribución de valor extremo tipo I, lo que implica que todos los productos son igualmente sustitutos entre sí. Esto genera la propiedad de independencia de alternativas irrelevantes (IIA), donde la razón de probabilidades de elegir entre dos productos no depende de la presencia de otros productos. En mercados con muchos productos, esto produce elasticidades cruzadas poco realistas, ya que un aumento en el precio de un producto redistribuye la demanda proporcionalmente entre todos los demás productos, independientemente de qué tan similares sean.

\end{solution}

%-------------------------------------------------------------------- Comente: nested logit
\needspace{8\baselineskip}
\tema{Comente: ¿resuelve el nested logit los problemas del logit?}
\question \textbf{Comente.} El modelo \emph{nested logit} resuelve completamente los problemas del modelo logit al permitir diferentes grados de sustitución entre productos.
\needspace{5\baselineskip}
\begin{solution}
Falso. El \emph{nested logit} mejora las elasticidades estimadas al incorporar grupos de marcas (los bienes son más sustitutos dentro de un grupo que entre grupos), relajando la IIA a nivel de grupos. Sin embargo, no resuelve completamente los problemas: dentro de cada grupo se mantiene la IIA y la estructura de grupos debe ser impuesta por el investigador; para patrones de sustitución más flexibles se requieren modelos con coeficientes aleatorios.

\end{solution}

%-------------------------------------------------------------------- Heladerías
\needspace{8\baselineskip}
\tema{Hotelling con costos de transporte lineales: las heladerías}
\question Dos heladerías compiten en una playa de 10 km de longitud. La Heladería 1 se ubica en el km 1 y la Heladería 2 en el km 5. Los 500 consumidores se distribuyen uniformemente a lo largo de la playa. Cada consumidor valora el helado en $v = 50$ pesos y enfrenta un costo de transporte de $t = 3$ pesos por kilómetro recorrido. Ambas heladerías tienen un costo marginal constante $c = 3$ pesos por helado.
\begin{parts}
\part Encuentre la función de demanda de cada heladería en función de los precios $p_1$ y $p_2$.
\needspace{5\baselineskip}
\begin{solution}
Primero encontramos al consumidor indiferente, ubicado en la posición $x$. Este consumidor obtiene la misma utilidad comprando en cualquier heladería:
\[
  v - p_1 - 3(x - 1) = v - p_2 - 3(5 - x)
  \;\Longrightarrow\;
  6x = p_2 - p_1 + 18
  \;\Longrightarrow\;
  x = 3 + \frac{p_2 - p_1}{6}.
\]
Las funciones de demanda son
\[
  q_1 = 500\cdot\frac{x}{10} = 50x = 150 + \frac{25}{3}(p_2 - p_1), \qquad
  q_2 = 500\cdot\frac{10 - x}{10} = 350 - \frac{25}{3}(p_2 - p_1).
\]

\end{solution}

\part Encuentre el equilibrio de Nash en precios y calcule los beneficios de cada heladería.
\needspace{5\baselineskip}
\begin{solution}
Heladería 1: $\max_{p_1}\ (p_1 - 3)\bigl[150 + \frac{25}{3}(p_2 - p_1)\bigr]$, con CPO $175 + \frac{25}{3}p_2 - \frac{50}{3}p_1 = 0$, es decir, $p_1 = 10{,}5 + 0{,}5p_2$.

Heladería 2: $\max_{p_2}\ (p_2 - 3)\bigl[350 - \frac{25}{3}(p_2 - p_1)\bigr]$, con CPO $375 - \frac{50}{3}p_2 + \frac{25}{3}p_1 = 0$, es decir, $p_2 = 22{,}5 + 0{,}5p_1$.

Resolviendo el sistema:
\[
  p_1^* = 29, \qquad p_2^* = 37, \qquad q_1^* = 150 + \tfrac{25}{3}\cdot 8 \approx 216{,}7, \qquad q_2^* \approx 283{,}3,
\]
\[
  \pi_1^* = (29 - 3)\cdot 216{,}7 \approx 5633, \qquad \pi_2^* = (37 - 3)\cdot 283{,}3 \approx 9633 .
\]

\end{solution}

\part Si ambas heladerías estuvieran ubicadas en el mismo punto (km 5), ¿cuál sería el equilibrio en precios?
\needspace{5\baselineskip}
\begin{solution}
Si ambas heladerías estuvieran en la misma ubicación, los productos serían homogéneos y la competencia en precios sería à la Bertrand: ambas cobran el precio igual al costo marginal, $p_1 = p_2 = c = 3$, y los beneficios son cero.

\end{solution}

\part Compare los resultados de (b) y (c). ¿Qué efecto tiene la diferenciación horizontal en los precios y beneficios?
\needspace{5\baselineskip}
\begin{solution}
La diferenciación horizontal permite a las empresas relajar la competencia en precios. Cada heladería tiene cierto poder de mercado local debido a los costos de transporte, lo que les permite cobrar precios sustancialmente por encima del costo marginal y obtener beneficios positivos. Sin diferenciación, la competencia perfecta elimina todos los beneficios.

\end{solution}

\part ¿Cómo afectaría una reducción en el costo de transporte $t$ a los precios de equilibrio? Explique la intuición económica.
\needspace{5\baselineskip}
\begin{solution}
Si el costo de transporte $t$ disminuye, los productos se vuelven más sustitutos entre sí: los consumidores estarían más dispuestos a viajar hacia la heladería que ofrece menores precios. Los precios de equilibrio serían menores, la competencia se intensificaría y los beneficios de ambas heladerías disminuirían. En el límite, cuando $t \to 0$, convergeríamos al caso sin diferenciación: $p_1 = p_2 = c = 3$. El costo de transporte actúa como una ``barrera natural'' que otorga poder de mercado local a cada heladería.

\end{solution}
\end{parts}

\end{unidad}

\unidadtitulo{Monopolio multiproducto}
\begin{unidad}
%-------------------------------------------------------------------- Lerner
\needspace{8\baselineskip}
\tema{Comente: índice de Lerner multiproducto}
\question \textbf{Comente.} El índice de Lerner corresponde a una medida de poder de mercado, que indica cuánto puede cobrar una empresa por sobre su costo marginal. Este siempre se puede expresar como el inverso de la elasticidad precio-demanda por el bien.
\needspace{5\baselineskip}
\begin{solution}
Falso. Si bien la definición del índice es correcta, pues muestra cuánto se puede desviar una firma del equilibrio de competencia perfecta ($P=c$), no siempre se cumple que sea el inverso de la elasticidad. Generalmente es correcto decir que $L = \frac{1}{|\varepsilon|}$, pero existen casos en que esto no se cumple, como por ejemplo un monopolio multiproducto. En este caso, como la demanda de cada bien depende del precio del otro, el índice es distinto, pues el monopolista internaliza las decisiones que toma en ambos bienes para no generar pérdidas en sus beneficios. De la CPO respecto de $p_1$, $D_1 + (p_1-c_1)\frac{\partial D_1}{\partial p_1} + (p_2-c_2)\frac{\partial D_2}{\partial p_1} = 0$, se obtiene
\[
  L_1 = \frac{p_1 - c_1}{p_1} = \frac{1}{|\varepsilon_{11}|} - \frac{p_2 - c_2}{p_1}\cdot\frac{\partial D_2/\partial p_1}{\partial D_1/\partial p_1}.
\]
Cuando las demandas de los bienes del monopolio multiproducto están relacionadas, la regla de la elasticidad inversa no se cumple. En detalle:
\begin{itemize}
  \item Si los bienes son \emph{sustitutos} ($\partial D_2/\partial p_1 > 0$), entonces $L^{\text{Multiproducto}}_1 > \frac{1}{|\varepsilon|} = L^{\text{Uniproducto}}_1$.
  \item Si los bienes son \emph{complementarios} ($\partial D_2/\partial p_1 < 0$), entonces $L^{\text{Multiproducto}}_1 < \frac{1}{|\varepsilon|} = L^{\text{Uniproducto}}_1$.
\end{itemize}

\end{solution}

%-------------------------------------------------------------------- Bienes diferenciados
\needspace{8\baselineskip}

%-------------------------------------------------------------------- Sillas y mesas
\needspace{8\baselineskip}
\tema{Monopolio multiproducto: sillas y mesas}
\question En el país «Muy muy lejano» existe un solo productor de sillas ($q_1$) y mesas ($q_2$) que abastece a todo el mercado. Las funciones de demanda de ambos bienes son
\begin{align*}
  q_1 &= 200 - 4p_1 - p_2 \\
  q_2 &= 180 - 3p_2 - p_1
\end{align*}
Adicionalmente se sabe que los costos de producción son separables, con funciones $CT_1 = 4q_1 + 200$ y $CT_2 = 2q_2 + 500$.
\begin{parts}
\part ¿Qué tipo de bienes son las sillas y las mesas?
\needspace{5\baselineskip}
\begin{solution}
Debemos derivar cruzado:
\[
  \frac{\partial q_1}{\partial p_2} = -1 < 0, \qquad \frac{\partial q_2}{\partial p_1} = -1 < 0 .
\]
Dado que las derivadas cruzadas son negativas, son bienes \emph{complementarios}: un aumento en el precio de las sillas induce una caída en la demanda de mesas (y viceversa).

\end{solution}

\part Encuentre el equilibrio del monopolio para ambos mercados.
\needspace{5\baselineskip}
\begin{solution}
Resolvemos el problema de maximización:
\[
  \max_{p_1,p_2}\ \pi = p_1(200 - 4p_1 - p_2) + p_2(180 - 3p_2 - p_1) - (4q_1 + 200) - (2q_2 + 500).
\]
CPO:
\begin{gather*}
  \frac{\partial \pi}{\partial p_1} = 218 - 8p_1 - 2p_2 = 0, \\
  \frac{\partial \pi}{\partial p_2} = 190 - 2p_1 - 6p_2 = 0.
\end{gather*}
Despejando $p_2$ de la primera, $p_2 = \frac{218 - 8p_1}{2}$, y reemplazando en la segunda se obtiene
\[
  p_1^* \approx 21{,}09, \qquad p_2^* \approx 24{,}64, \qquad q_1^* = 91, \qquad q_2^* = 85 .
\]

\end{solution}

\part Obtenga el beneficio de este monopolio multiproducto. Además, mencione cómo cambiarían los beneficios si el monopolista actuara de forma independiente en cada mercado.
\needspace{5\baselineskip}
\begin{solution}
\[
  \pi^M = 21{,}09\cdot 91 + 24{,}64\cdot 85 - (4\cdot 91 + 200) - (2\cdot 85 + 500) \approx 2779 .
\]
Si el monopolio decide independientemente el precio de cada bien (ignorando el efecto sobre el otro mercado), las CPO pasan a ser $216 - 8p_1 - p_2 = 0$ y $186 - p_1 - 6p_2 = 0$, con $p_1 \approx 23{,}62$, $p_2 \approx 27{,}06$ y un beneficio total de $\approx 2724 < 2779$.

El beneficio disminuye porque los bienes son complementarios. Cuando un monopolio posee los mercados de 2 bienes complementarios, el producir de manera conjunta le permite disminuir el precio de un bien más allá de lo que haría de forma independiente, con el objetivo de incentivar la demanda del otro bien, maximizando así su beneficio total.

\end{solution}
\end{parts}

%-------------------------------------------------------------------- Comente: monopolista nuevos productos
\needspace{8\baselineskip}
\tema{Comente: ¿conviene a un monopolista introducir nuevos productos?}
\question \textbf{Comente.} A un monopolista no le conviene introducir nuevos productos al mercado, pues se crea competencia con el mismo.
\needspace{5\baselineskip}
\begin{solution}
Falso. El monopolista introduce nuevos productos ya sea porque las demandas están relacionadas o porque tiene reducción de costos. Si las demandas están relacionadas, le conviene introducir marcas nuevas, es decir, bienes sustitutos, y así aprovechar la demanda de ambos bienes ante cambios de precios (internalizando que sus bienes compiten entre sí). El monopolio intertemporal es un caso de monopolio multiproducto, donde el monopolio debe considerar que es competencia consigo mismo en el tiempo y por tanto elige los precios de ambos períodos para maximizar la utilidad conjunta.

\end{solution}

\end{unidad}

\unidadtitulo{Concentración, poder de mercado y fusiones}
\begin{unidad}
%-------------------------------------------------------------------- Comente: IHH
\needspace{8\baselineskip}
\tema{Comente: IHH y número de firmas}
\question \textbf{Comente.} El índice de Herfindahl-Hirschman (IHH) es el indicador más empleado para medir concentración en una industria. Si en un mercado existen muchas firmas, el índice IHH toma un valor cercano a cero y por lo tanto no existiría concentración de mercado.
\needspace{5\baselineskip}
\begin{solution}
Falso. El IHH no solo mide el número de firmas en la industria para determinar si existe concentración de mercado: además mide la dispersión en la producción de cada firma (qué tan asimétrica es la estructura de mercado). Matemáticamente, el índice se puede definir como
\[
  IHH(N) = \sum_{i=1}^{N}s_i^2 = \frac{1}{N} + N\sigma^2,
\]
donde $\sigma^2 = Var(s_i)$ es la varianza de las participaciones de mercado, por lo que el indicador aumenta a medida que hay mayor dispersión en la producción de cada firma. Así, aun con muchas firmas, una muy grande y muchas pequeñas puede dar un IHH alto.

\end{solution}

%-------------------------------------------------------------------- Comente: aumento del IHH
\needspace{8\baselineskip}
\tema{Comente: aumento del IHH y bienestar de los consumidores}
\question \textbf{Comente.} Un aumento en el Índice de Herfindahl-Hirschman (IHH) indica que hubo una disminución de firmas y de intensidad competitiva en el mercado. Esto perjudica a los consumidores ya que implica un aumento en el precio del bien o servicio de la industria en cuestión.
\needspace{5\baselineskip}
\begin{solution}
No necesariamente. En primer lugar, el índice puede aumentar debido a una mayor asimetría entre las firmas, y no solo a que hayan menos. Además, que la concentración de mercado haya aumentado no siempre implica una menor intensidad competitiva, ya que esto también dependerá de la elasticidad precio de la demanda. Un aumento en el IHH puede provocarse debido a la fusión de un número de firmas. Si tal fusión conlleva ganancias de eficiencia lo suficientemente altas, entonces el precio, aun cuando el mercado esté más concentrado, disminuirá, beneficiando a los consumidores. Otro ejemplo que contradice la afirmación es que empresas innovadoras o más eficientes aumenten su participación de mercado y, a la vez, aumenten su competencia, llevando a un menor precio para los consumidores.

\end{solution}

%-------------------------------------------------------------------- Comente: IHH y Lerner
\needspace{8\baselineskip}
\tema{Comente: IHH e índice de Lerner en Cournot}
\question \textbf{Comente.} Para el modelo de Cournot, la relación entre el índice de Herfindahl-Hirschman (IHH) y el índice de Lerner establece que una firma tiene alto poder de mercado si la elasticidad precio de la demanda es alta cuando la industria está concentrada.
\needspace{5\baselineskip}
\begin{solution}
Falso. En el modelo de Cournot, al juntar el IHH ($\sum s_i^2$) con el índice de Lerner individual ($L_i = \frac{s_i}{|\varepsilon|}$) se obtiene la siguiente relación:
\[
  \bar L = \sum_{i=1}^{N}s_i L_i = \sum_{i=1}^{N}\frac{s_i^2}{|\varepsilon|} = \frac{IHH}{|\varepsilon|},
\]
donde $L_i$ y $\bar L$ son el índice de Lerner individual y agregado, respectivamente. Esto muestra que, pese a que una industria esté muy concentrada (alto IHH), una firma podría tener bajo poder de mercado si la demanda es muy elástica. En otras palabras, una firma tendría bajo poder de mercado cuando: (i) está en un mercado muy atomizado y/o (ii) enfrenta una demanda muy elástica. Lo correcto es entonces que el poder de mercado es alto cuando la elasticidad es \emph{baja} y la industria está concentrada.

\end{solution}

%-------------------------------------------------------------------- Comente: prohibir fusiones
\needspace{8\baselineskip}
\tema{Comente: las fusiones deben prohibirse}
\question \textbf{Comente.} Las fusiones deben ser prohibidas porque solo empeoran a los consumidores.
\needspace{5\baselineskip}
\begin{solution}
Falso. Si bien en algunos casos pueden generar concentraciones en mercados que pueden dañar a los consumidores, existen fusiones que pueden generar una ganancia de eficiencia tal que el precio baje, la cantidad aumente y los consumidores se vean beneficiados. Por lo tanto, las fusiones no siempre empeoran a los consumidores.

\end{solution}

%-------------------------------------------------------------------- Comente: participación de mercado y fusión
\needspace{8\baselineskip}
\tema{Comente: fusiones entre firmas grandes}
\question Explique por qué razón las fusiones entre empresas con mayor participación de mercado a priori se consideran que tienen mayor posibilidad de incrementar los precios que aquellas entre firmas con menor participación de mercado. Emplee un modelo si es necesario.
\needspace{5\baselineskip}
\begin{solution}
Si las firmas venden un producto homogéneo, es decir, sin diferenciación, aquellas de menor costo tendrán mayor participación de mercado. Por lo tanto, una fusión entre firmas de mayor participación de mercado implica que las de menores costos se fusionan y, por lo tanto, la principal firma competidora de otra deja de competir, razón por la cual los precios tienden a subir más. Se puede emplear un modelo de Cournot o Bertrand para demostrar el resultado.

\end{solution}

%-------------------------------------------------------------------- Fusión Cournot
\needspace{8\baselineskip}
\tema{Fusión en competencia tipo Cournot (zapatos de fútbol)}
\question Recientemente, con el objetivo de ganar la Champions League, Arturo Vidal fichó por el FC Barcelona. Y para estar lo más preparado posible y enfocarse solo en entrenar, le delega a usted la administración de su firma de zapatos de fútbol. En este mercado solo existen 3 firmas que compiten en cantidad (incluida la de King Arturo) y que tienen la misma función de costos totales: $C(q) = 20q$. Además se sabe que la demanda está dada por $Q(p) = 200 - p$.
\begin{parts}
\part Una de las firmas competidoras le ofrece realizar una fusión. Demuestre que lo más conveniente es rechazar la oferta.
\needspace{5\baselineskip}
\begin{solution}
Para evaluar si conviene o no la oferta, se debe comparar la situación de la firma pre y post fusión.

\textbf{Pre-fusión.} Cada firma maximiza sus beneficios, $\max_{q_i}\ \pi_i = (200 - q_i - q_j - q_k - 20)q_i$, de donde se obtiene la función de reacción
\[
  q_i = \frac{180 - q_j - q_k}{2}.
\]
En equilibrio simétrico: $q_i = q_j = q_k = \frac{180}{4} = 45$, $Q = 135$, $P = 65$. Los beneficios de cada firma son $\pi_i = (65-20)\cdot 45 = 2025$.

\textbf{Post-fusión.} La firma fusionada ($q_f = q_i + q_k$) y la otra firma resuelven $\max_{q_f}\ (200 - q_f - q_j - 20)q_f$, con funciones de reacción $q_f = \frac{180 - q_j}{2}$ y $q_j = \frac{180 - q_f}{2}$. En equilibrio,
\[
\begin{gathered}
  q_f = q_j = \frac{180}{3} = 60, \quad Q = 120, \quad P = 80, \\ \pi_f = \pi_j = (80-20)\cdot 60 = 3600 .
\end{gathered}
\]
Comparando el beneficio de la empresa fusionada con la suma de sus beneficios antes de la fusión:
\[
  \pi_f = 3600 < \pi_i + \pi_k = 2025 + 2025 = 4050 .
\]
Dada esta estructura de costos y la competencia en cantidades, no es conveniente la fusión. Por lo que lo mejor es rechazar la oferta.

\end{solution}

\part Ahora suponga que la función de costos de las firmas cambia a $C(q) = 20q + q^2$. Analice si en este caso es conveniente o no la fusión.
\needspace{5\baselineskip}
\begin{solution}
Igual que antes, comparamos la situación de la firma pre y post fusión.

\textbf{Pre-fusión.} Cada firma maximiza $\pi_i = (200 - q_i - q_j - q_k)q_i - 20q_i - q_i^2$, con función de reacción $q_i = \frac{180 - q_j - q_k}{4}$. En equilibrio:
\[
  q_i = q_j = q_k = \frac{180}{6} = 30, \qquad Q = 90, \qquad P = 110,
\]
\[
  \pi_i = 110\cdot 30 - 20\cdot 30 - 30^2 = 3300 - 600 - 900 = 1800 .
\]

\textbf{Post-fusión.} La firma fusionada distribuye su producción entre sus dos plantas según $CMg_i(q_i) = CMg_k(q_k)$, es decir $20 + 2q_i = 20 + 2q_k$, por lo que $q_i = q_k = \frac{q_{ik}}{2}$. La función de costos de la firma fusionada es entonces
\[
  C_{ik}(q_{ik}) = 2\Bigl[20\frac{q_{ik}}{2} + \Bigl(\frac{q_{ik}}{2}\Bigr)^2\Bigr] = 20q_{ik} + \frac{q_{ik}^2}{2}.
\]
Notar que la firma fusionada presenta una ventaja en su estructura de costos con respecto a la otra firma: $C_{ik}(q) < C(q)$. Ahora los problemas de maximización son distintos:
\[
  \max_{q_{ik}}\ (200 - q_{ik} - q_j)q_{ik} - 20q_{ik} - \frac{q_{ik}^2}{2}, \qquad
  \max_{q_j}\ (200 - q_{ik} - q_j)q_j - 20q_j - q_j^2 .
\]
Las funciones de reacción son
\[
  q_{ik} = \frac{180 - q_j}{3}, \qquad q_j = \frac{180 - q_{ik}}{4}.
\]
Reemplazando una en la otra, se obtiene el equilibrio
\[
  q_{ik} = \frac{540}{11} \approx 49{,}1, \qquad q_j = \frac{360}{11} \approx 32{,}7, \qquad Q \approx 81{,}8, \qquad P \approx 118{,}2 .
\]
Los beneficios son $\pi_{ik} = (P-20)q_{ik} - \frac{q_{ik}^2}{2} \approx 3615$ y $\pi_j = (P-20)q_j - q_j^2 \approx 2142$. Comparando:
\[
  \pi_{ik} \approx 3615 > \pi_i + \pi_k = 1800 + 1800 = 3600 .
\]
Ahora sí conviene aceptar la oferta de fusión. Pero notar que los consumidores están peor que antes de la fusión, pues el precio aumentó ($110 \to 118{,}2$) y la cantidad del bien disminuyó ($90 \to 81{,}8$).

\end{solution}
\end{parts}

%-------------------------------------------------------------------- Fusión Bertrand
\needspace{8\baselineskip}
\tema{Fusión en competencia tipo Bertrand (almacenes)}
\question En el mercado de los almacenes existen 3 firmas: Siette, Buen Market y Próximo, las cuales tienen los siguientes costos marginales: $c_S = c_B = 1$ y $c_P = 2$, respectivamente. Las firmas compiten en precios y ofrecen un producto homogéneo. Además, usted sabe que la demanda está dada por la función $Q(P) = 10 - P$.
\begin{parts}
\part Siette está evaluando realizar una fusión con Buen Market para así obtener más beneficios. Evalúe si le conviene o no la fusión.
\needspace{5\baselineskip}
\begin{solution}
Para saber si a Siette le conviene o no la fusión, se deben comparar los beneficios de dicha empresa antes y después de la fusión.

\textbf{Antes de la fusión.} Dado que las empresas compiten en precios, el precio de equilibrio será el costo marginal de la segunda empresa más eficiente, es decir, $P = c_S = c_B = 1$. La firma Próximo queda fuera del mercado al ser su costo marginal mayor al precio de equilibrio ($c_P = 2 > P = 1$). Por lo tanto, en equilibrio:
\[
  Q = 10 - 1 = 9, \qquad q_S = q_B = \frac{9}{2} = 4{,}5, \qquad \pi_S = \pi_B = 0 .
\]

\textbf{Después de la fusión.} Dado que se fusionan las firmas más eficientes y compiten contra la menos eficiente, el precio de equilibrio estará dado por $P = c_P - \epsilon \approx 2$, dejando fuera a Próximo. Por lo tanto, en equilibrio:
\[
  Q = q_{SB} = 10 - 2 = 8, \qquad \pi_{SB} = (2-1)\cdot 8 = 8 .
\]
Al comparar los beneficios antes y después de la fusión, vemos que son mayores con la fusión. Por ende, a Siette le conviene fusionarse con Buen Market.

\end{solution}

\part ¿Cómo cambia el bienestar del consumidor con esta fusión entre Siette y Buen Market?
\needspace{5\baselineskip}
\begin{solution}
Hay que calcular el excedente del consumidor antes y después de la fusión.

\textbf{Antes:} $P = 1$ y $Q = 9$, por lo que $EC_{\text{antes}} = \frac{(10-1)\cdot 9}{2} = \frac{81}{2} = 40{,}5$.

\textbf{Después:} $P = 2$ y $Q = 8$, por lo que $EC_{\text{después}} = \frac{(10-2)\cdot 8}{2} = \frac{64}{2} = 32$.

El excedente después de la fusión es menor que antes de esta, por lo que el bienestar del consumidor disminuye.

\end{solution}
\end{parts}

%-------------------------------------------------------------------- HHI cerveza
\needspace{8\baselineskip}
\tema{Cálculo del HHI: industria de la cerveza}
\question Suponga que usted es nombrado miembro del Tribunal de Defensa de la Libre Competencia. El primer caso que enfrenta es la fusión de las dos empresas más grandes de la industria de la cerveza: CCU tiene el 50\% del mercado, Cervecerías Chile un 30\%, y el 20\% restante está repartido en partes iguales en las 4 empresas restantes. Calcule el HHI.
\needspace{5\baselineskip}
\begin{solution}
El índice de Herfindahl-Hirschman se calcula como
\[
  HHI = 10\,000\sum_{i=1}^{n}s_i^2 = 10\,000\bigl[(0{,}5)^2 + (0{,}3)^2 + 4\cdot(0{,}05)^2\bigr] = 10\,000\cdot 0{,}35 = 3500 .
\]
(Si CCU y Cervecerías Chile se fusionaran, el HHI pasaría a $10\,000[(0{,}8)^2 + 4(0{,}05)^2] = 6500$, un aumento de $3000$ puntos.)

\end{solution}

%-------------------------------------------------------------------- Aprobar fusión
\needspace{8\baselineskip}
\tema{Comente: ¿debe aprobarse la fusión CCU--Cervecerías Chile?}
\question \textbf{Comente.} Suponga que CCU se quiere fusionar con Cervecerías Chile. ¿Cree usted que debería aprobarse la fusión? Argumente su respuesta.
\needspace{5\baselineskip}
\begin{solution}
En general, se argumenta que una fusión traería economías de escala y sinergias entre las empresas, y además ahorro de costos y publicidad. Para que la fusión sea conveniente debe aumentar el bienestar social agregado de los consumidores y de los productores. Si ambas empresas están dispuestas a fusionarse, sus utilidades deberían aumentar (sea por el motivo que sea). En el caso de los consumidores, estos se verán beneficiados si bajan los precios y tienen acceso a mejor y más barata cerveza. Para aprobar la fusión se debe considerar si en el mercado existen o no barreras a la entrada. Si hay barreras de entrada, puede no aprobarse la fusión, ya que la nueva empresa puede comportarse como un monopolio.

\end{solution}

%-------------------------------------------------------------------- Smartphones
\needspace{8\baselineskip}
\tema{Comente: fusión Samsung--Apple y precios de los rivales}
\question \textbf{Comente.} Las cuotas de mercado para \emph{smartphones} en el mercado mundial son: Samsung tiene un 20,9\%, Apple un 11,7\%, Huawei un 12,7\%, Lenovo un 4,2\%, LG un 4,8\%, y otros un 45,7\%. Suponga que Samsung y Apple deciden fusionarse. ¿Qué espera que suceda con los precios del resto de los vendedores de \emph{smartphones} (Huawei, Lenovo y LG)? ¿Su respuesta sería distinta si la fusión permitiese importantes ahorros en los costos de producción? Fundamente su respuesta.
\needspace{5\baselineskip}
\begin{solution}
Sin ahorro de costos, se espera que los rivales que no se fusionan suban sus precios. La razón es que las firmas fusionadas aumentan sus precios (al internalizar las elasticidades precio cruzadas) y, dado que los precios de bienes sustitutos son complementos estratégicos, los rivales tenderán a subir los suyos. Con ahorro de costos, el resultado no es obvio: si los ahorros de costos son importantes, la empresa fusionada baja los precios y sus rivales tienden a bajar los suyos (complementos estratégicos).

\end{solution}

%-------------------------------------------------------------------- Fusión transporte aéreo
\needspace{8\baselineskip}
\tema{Fusiones en el transporte aéreo}
\question Las empresas de transporte aéreo $X$, $Y$ y $Z$ son las únicas que operan en la ruta Santiago--Lima. La demanda semanal por transporte es igual a $Q = 130 - P$, donde $Q$ representa el número total de viajes semanales. La función de costo total de cada empresa es $C_i = 50q_i + 300$, con $i = X, Y, Z$, donde $q_i$ es el número de viajes semanales hechos por la empresa.
\begin{parts}
\part Debido a la alta contaminación emitida por las aerolíneas, el Ministerio de Transporte estableció, por empresa, realizar solamente 32 viajes semanales como máximo. Asumiendo que la competencia es por viaje realizado, determine el número de viajes que realiza cada empresa, el precio de equilibrio, los beneficios de cada una, y si la restricción propuesta por el Ministerio afecta a las empresas.
\needspace{5\baselineskip}
\begin{solution}
Resolviendo el problema para la firma $X$ sin imponer ninguna restricción:
\[
\begin{gathered}
  \max_{q_X}\ \pi_X = (130 - q_X - q_Y - q_Z - 50)q_X - 300 \\
  \;\Rightarrow\; 80 - 2q_X - q_Y - q_Z = 0 \\
  \;\Rightarrow\; q_X(q_Y,q_Z) = \frac{80 - q_Y - q_Z}{2}.
\end{gathered}
\]
Por simetría, $q_Y$ y $q_Z$ tienen funciones de reacción análogas. Sumando las tres: $2Q = 240 - 2Q$, es decir, $Q = 60$, el precio es $P = 130 - 60 = 70$ y cada firma produce $q_X = q_Y = q_Z = 20$. La restricción impuesta por el Ministerio de Transporte no está activa ($20 < 32$). Los beneficios son
\[
  \pi_X = \pi_Y = \pi_Z = (70 - 50)\cdot 20 - 300 = 100 .
\]

\end{solution}

\part Considere ahora que las empresas $X$ e $Y$ han decidido fusionarse, formando el grupo $XY$. La idea de las aerolíneas es poder ahorrar costos, ya que evitarían duplicar los costos fijos y también tendrían ganancias de eficiencia que les permitirían disminuir en 4\% sus costos variables. Esta disminución de costos está asociada a menores emisiones de agentes contaminantes, por lo que ya no se restringe el número de viajes. Determine las cantidades que produciría cada firma en el equilibrio, el precio de mercado y los beneficios de cada firma. Por último, determine si es rentable para las empresas llevar a cabo la fusión.
\needspace{5\baselineskip}
\begin{solution}
La firma fusionada tiene costos totales $C_{XY} = 50\cdot 0{,}96\,q_{XY} + 300 = 48q_{XY} + 300$. Las funciones de reacción son
\[
  q_Z(q_{XY}) = \frac{80 - q_{XY}}{2}, \qquad q_{XY}(q_Z) = \frac{82 - q_Z}{2}.
\]
Reemplazando $q_{XY}$ en $q_Z$: $4q_Z = 160 - 82 + q_Z$, es decir $q_Z = 26$ y $q_{XY} = 28$. La cantidad de mercado es $Q = 54$ y el precio $P = 130 - 54 = 76$. Los beneficios son
\[
  \pi_{XY} = (76 - 48)\cdot 28 - 300 = 484, \qquad \pi_Z = (76 - 50)\cdot 26 - 300 = 376 .
\]
La firma fusionada gana $484 > 100 + 100 = 200$, por lo que la fusión es rentable; y la firma $Z$, que no se fusiona, también gana más que antes ($376 > 100$).

\end{solution}

\part Si usted trabaja en la División de Fusiones de la FNE, y su labor es maximizar el bienestar social, ¿estaría de acuerdo con aprobar la fusión en los términos señalados? Justifique.
\needspace{5\baselineskip}
\begin{solution}
De los resultados anteriores es claro que al fusionarse dos firmas el precio se ve incrementado ($70 \to 76$) y la cantidad total ofrecida al mercado disminuye ($60 \to 54$), por lo que los consumidores están peor (se reduce el excedente del consumidor). Aunque hay ahorros de costos, el aumento de beneficios de las firmas no compensa la pérdida de los consumidores en términos de bienestar de los consumidores, por lo que en el estándar de excedente del consumidor no se aprobaría la fusión.

\end{solution}
\end{parts}

%-------------------------------------------------------------------- Análisis de fusión 3 firmas costos cuadráticos
\needspace{8\baselineskip}
\tema{Análisis de fusión con costos cuadráticos (tres firmas)}
\question Existe un mercado compuesto por tres firmas inicialmente idénticas que ofrecen productos homogéneos, las que poseen una función de costo total $C(q) = 0{,}5q^2$. La demanda del mercado está dada por $Q(P) = 120 - P$. La competencia es en cantidades.
\begin{parts}
\part Establezca las condiciones de primer orden de cada firma y la función de reacción de cada una de ellas (por ejemplo, $q_1(q_2,q_3)$).
\needspace{5\baselineskip}
\begin{solution}
La función a maximizar para cada firma es $\pi_i = (120 - q_1 - q_2 - q_3)q_i - 0{,}5q_i^2$. La condición de primer orden es $120 - 2q_i - q_j - q_k - q_i = 0$, por ende la función de reacción es
\[
  q_i^* = \frac{120 - q_j - q_k}{3}.
\]

\end{solution}

\part Obtenga el equilibrio de Nash (emplee la condición de simetría entre las firmas).
\needspace{5\baselineskip}
\begin{solution}
Con simetría, $3q = 120 - 2q$, de modo que $q_1^* = q_2^* = q_3^* = 24$. El precio es $P^* = 120 - 72 = 48$ y los beneficios son $\pi_i = 48\cdot 24 - 0{,}5\cdot 24^2 = 864$.

\end{solution}

\part Suponga que las firmas 2 y 3 se fusionan, pasando ambas a ser propiedad de un único dueño. ¿Sigue el dueño de la empresa fusionada operando ambas plantas anteriores?
\needspace{5\baselineskip}
\begin{solution}
Primero, debemos calcular la nueva estructura de costos que enfrenta la firma fusionada. Para esto igualamos los costos marginales de las plantas como condición de optimalidad: $q_2 = q_3$, de modo que $q_{23} = q_2 + q_3 = 2q_2$ y cada planta produce $\frac{q_{23}}{2}$. En el caso de que decida operar con las dos plantas, los costos totales son
\[
  C(q_{23}) = 2\cdot 0{,}5\Bigl(\frac{q_{23}}{2}\Bigr)^2 = \frac{q_{23}^2}{4},
\]
mientras que operar con una sola planta costaría $0{,}5q_{23}^2$. Por lo tanto, la fusión permite una reducción en los costos de producción, y dado que no existen restricciones de capacidad, al dueño le conviene seguir operando ambas plantas (repartiendo la producción en partes iguales).

\end{solution}

\part Obtenga las condiciones de primer orden de la nueva firma fusionada y compare la función de reacción de ésta con la de las dos firmas anteriores (2 y 3). ¿Cuándo producen más las firmas 2 y 3, antes o después de la fusión?
\needspace{5\baselineskip}
\begin{solution}
La firma fusionada maximiza $\pi_{23} = (120 - q_1 - q_{23})q_{23} - \frac{q_{23}^2}{4}$. La condición de primer orden es
\[
\begin{gathered}
  120 - q_1 - 2q_{23} - \frac{q_{23}}{2} = 0 \\
  \Longrightarrow\; q_{23}(q_1) = \frac{2}{5}(120 - q_1).
\end{gathered}
\]
Antes de la fusión, y por simetría entre 2 y 3, la producción conjunta de ambas firmas dada $q_1$ era $q_2 + q_3 = \frac{120 - q_1}{2}$ (pues $q_2 = q_3 = \frac{120 - q_1}{4}$). Como $\frac{2}{5} < \frac{1}{2}$, después de la fusión las firmas 2 y 3 producen menos en conjunto que cuando actuaban por separado.

\end{solution}

\part Analíticamente, concluya si la fusión aumenta o disminuye la cantidad producida y los beneficios de la firma que no se fusiona (firma 1).
\needspace{5\baselineskip}
\begin{solution}
Para la firma no fusionada, la CPO viene de $\pi_1 = (120 - q_1 - q_{23})q_1 - 0{,}5q_1^2$, es decir $q_1 = \frac{120 - q_{23}}{3}$. Como la cantidad producida por la firma fusionada es menor que la suma de las cantidades producidas por 2 y 3 por separado, la cantidad producida por la firma 1 (no fusionada) es mayor que en el caso donde compiten las 3 firmas.

Resolviendo el sistema: $q_1 = \frac{120 - q_{23}}{3}$ y $q_{23} = \frac{2}{5}(120 - q_1)$,
\[
  q_1^* = \frac{360}{13} \approx 27{,}7, \qquad q_{23}^* = \frac{480}{13} \approx 36{,}9, \qquad Q^* = \frac{840}{13}, \qquad P^* = \frac{720}{13} \approx 55{,}4 .
\]
Los beneficios son
\[
  \pi_1 = \frac{194\,400}{169} \approx 1150 > 864, \qquad
  \pi_{23} = P^*q_{23}^* - \frac{(q_{23}^*)^2}{4} = \frac{288\,000}{169} \approx 1704 .
\]
Es decir, los beneficios de la firma no fusionada aumentan ($864 \to 1150$). Notar además que la firma fusionada gana $1704 < 2\cdot 864 = 1728$: la fusión \emph{no} es rentable para las firmas fusionadas, y sí para la rival externa (paradoja de las fusiones).

\end{solution}
\end{parts}

%-------------------------------------------------------------------- Comente: efectos de fusiones
\needspace{8\baselineskip}
\tema{Comente: efectos unilaterales y coordinados de una fusión}
\question \textbf{Comente.} Una fusión puede generar efectos negativos sobre la competencia incluso si no se llega a un acuerdo explícito entre empresas.
\needspace{5\baselineskip}
\begin{solution}
Verdadero. Las fusiones horizontales pueden disminuir la competencia por dos medios principales:
\begin{itemize}
  \item \emph{Efectos unilaterales:} aumento de precios o reducción de la calidad del producto dado el mayor poder de mercado producto de la concentración.
  \item \emph{Efectos coordinados:} la fusión puede facilitar una acomodación tácita entre las empresas que permanecen en el mercado, volviendo más atractiva la estrategia de no competir agresivamente. Esto ocurre, por ejemplo, cuando la concentración reduce el número de actores, incrementa la simetría entre ellos y mejora la transparencia del mercado, condiciones que facilitan la coordinación implícita.
\end{itemize}
También las fusiones que aumentan el poder de mercado de las firmas podrían llevar a otras consecuencias negativas, como disminución de la calidad del producto o menor innovación o inversión.

\end{solution}

%-------------------------------------------------------------------- Comente: HHI disminuye
\needspace{8\baselineskip}
\tema{Comente: el HHI disminuye con una fusión}
\question \textbf{Comente.} El índice de Herfindahl-Hirschman (HHI) disminuye si dos firmas se fusionan. Verdadero o falso.
\needspace{5\baselineskip}
\begin{solution}
Falso. Una fusión aumenta la concentración y, por tanto, aumenta el HHI, ya que este es la suma de los cuadrados de las participaciones de mercado: $(s_i + s_j)^2 = s_i^2 + s_j^2 + 2s_is_j > s_i^2 + s_j^2$ (el aumento es $2s_is_j$).

\end{solution}

%-------------------------------------------------------------------- Comente: presión al alza de precios
\needspace{8\baselineskip}
\tema{Comente: presión al alza de precios y sustitutos}
\question \textbf{Comente.} La presión al alza de precios será mayor si los productos fusionados son altamente sustitutos.
\needspace{5\baselineskip}
\begin{solution}
Verdadero. Mientras más sustitutos sean dos bienes de distintas empresas, hay una mayor competencia entre ellos; al fusionarse estas dos firmas, esta competencia se disipa completamente (se internaliza la externalidad), empujando los precios al alza.

\end{solution}

%-------------------------------------------------------------------- Comente: abuso de posición dominante
\needspace{8\baselineskip}
\tema{Comente: abuso de posición dominante sin precios monopólicos}
\question \textbf{Comente.} Una empresa puede abusar de su posición dominante sin necesidad de fijar precios monopólicos.
\needspace{5\baselineskip}
\begin{solution}
Verdadero. El abuso de posición dominante puede manifestarse de múltiples formas, como imponer condiciones contractuales restrictivas, excluir competidores mediante pagos o acuerdos de exclusividad, o manipular el acceso a canales de distribución. En el caso de Google, se observó una estrategia sistemática para mantener su dominio a través de prácticas que limitaban la competencia, sin necesidad de fijar precios altos directamente. Por ejemplo, pagar altas sumas de dinero para dejar a Chrome como el navegador por defecto en celulares nuevos.

\end{solution}

%-------------------------------------------------------------------- Comente: innovación y acceso
\needspace{8\baselineskip}
\tema{Comente: las autoridades evalúan más que precios}
\question \textbf{Comente.} Las autoridades de competencia evalúan no solo el impacto en precios, sino también en la innovación y en el acceso al mercado.
\needspace{5\baselineskip}
\begin{solution}
Verdadero. La defensa de la libre competencia busca preservar el dinamismo del mercado, no solo evitar alzas de precios. El caso europeo contra Google se basó, entre otros puntos, en que su conducta restringía la capacidad de otras empresas para innovar y competir en igualdad de condiciones, afectando a largo plazo la variedad y calidad de los servicios digitales ofrecidos.

\end{solution}

%-------------------------------------------------------------------- Comente: maverick
\needspace{8\baselineskip}
\tema{Comente: eliminación de una empresa maverick}
\question \textbf{Comente.} La eliminación de una empresa \emph{maverick} puede incrementar el riesgo de coordinación entre los actores restantes del mercado.
\needspace{5\baselineskip}
\begin{solution}
Verdadero. Una empresa \emph{maverick} es aquella que compite agresivamente y rompe posibles equilibrios de precios altos. Su eliminación por vía de una fusión puede hacer más estable un acuerdo tácito de precios altos entre las firmas que permanecen, al remover el incentivo o la capacidad para competir agresivamente.

\end{solution}

%-------------------------------------------------------------------- Fusión sin eficiencias
\needspace{8\baselineskip}
\tema{Fusión sin eficiencias en Cournot (tres firmas)}
\question Existen tres firmas idénticas que producen a un costo marginal constante igual a $c$. Suponemos una demanda lineal $P(Q) = a - Q$, donde $Q = q_1 + q_2 + q_3$. Antes de la fusión las tres firmas deciden sus cantidades en forma independiente. Después de la fusión, las firmas 2 y 3 pasan a ser propiedad de una misma empresa.
\begin{parts}
\part Determine las funciones de reacción de las tres firmas antes de la fusión y obtenga el precio de equilibrio.
\needspace{5\baselineskip}
\begin{solution}
La firma 1 maximiza $\pi_1 = (a - q_1 - q_2 - q_3 - c)q_1$, con CPO $a - 2q_1 - q_2 - q_3 - c = 0$:
\[
  q_1(q_2,q_3) = \frac{(a-c) - (q_2 + q_3)}{2},
\]
y análogamente para las demás. Por simetría, $q_i = q$ y $4q = a - c$:
\[
  q^* = \frac{a - c}{4}, \qquad P = a - \frac{3(a-c)}{4} = \frac{a + 3c}{4}.
\]

\end{solution}

\part Las firmas 2 y 3 se fusionan y deciden la producción que maximice beneficios de forma conjunta. Encuentre las funciones de reacción y el precio en equilibrio.
\needspace{5\baselineskip}
\begin{solution}
La firma fusionada maximiza $\pi_{2\&3} = (a - q_1 - q_2 - q_3 - c)(q_2 + q_3)$. Las CPO respecto de $q_2$ y $q_3$ son $a - c - q_1 - 2q_2 - 2q_3 = 0$, de modo que solo importa $q_{23} = q_2 + q_3$: $q_{23} = \frac{(a - c) - q_1}{2}$, mientras que $q_1 = \frac{(a-c) - q_{23}}{2}$. Resolviendo,
\[
  q_1^* = \frac{a - c}{3}, \qquad q_2^* + q_3^* = \frac{a - c}{3} \;\;\Longrightarrow\;\; q_2^* = q_3^* = \frac{a - c}{6}.
\]
(Es distinto al caso anterior: la fusión equivale a que una empresa desaparece del mercado.) El precio es $P = a - \frac{2(a-c)}{3} = \frac{a + 2c}{3}$.

\end{solution}

\part Con los siguientes valores compare la cantidad producida por cada firma y el precio de equilibrio antes y después de la fusión: $a = 100$, $c = 20$.
\needspace{5\baselineskip}
\begin{solution}
\emph{Antes de la fusión:} $q_1 = q_2 = q_3 = \frac{100 - 20}{4} = 20$ y $P = 100 - 60 = 40$.

\emph{Después de la fusión:} $q_1 = \frac{80}{3} \approx 26{,}7$, $q_2 = q_3 = \frac{80}{6} \approx 13{,}3$, $Q \approx 53{,}3$ y $P = \frac{100 + 40}{3} \approx 46{,}7$.

El precio sube ($40 \to 46{,}7$) y la cantidad total cae ($60 \to 53{,}3$). Además, la entidad fusionada obtiene $\pi_{2\&3} = (46{,}7 - 20)\cdot 26{,}7 \approx 711 < 800 = 2\cdot 400$ (cada firma ganaba $20\cdot 20 = 400$ antes), de modo que la fusión \emph{no} es rentable para las firmas fusionadas, mientras que beneficia a la firma 1 (de $400$ a $\approx 711$).

\end{solution}

\part Dé una intuición económica de los cambios que encontró en el ítem (c).
\needspace{5\baselineskip}
\begin{solution}
Al fusionarse, las firmas 2 y 3 internalizan el efecto de la producción de la otra sobre el precio y deciden producir menos que en el caso donde cada una maximiza por separado: actúan de manera más cooperativa. La firma 1 responde de manera óptima a una menor producción del competidor produciendo más que en el caso sin fusión. Este resultado ilustra la crítica de Salant, Switzer y Reynolds (1983): en el modelo estándar de Cournot, una fusión entre firmas simétricas no es rentable a menos que involucre $80\%$ o más de las empresas de la industria; por ello, estos modelos simples no logran reflejar la naturaleza básica de las fusiones (se requieren modelos de diferenciación de producto o de capacidad).

\end{solution}
\end{parts}

%-------------------------------------------------------------------- Caso Meta
\needspace{8\baselineskip}
\tema{Caso Meta: adquisiciones de Instagram y WhatsApp}
\question La Federal Trade Commission (FTC) demandó a Meta (ex-Facebook) denunciando que la compañía mantuvo una posición monopólica ilegal en el mercado de \emph{personal social networking services} a través de las adquisiciones de Instagram y WhatsApp. En 2008, Zuckerberg escribió: ``\emph{It is better to buy than compete}''. Meta es una de las empresas con mayor participación de mercado en redes sociales (78\% de los usuarios mensualmente activos, 77\% de los usuarios diariamente activos y 85\% del tiempo en pantalla). Discuta (i) los antecedentes de las adquisiciones, (ii) el daño a la competencia y (iii) la definición del mercado relevante.
\needspace{5\baselineskip}
\begin{solution}
\textbf{(i) Antecedentes.} En 2012, Facebook pudo haber visto a Instagram como una amenaza futura, por lo que comprar y escalar a los competidores era ideal en ese momento. Por otro lado, la compra de Instagram puede verse como una barrera de entrada a potenciales entrantes, dañando la competencia. Factores similares se encuentran en el caso de WhatsApp.

\textbf{(ii) Daño a la competencia.} El daño a la competencia no solo se da en aumento de precios, sino también en reducción de la calidad del producto o servicio, menos inversión en innovación y otras malas prácticas. Meta podría estar extrayendo rentas monopólicas: su WACC es alrededor del 10\% y su retorno anual es del 36--41\%.

\textbf{(iii) Mercado relevante.} A grandes rasgos, la FTC argumenta que el mercado relevante incluye las aplicaciones dedicadas a conectarse con amigos y familia (WhatsApp, Instagram, Facebook). Zuckerberg se defiende con una definición de mercado más amplia, y propone que Facebook compite con YouTube o TikTok. De la definición de mercado depende la participación de mercado y, con ella, el HHI y el poder de mercado imputado.

\end{solution}

\end{unidad}

\unidadtitulo{Discriminación de precios}
\begin{unidad}
%-------------------------------------------------------------------- Comente: 1er vs 3er grado
\needspace{8\baselineskip}
\tema{Comente: discriminación de primer vs tercer grado}
\question \textbf{Comente.} ¿Cuál es la diferencia entre discriminación de primer y tercer grado?
\needspace{5\baselineskip}
\begin{solution}
En la discriminación de primer grado la firma puede fijar el precio que maximice su excedente, dejando nada del excedente al consumidor (cobra a cada consumidor su precio de reserva). En la discriminación de tercer grado la firma explota las características observables del comprador para cobrar precios diferenciados según el grupo.

\end{solution}

%-------------------------------------------------------------------- Comente: mercados
\needspace{8\baselineskip}
\tema{Comente: mercados con discriminación de primer y tercer grado}
\question \textbf{Comente.} Caracterice los mercados en los que suele haber discriminación de primer grado y de tercer grado.
\needspace{5\baselineskip}
\begin{solution}
Los mercados en donde suele haber discriminación de primer grado son:
\begin{itemize}
  \item Mercados donde el número de compradores (\emph{customers}) es relativamente pequeño y el vendedor posee considerable información sobre los compradores.
  \item Industrias como el concreto fresco, aviones de pasajeros, software especializado para empresas, etc.
  \item A pesar de que existe un precio de lista, cada cliente recibe un descuento que se negocia: el precio final depende de la disposición a pagar del cliente y de su poder de negociación.
\end{itemize}
Por otro lado, los mercados donde suele haber discriminación de tercer grado son mercados más grandes, donde los costos de información son mayores; por ejemplo, los pasajes del transporte público.

\end{solution}

%-------------------------------------------------------------------- Comente: captar todo el excedente
\needspace{8\baselineskip}
\tema{Comente: captar todo el excedente del consumidor}
\question \textbf{Comente.} Es imposible que una firma monopólica logre captar todo el excedente del consumidor, puesto que los precios y la estrategia que impongan dependerán también de factores como la elasticidad de la demanda.
\needspace{5\baselineskip}
\begin{solution}
Es verdad que en la práctica es difícil que se logre captar todo el excedente, pero si el monopolio tiene información perfecta sobre la disposición a pagar de los consumidores podría imponer una estrategia de discriminación de precios de primer grado, donde cada individuo paga el máximo que está dispuesto a pagar por el producto en cuestión, generando que el excedente sea nulo para los consumidores y que la firma se lleve todo el excedente del mercado.

\end{solution}

%-------------------------------------------------------------------- Comente: discriminar o no
\needspace{8\baselineskip}
\tema{Comente: ¿discriminar o no discriminar?}
\question \textbf{Comente.} ¿Discriminar o no discriminar? Discuta los principios básicos para responder esta pregunta.
\needspace{5\baselineskip}
\begin{solution}
Cualquier baja en la producción suele llevar a que el bienestar total baje. Una discriminación que aumenta la demanda del producto o servicio suele aumentar el bienestar (no hay consumidor más triste que el que no consume). Los principios básicos son:
\begin{itemize}
  \item Si la producción total cae con la discriminación, el bienestar total disminuye.
  \item Si un monopolista que no discrimina cierra un mercado, es mejor discriminar.
\end{itemize}

\end{solution}

%-------------------------------------------------------------------- Museo
\needspace{8\baselineskip}
\tema{Discriminación de tercer grado: el museo}
\question En un pueblo del sur de Chile hay un único museo que recibe a visitantes nacionales y extranjeros, quienes presentan una mayor valoración del museo. El costo marginal de producción del museo es 1 por cada visitante y no hay costos fijos. Las demandas de extranjeros y nacionales son
\begin{align*}
  Q_E &= 10 - P_E \\
  Q_N &= 8 - P_N
\end{align*}
\begin{parts}
\part Suponga que el museo monopólico fija un precio uniforme. ¿Qué precio cobra y qué beneficio obtiene? ¿Qué condición debe cumplirse para que se venda a ambos grupos?
\needspace{5\baselineskip}
\begin{solution}
La firma maximiza beneficios sobre la demanda agregada, $Q_T = 18 - 2P$:
\[
\begin{gathered}
  \max_P\ \pi = (P - c)(18 - 2P) \\
  \;\Rightarrow\; 18 - 4P + 2c = 0 \\
  \;\Rightarrow\; P^U = 5, \\
  \qquad \pi^U = (5-1)(18 - 2\cdot 5) = 32 .
\end{gathered}
\]
Para que se venda a ambos grupos debe cumplirse la restricción de participación del grupo de menor valoración: los nacionales no demandarían ninguna entrada a precios mayores a 8, por lo que debe cumplirse $P \leq 8$ (la demanda agregada $18-2P$ solo es válida en ese tramo). Como $P^U = 5 < 8$, se vende a ambos.

\end{solution}

\part Suponga ahora que el monopolio puede discriminar en tercer grado. Calcule los precios de cada mercado y los beneficios del museo.
\needspace{5\baselineskip}
\begin{solution}
Para los extranjeros el museo resuelve
\[
\begin{gathered}
  \max_{P_E}\ (P_E - c)(10 - P_E) \\
  \;\Rightarrow\; 10 - 2P_E + c = 0 \\
  \;\Rightarrow\; P_E = \frac{10 + c}{2} = 5{,}5, \qquad \pi_E = (5{,}5-1)(10-5{,}5) = 20{,}25 .
\end{gathered}
\]
Y de la misma manera, para los nacionales,
\[
  \max_{P_N}\ (P_N - c)(8 - P_N)
  \;\Rightarrow\; P_N = \frac{8 + c}{2} = 4{,}5, \qquad \pi_N = (4{,}5-1)(8-4{,}5) = 12{,}25 .
\]
Los beneficios totales son $\pi^D = \pi_E + \pi_N = 32{,}5$.

\end{solution}

\part Un compañero le menciona que el museo debería enfocarse solo en el grupo de alta valoración. Demuéstrele que al museo no le conviene enfocarse solo en una parte del mercado, cuando tiene la opción de discriminar precios.
\needspace{5\baselineskip}
\begin{solution}
Si el museo se enfocara solo en los extranjeros, fijaría un precio en que solo ellos demanden una cantidad positiva ($P = 8$). Los beneficios serían entonces
\[
  \pi_E = (8-1)(10-8) = 14,
\]
lo cual es bastante menor que al no discriminar ($32$) o al discriminar en tercer grado ($32{,}5$).

\end{solution}

\part Comparando los beneficios del monopolio con la estrategia de precio uniforme y de discriminación, ¿cuál le conviene más? Además, calcule los excedentes de los consumidores extranjeros y nacionales. ¿Cuál grupo se beneficia de la discriminación de tercer grado y cuál se perjudica?
\needspace{5\baselineskip}
\begin{solution}
Como $\pi^D = 32{,}5 > \pi^U = 32$, al museo le conviene discriminar en tercer grado (aunque la ganancia es marginal).

Los excedentes de los consumidores con precio uniforme (U) y con precio diferenciado (D) son:
\begin{align*}
  EC_E^U &= \frac{5(10-5)}{2} = 12{,}5, & EC_E^D &= \frac{4{,}5(10-5{,}5)}{2} = 10{,}125, \\
  EC_N^U &= \frac{3(8-5)}{2} = 4{,}5, & EC_N^D &= \frac{3{,}5(8-4{,}5)}{2} = 6{,}125 .
\end{align*}
El grupo de mayor valoración (extranjeros) sale perjudicado, mientras que el de menor valoración (nacionales) está mejor siendo discriminado. En cuanto al excedente total, los consumidores están peor con discriminación ($17 \to 16{,}25$).

\end{solution}

\part ¿Qué pasa con el beneficio social? Calcule también el producto total ofrecido en cada estrategia y comente su relación con los cambios en bienestar o ineficiencias que puedan estar ocurriendo. ¿Qué otros factores explican los cambios en bienestar al pasar de una estrategia a otra?
\needspace{5\baselineskip}
\begin{solution}
El beneficio social para cada estrategia corresponde a
\[
  ES^U = \pi^U + EC^U = 32 + 17 = 49, \qquad ES^D = \pi^D + EC^D = 32{,}5 + 16{,}25 = 48{,}75 .
\]
Se concluye que en este caso el bienestar social no aumenta con la discriminación de precios. La producción total de cada estrategia es
\[
  Q^U = 18 - 2\cdot 5 = 8, \qquad Q^D = Q_E^D + Q_N^D = (10 - 5{,}5) + (8 - 4{,}5) = 4{,}5 + 3{,}5 = 8 .
\]
Como el producto total no cambia, no hay un efecto de expansión de la producción que compense; el bienestar cae por una \emph{mala asignación} del bien entre los consumidores: con precios distintos hay consumidores con mayor valoración que consumen menos y consumidores con menor valoración que consumen más (se pierde eficiencia en la asignación).

\end{solution}
\end{parts}

%-------------------------------------------------------------------- Comente: condiciones discriminación
\needspace{8\baselineskip}
\tema{Comente: condiciones para discriminar precios}
\question \textbf{Comente.} La discriminación de precios solo es posible si es que se cumplen las siguientes condiciones: (i) es posible establecer que existen grupos con diferentes demandas o disposiciones a pagar en el mercado y (ii) es posible identificar explícitamente a los individuos que pertenecen a cada grupo.
\needspace{5\baselineskip}
\begin{solution}
Falso. La discriminación de precios de segundo grado no requiere que se cumpla la condición (ii). Se pueden diseñar ofertas o planes tales que cada consumidor se autoseleccione en el servicio que va dirigido a su propio grupo.

\end{solution}

%-------------------------------------------------------------------- Tarifas en dos partes
\needspace{8\baselineskip}
\tema{Tarifas en dos partes (dos tipos de clientes)}
\question Una empresa ofrece dos planes de servicio para sus clientes, los que consisten en un monto fijo $F$ y un precio $P$ por unidad. Esta empresa se enfrenta a dos mercados distintos con las siguientes funciones de demanda:
\begin{align*}
  Q_1 &= 100 - P \\
  Q_2 &= 80 - P
\end{align*}
Supondremos que ofrecer el servicio a cualquier mercado tiene un costo de cero. Además, la firma tiene la opción de fijar una tarifa fija $F$ además de una tarifa por unidad $P$.
\begin{parts}
\part Si usted puede separar a los clientes de cada mercado, ¿cuál sería la tarifa $F$ y cuál la tarifa por unidad $P$ para cada tipo de cliente? ¿Cuánto sería el beneficio de la empresa?
\needspace{5\baselineskip}
\begin{solution}
Si puedo separar a los clientes de cada mercado significa que sé quién es quién (y de esa forma puedo evitar el arbitraje). Lo más lógico es extraer a cada individuo su máxima disposición a pagar, que es igual al excedente máximo de cada individuo. El $F$ óptimo busca apropiarse del excedente del consumidor que queda al cobrar una tarifa por unidad $P$:
\[
  F_1 = \frac{(100 - P_1)^2}{2}, \qquad \pi_1 = F_1 + P_1(100 - P_1) = \frac{(100-P_1)^2}{2} + (100 - P_1)P_1 .
\]
Esta expresión se maximiza con $P_1 = 0$ (cobrar al costo marginal, para maximizar el excedente total que luego se extrae), por lo tanto $F_1 = \frac{100^2}{2} = 5000$. Con la misma lógica, $P_2 = 0$ y
\[
  F_2 = \frac{80^2}{2} = 3200 .
\]
Los beneficios totales de la empresa serán $\pi = 5000 + 3200 = 8200$.

\end{solution}

\part Suponga que no puede separar a los clientes y la tarifa fija $F$ está restringida a cero. ¿Cuál sería la tarifa $P$ que cobraría? ¿Cuánto sería el beneficio de la empresa?
\needspace{5\baselineskip}
\begin{solution}
Si no se puede separar a los clientes, se cobrará un precio uniforme en ambos mercados. Como $F=0$, solo se debe decidir $P$. Sumando las demandas horizontalmente, la demanda total es
\[
  D(P) =
  \begin{cases}
    100 - P & \text{si } P > 80 \\
    180 - 2P & \text{si } P \leq 80 .
  \end{cases}
\]
El tramo donde maximizará la firma es el segundo (si se maximiza en el primero se obtiene $P = 50$, pero la condición del tramo es $P>80$):
\[
  \max_P\ \pi = (180 - 2P)P \;\Rightarrow\; P = 45, \quad Q = 90, \quad \pi = 4050 .
\]
Los beneficios son menores que en la parte (a).

\end{solution}

\part Suponga ahora que cobra una tarifa en dos partes, es decir, $F$ y $P$, pero sin diferenciar al tipo de cliente, sino que para todos lo mismo. ¿Cuál sería la tarifa fija $F$ y cuál la tarifa por unidad $P$? ¿Cuánto sería el beneficio de la empresa?
\needspace{5\baselineskip}
\begin{solution}
En este caso no podemos cobrar una tarifa fija de $5000$ como en el primer caso, pues los clientes de demanda baja no consumirían. Por lo tanto, cobramos la tarifa fija más alta posible que sigue siendo aceptada por ambos, que es precisamente el excedente de los consumidores de demanda baja:
\[
  F = \frac{(80 - P)^2}{2}.
\]
La empresa maximiza
\[
\begin{gathered}
  \max_P\ \pi = 2F + P(100 - P) + P(80 - P) = (80-P)^2 + P(180 - 2P) \\
  \;\Rightarrow\; \\
  \frac{\partial\pi}{\partial P} = -2(80-P) + 180 - 4P = 20 - 2P = 0 \\
  \;\Rightarrow\; P = 10 .
\end{gathered}
\]
La tarifa fija correspondiente es $F = \frac{(80-10)^2}{2} = 2450$, y los beneficios son $\pi = 2\cdot 2450 + 10\cdot(90 + 70) = 6500$.

\end{solution}

\part ¿De qué forma la tarifa en dos partes le permite discriminar? ¿Qué restricciones enfrenta al maximizar beneficios esta empresa?
\needspace{5\baselineskip}
\begin{solution}
El beneficio con tarifa en dos partes ($6500$) es mayor que cobrar solo un precio unitario ($4050$). Esto se produce porque a través de la tarifa en dos partes se puede inducir a los distintos tipos de individuos presentes en el mercado (demanda alta vs.\ demanda baja) a autoseleccionarse de acuerdo a los menús de precios y tarifas ofrecidos. Esta discriminación será posible si los individuos presentes en el mercado tienen demandas distintas. Con respecto a las restricciones que enfrenta la firma a la hora de maximizar, existen dos tipos:
\begin{enumerate}
  \item[(a)] \emph{Restricción de participación:} la tarifa y el precio deben ser lo suficientemente atractivos para que los individuos a los cuales se quiere vender, compren.
  \item[(b)] \emph{Restricción de compatibilidad de incentivos:} puede que los individuos de demanda alta tengan el incentivo a comprar el paquete de menor costo. Por lo tanto, se debe asegurar que los clientes no se hagan pasar unos por otros.
\end{enumerate}

\end{solution}
\end{parts}

%-------------------------------------------------------------------- Comente: 2º grado tarifa en dos partes
\needspace{8\baselineskip}
\tema{Comente: discriminación de segundo grado}
\question \textbf{Comente.} En una discriminación de segundo grado el monopolista no observa ninguna característica del consumidor que le permita aplicar una tarifa en dos partes similar a la discriminación perfecta.
\needspace{5\baselineskip}
\begin{solution}
Verdadero. Bajo discriminación de segundo grado el monopolista no observa directamente alguna característica que le permita separar a los consumidores, pero puede establecer una tarifa en dos partes con un cargo variable superior al costo marginal y un cargo fijo equivalente al excedente que obtendría el consumidor con menor disponibilidad a pagar. Esta tarifa difiere de la aplicada bajo discriminación perfecta en que allí el cargo variable es igual al costo marginal y el cargo fijo puede ser diferenciado.

\end{solution}

%-------------------------------------------------------------------- Comente: esquemas de tarifación
\needspace{8\baselineskip}
\tema{Comente: tipos de discriminación (cine, buffet, transporte)}
\question \textbf{Comente.} Es interesante observar que los esquemas de tarificación difieren en distintas actividades. Describa el tipo de discriminación de precios y por qué se usa (o no se usa) en los siguientes casos:
\begin{parts}
\part Un cine durante un día normal.
\needspace{5\baselineskip}
\begin{solution}
Al haber descuentos para estudiantes y tercera edad, esto constituye una discriminación de tercer grado. Con las salas \emph{premium} de algunos cines se ofrecen paquetes distintos, lo que es discriminación de segundo grado para que los clientes se autoseleccionen según su disposición a pagar.

\end{solution}

\part Restaurantes con buffet (se puede comer cuanto se desea).
\needspace{5\baselineskip}
\begin{solution}
Discriminación de segundo grado. Los clientes se autoseleccionan; en general se cobra una tarifa fija porque no se puede controlar el consumo.

\end{solution}

\part Transporte público con pasajes especiales para estudiantes.
\needspace{5\baselineskip}
\begin{solution}
Discriminación de tercer grado, pues los pases escolares de los estudiantes son observables.

\end{solution}
\end{parts}

%-------------------------------------------------------------------- Comente: plan de telefonía
\needspace{8\baselineskip}
\tema{Comente: planes de telefonía móvil}
\question \textbf{Comente.} Suponga que usted es una/un feliz ayudante de la FEN-UChile que se dispone a recibir su pago de ayudantía en los próximos días. Usted ha decidido destinar este dinero a renovar su celular, el que actualmente tiene un plan de 200 minutos e internet ilimitado. Con esta idea se dirige a la compañía de telefonía móvil y cotiza un nuevo equipo. La vendedora le hace el siguiente comentario: «El nuevo equipo cuesta \$50.000. Sin embargo, su plan actual de \$25.990 está descontinuado. Dado esto, lo homologaremos al más cercano de los nuevos sin perjudicar su pago. Este corresponde al plan de \$23.990, que consta de 150 minutos y 1 Gb de internet. Adicionalmente, le ofrecemos el plan de \$29.990 con 300 minutos y 2,5 Gb de internet». ¿Qué estrategia está utilizando la compañía móvil? ¿De qué forma funciona?
\needspace{5\baselineskip}
\begin{solution}
La compañía móvil está utilizando una estrategia de discriminación de segundo grado. La empresa no puede exigir que el consumidor esté en cierto plan, pero ofrece dos planes distintos donde el cliente tiene que autoseleccionarse. En particular, lo que hace la firma es perjudicar la opción que corresponde a los consumidores de menor valoración, disminuyendo el número de minutos y ofreciendo una cantidad limitada de internet móvil. Dado lo anterior, el excedente que puede obtener el individuo de mayor valoración escogiendo este plan disminuye. Esto puede empujar a que este decida optar por el plan más costoso.

\end{solution}

%-------------------------------------------------------------------- Comente: 3er grado bienestar
\needspace{8\baselineskip}
\tema{Comente: discriminación de tercer grado y bienestar}
\question \textbf{Comente.} Un economista afirma: «La discriminación de precios de tercer grado resulta en que algunos consumidores pagan un precio mayor que lo que hubieran pagado si el monopolista cobrara un precio uniforme. Por lo mismo, la discriminación de precios disminuye el bienestar y debería estar prohibida.» ¿Usted está de acuerdo con dicha afirmación? Fundamente su respuesta.
\needspace{5\baselineskip}
\begin{solution}
No estoy de acuerdo. La discriminación de precios puede aumentar o reducir el bienestar. Con discriminación de precios, los consumidores con mayor disposición a pagar (menor elasticidad de demanda) pagan precios mayores, pero los consumidores con menor disposición a pagar (mayor elasticidad de demanda) pagan precios menores. Al final del día, el bienestar puede aumentar o disminuir: depende de si la producción total aumenta (por ejemplo, si se abren mercados que no se servían) o no.

\end{solution}

%-------------------------------------------------------------------- Comente: aerolíneas
\needspace{8\baselineskip}
\tema{Comente: discriminación de precios en aerolíneas}
\question \textbf{Comente.} Debería prohibirse la discriminación de precios por parte de las aerolíneas, ya que esto les permite abusar de su poder de mercado.
\needspace{5\baselineskip}
\begin{solution}
Falso. En el extremo, incluso un monopolio discriminador en primer grado produce la cantidad socialmente óptima, por lo que en la medida que mejor pueda discriminar, disminuirá la pérdida social de eficiencia. Además, en la industria aérea existen costos fijos importantes, pero los costos variables y marginales son muy pequeños (una vez que el vuelo se programó y se pagó, el costo de subir pasajeros es prácticamente cero). En este caso, una efectiva discriminación de precios hace que las aerolíneas puedan mejorar sus factores de ocupación, disminuyendo sus costos medios y por ende pudiendo cobrar menores tarifas.

\end{solution}

%-------------------------------------------------------------------- AeroStrong
\needspace{8\baselineskip}
\tema{Discriminación de precios en aerolíneas: AeroStrong y AeroWeak}
\question La compañía aérea AeroStrong tiene dos tipos de pasajeros según la motivación del viaje: turismo y negocios. Las demandas de cada grupo son
\begin{align*}
  Q_t &= 20 - 4P_t \\
  Q_n &= 22 - P_n
\end{align*}
Los costos de movilizar un pasajero son independientes del motivo del viaje y es \$2.
\begin{parts}
\part Si la aerolínea puede evitar el arbitraje entre viajeros, ¿cuánto le cobraría por el pasaje a cada tipo de pasajero? En términos prácticos, ¿cómo la aerolínea podría evitar el arbitraje?
\needspace{5\baselineskip}
\begin{solution}
Si la empresa puede evitar el arbitraje entre los viajeros, cobrará un precio que maximice sus beneficios en cada mercado. Para los turistas:
\[
  \max_{P_t}\ (P_t - 2)(20 - 4P_t) \;\Rightarrow\; 28 - 8P_t = 0 \;\Rightarrow\; P_t = 3{,}5 .
\]
Para los pasajeros de negocios:
\[
  \max_{P_n}\ (P_n - 2)(22 - P_n) \;\Rightarrow\; 24 - 2P_n = 0 \;\Rightarrow\; P_n = 12 .
\]
Por ejemplo, si la aerolínea vende pasajes empaquetados ida y vuelta, podría identificar el propósito del viaje según si el pasajero permanece el fin de semana en el destino o no.

\end{solution}

\part Si la aerolínea se ve impedida de discriminar entre pasajeros, ¿cuál sería el precio que cobraría por los pasajes?
\needspace{5\baselineskip}
\begin{solution}
El monopolio fija un precio uniforme sobre la demanda agregada:
\[
  Q =
  \begin{cases}
    42 - 5P & \text{si } P \leq 5 \\
    22 - P & \text{si } P > 5
  \end{cases}
\]
El tramo donde maximizará la firma es el segundo: si maximizara en el primero el precio sería $P = 5{,}2$, pero la condición del tramo es $P \leq 5$. Entonces resuelve $\max_P\ (P-2)(22-P)$, cuya CPO es $24 - 2P = 0$, es decir, $P = 12$ (el grupo de turismo no consume).

\end{solution}

\part La aerolínea de bajo costo AeroWeak ingresa a competir en el segmento de pasajeros de turismo solamente, puesto que los pasajeros de negocios no la consideran una alternativa para sus viajes. El costo por pasajero de AeroWeak es \$1. ¿Qué es preferible para la entrante: que AeroStrong discrimine precios, que fije un precio único, o es indiferente dado que AeroWeak no compite en el segmento de negocios? Fundamente su respuesta.
\needspace{5\baselineskip}
\begin{solution}
Si AeroStrong fija un precio único ($12$), solamente abastecerá al grupo de negocios, por lo que AeroWeak podrá actuar como un monopolio en el grupo de turistas:
\[
\begin{gathered}
  \max_{P_t}\ (P_t - 1)(20 - 4P_t) \;\Rightarrow\; 24 - 8P_t = 0 \;\Rightarrow\; P_t = 3, \quad Q_t = 8, \\ \pi = 2\cdot 8 = 16 .
\end{gathered}
\]
Si AeroStrong discrimina precios, ambas empresas competirán en el segmento de turistas. Suponiendo que los turistas son indiferentes entre ambas aerolíneas y eligen la más barata, compiten à la Bertrand, y como AeroWeak es más eficiente fija $P_t = 2 - \epsilon$, con beneficios $\approx (2-1)\cdot(20-8) = 12 < 16$.

Por lo tanto, AeroWeak prefiere que AeroStrong fije un precio único, ya que le permite fijar un precio monopólico en el mercado de turistas.

\end{solution}
\end{parts}

%-------------------------------------------------------------------- Tarifas Clavistel
\needspace{8\baselineskip}
\tema{Tarifas en dos partes: Clavistel (variante numérica)}
\question La empresa de telefonía móvil Clavistel ofrece dos planes de servicio para sus clientes, los que consisten en un monto fijo $F$ y un precio $P$ por unidad del bien. El servicio a cualquier mercado tiene un costo de cero. Además, esta empresa se enfrenta a dos mercados distintos con las siguientes funciones de demanda: $q_1 = 40 - p$ y $q_2 = 20 - p$.
\begin{parts}
\part Si la empresa puede separar a los clientes de cada mercado, ¿cuál sería la tarifa $F$ y cuál la tarifa por unidad $P$ para cada tipo de cliente? Obtenga los beneficios de la empresa con esas tarifas.
\needspace{5\baselineskip}
\begin{solution}
Si la firma puede separar a sus clientes, le extraerá su máxima disposición a pagar, igual al excedente máximo de cada individuo. Para el cliente 1: $\frac{40\cdot 40}{2} = 800$. Con $F_1 = \frac{(40 - p_1)^2}{2}$, el beneficio es $F_1 + (40 - p_1)p_1$, que se maximiza en $p_1 = 0$ y $F_1 = 800$. Para el cliente 2, $p_2 = 0$ y $F_2 = \frac{20^2}{2} = 200$. Los beneficios totales son $800 + 200 = 1000$.

\end{solution}

\part Suponga ahora que la empresa no puede separar a los clientes y la tarifa fija $F$ está restringida a cero. ¿Cuál sería la tarifa $P$ que cobraría? ¿Cuánto sería el beneficio de la empresa?
\needspace{5\baselineskip}
\begin{solution}
Se cobra un precio uniforme en ambos mercados. La demanda total es $D(p) = 40 - p$ si $p > 20$ y $D(p) = 60 - 2p$ si $p \leq 20$. El tramo donde maximizará es el segundo (en el primero se tendría $p = 20$, justo en el borde, con beneficios menores): $\max_p\ (60 - 2p)p$, por lo tanto $P = 15$, $Q = 30$ y los beneficios son $450 < 1000$.

\end{solution}

\part Suponga ahora que cobra una tarifa en dos partes, es decir, $F$ y $P$, pero sin diferenciar al tipo de cliente. ¿Cuál sería la tarifa $F$ y $P$ que cobraría? ¿Cuánto sería el beneficio de la empresa?
\needspace{5\baselineskip}
\begin{solution}
No podemos cobrar una tarifa fija de 800, pues los clientes de demanda baja no consumirían; por lo tanto, cobramos la tarifa fija más alta posible, que es precisamente el excedente de los consumidores de demanda baja, $F = \frac{(20-p)^2}{2}$. La empresa maximiza
\[
\begin{gathered}
  \max_p\ \pi = 2\cdot\frac{(20-p)^2}{2} + (40 - p)p + (20 - p)p = (20-p)^2 + p(60 - 2p) \\
  \;\Rightarrow\; 20 - 2p = 0 \;\Rightarrow\; p = 10 .
\end{gathered}
\]
Entonces $F = \frac{10^2}{2} = 50$ y los beneficios son $\pi = 2\cdot 50 + 10\cdot(30 + 10) = 500$.

\end{solution}

\part ¿Qué restricciones enfrenta la empresa al maximizar beneficios discriminando con tarifa en dos partes?
\needspace{5\baselineskip}
\begin{solution}
La empresa enfrenta restricciones de dos tipos:
\begin{enumerate}
  \item[(a)] \emph{Restricción de participación:} debe hacer la tarifa y el precio lo suficientemente atractivos para que los individuos a los cuales se quiere vender, compren.
  \item[(b)] \emph{Restricción de autoselección (compatibilidad de incentivos):} no se debe permitir que los consumidores de un grupo usen los planes de otro.
\end{enumerate}

\end{solution}
\end{parts}

%-------------------------------------------------------------------- Tipos de discriminación
\needspace{8\baselineskip}
\tema{Los tres tipos de discriminación de precios}
\question Describa y explique los tres tipos de discriminación que puede hacer un monopolista según distintos casos.
\needspace{5\baselineskip}
\begin{solution}
\textbf{Segundo grado.} Discriminación por medio de la autoselección de los individuos según su comportamiento; por ejemplo, tarifas:
\begin{itemize}
  \item Tarifa lineal: $T(q) = pq$.
  \item Tarifa de dos partes: $T(q) = A + pq$, un cargo fijo más un cobro por unidad consumida.
  \item Tarifa no lineal: $T(q) = A_j + p_jq$ si $q_{j-1} < q < q_j$, con distintos cargos según tramos de consumo. Es común en bares, gimnasios, festivales de música, entre otros.
\end{itemize}
\textbf{Primer grado.} En el caso anterior el monopolista no conoce completamente la demanda por los individuos; si tuviera completo conocimiento sobre sus disposiciones a pagar, podría fijar tarifas que extrajeran todo su excedente. Asumiendo que no hay arbitraje, el monopolista podría tomar el precio de competencia ($p = c$) pero cobrar una suma fija según el individuo, que represente el excedente que obtendría:
\[
  T(q) =
  \begin{cases}
    A_i + cq & \text{si } q > 0 \\
    0 & \text{si no.}
  \end{cases}
\]
Se maximiza el bienestar social como en competencia perfecta, pero el monopolista se queda con todo el excedente. Esto suele estar presente en servicios más personalizados (\emph{customer markets}).

\textbf{Tercer grado.} A partir de información observable se pueden cobrar precios diferenciados. La información que tiene el monopolista es respecto a la demanda por grupos, pero no por individuo. El monopolista enfrenta un problema similar al del monopolista multiproducto:
\[
  \max_{p_1,\dots,p_n}\ \sum_{i=1}^{n}p_i q_i(p_i) - c\sum_{i=1}^{n}q_i(p_i).
\]
El ingreso marginal debe ser igual en cada grupo (con dos grupos, $i = 1,2$): $IMg_1 = c = IMg_2$. Desde el punto de vista de las elasticidades, $IMg_i = p_i\bigl(1 - \frac{1}{|\varepsilon_i|}\bigr) = c$, de modo que
\[
  p_i = \frac{c}{1 - 1/|\varepsilon_i|}.
\]
Mientras mayor (menor) sea la elasticidad en determinado grupo, se le cobra menos (más).

\end{solution}

%-------------------------------------------------------------------- Principios bienestar
\needspace{8\baselineskip}
\tema{Principios para evaluar el bienestar con discriminación}
\question ¿Cuáles son los principios básicos para evaluar cambios en el bienestar dada la discriminación?
\needspace{5\baselineskip}
\begin{solution}
Si la cantidad de equilibrio cae, entonces el bienestar total disminuye. Dadas dos demandas, puede que en equilibrio monopólico haya una ineficiencia por no servir cierto mercado; al discriminar y dejar de servir a más personas, se acentúa la ineficiencia. En caso de que la discriminación abra un mercado, por la misma lógica anterior debería aumentar el bienestar social.

\end{solution}

%-------------------------------------------------------------------- Demostración bienestar
\needspace{8\baselineskip}
\tema{Cotas al cambio en bienestar por discriminación de tercer grado}
\question Demuestre que el cambio en el bienestar al pasar de un monopolio no discriminante a un monopolio discriminante satisface
\[
  \sum_{i}(p_i - c)\Delta q_i \;\leq\; \Delta W \;\leq\; \sum_{i}(\bar p - c)\Delta q_i ,
\]
donde $\bar p$ es el precio uniforme, $p_i$ el precio discriminado en el mercado $i$, y $\Delta q_i = q_i(p_i) - q_i(\bar p)$ es el cambio en las ventas del mercado $i$ respecto a la situación sin discriminación. Interprete la desigualdad.
\needspace{5\baselineskip}
\begin{solution}
Primero establecemos la relación entre el precio, la demanda y el bienestar del consumidor. Con un precio $p$ se obtiene un excedente del consumidor $S_i(p) = \int_p^{p^*}q_i(s)\,ds$, el área bajo la curva de demanda desde $p$ hasta el precio donde la demanda es cero. Por lo tanto $S_i'(p) = -q_i(p)$ y $S_i''(p) = -q_i'(p) > 0$: el excedente del consumidor es convexo en el precio, lo que implica
\[
  S_i(p_i) - S_i(\bar p) \geq S_i'(\bar p)(p_i - \bar p) = -q_i(\bar p)(p_i - \bar p)
\]
\[
  \text{y} \qquad S_i(p_i) - S_i(\bar p) \leq S_i'(p_i)(p_i - \bar p) = -q_i(p_i)(p_i - \bar p).
\]
El cambio en el bienestar es
\[
  \Delta W = \sum_i\bigl[S_i(p_i) - S_i(\bar p)\bigr] + \sum_i\bigl[(p_i - c)q_i(p_i) - (\bar p - c)q_i(\bar p)\bigr].
\]
\emph{Cota inferior.} Usando la primera desigualdad, escribiendo $\bar q_i = q_i(\bar p)$ y $q_i = q_i(p_i)$:
\[
  \Delta W \geq \sum_i\bigl[-\bar q_i(p_i - \bar p) + (p_i - c)q_i - (\bar p - c)\bar q_i\bigr] = \sum_i\bigl[(p_i - c)q_i - (p_i - c)\bar q_i\bigr] = \sum_i(p_i - c)\Delta q_i .
\]
\emph{Cota superior.} Usando la segunda desigualdad:
\[
  \Delta W \leq \sum_i\bigl[-q_i(p_i - \bar p) + (p_i - c)q_i - (\bar p - c)\bar q_i\bigr] = \sum_i\bigl[(\bar p - c)q_i - (\bar p - c)\bar q_i\bigr] = \sum_i(\bar p - c)\Delta q_i .
\]
\textbf{Interpretación.} Se llega a los mismos principios discutidos antes: el cambio en bienestar está acotado por el cambio en las ventas valorado al margen de precio sobre costo. Si la discriminación lleva a un aumento de ventas por abrir un mercado ($\Delta q_i > 0$ en los mercados donde $p_i > c$), habrá un aumento en el bienestar; una condición necesaria para que el bienestar aumente es que la producción total aumente ($\sum\Delta q_i > 0$).

\end{solution}

%-------------------------------------------------------------------- Editorial
\needspace{8\baselineskip}
\tema{Discriminación intertemporal: tapa dura y tapa blanda}
\question Suponga una editorial con poder monopólico que está por lanzar un nuevo libro al mercado y que enfrenta dos tipos de compradores (hay una misma proporción de cada uno de los dos tipos). El primer tipo, que denominamos tipo 1, son pacientes y no les importa esperar hasta que salga la versión con tapa blanda. Su disposición máxima a pagar por el libro es $V_1 = 30$ con independencia de cuándo lo compran. El otro tipo de comprador, denominados tipo 2, son impacientes y quieren leer el libro lo antes posible. La disposición a pagar máxima de este grupo es $V_2 = 50e^{-\delta t}$, donde $t$ es el momento en el tiempo en que compra y lee el libro y $\delta$ es su tasa de descuento intertemporal. Así, estos compradores obtienen una utilidad máxima de 50 si compran y leen el libro inmediatamente ($t=0$). Cada consumidor compra solo un libro.

La editorial decide lanzar inmediatamente ($t=0$) un libro de tapa dura por un precio de $P_2$ para los lectores impacientes. Luego, tiene que decidir en qué momento del tiempo ($t$) lanzar una versión en tapa blanda a un precio menor $P_1$. El costo de imprimir un libro (sea de tapa dura o blanda) es irrelevante en comparación al precio, así que se asume que el costo marginal de producción es 0. Sin embargo, la editorial tiene una tasa de descuento para sus flujos igual a $r$ por cada unidad de tiempo $t$.
\begin{parts}
\part Escriba la restricción de racionalidad individual de los compradores tipo 1 (IR1), bajo el supuesto de que estos compran el libro de tapa blanda a un precio de $P_1$.
\needspace{5\baselineskip}
\begin{solution}
La disposición máxima de pago de los compradores tipo 1 es $V_1 = 30$, por lo que el precio del libro de tapa blanda no puede ser superior a esta cantidad:
\[
  IR_1:\quad V_1 - P_1 = 30 - P_1 \geq 0 .
\]

\end{solution}

\part Escriba la restricción de racionalidad individual de los compradores tipo 2 (IR2), bajo el supuesto de que estos compran el libro de tapa dura a un precio de $P_2$ (recordando que este libro se vende en $t=0$).
\needspace{5\baselineskip}
\begin{solution}
La disposición máxima de pago de los compradores tipo 2 en $t=0$ es $V_2|_{t=0} = 50$. Por lo tanto, el precio del libro de tapa dura no puede ser superior a 50:
\[
  IR_2:\quad V_2|_{t=0} - P_2 = 50 - P_2 \geq 0 .
\]

\end{solution}

\part Escriba la condición para que el lector tipo 1 prefiera esperar hasta que salga el libro de tapa blanda en vez de comprar el de tapa dura (IC1). Determine además la condición entre los precios para que se cumpla.
\needspace{5\baselineskip}
\begin{solution}
Para que los lectores tipo 1 prefieran comprar la tapa blanda en lugar de la tapa dura, el excedente neto al comprar el libro de tapa blanda debe ser mayor que el de tapa dura:
\[
  IC_1:\quad V_1 - P_1 \geq V_1 - P_2, \quad\text{es decir,}\quad 30 - P_1 \geq 30 - P_2 .
\]
Basta con que $P_2 \geq P_1$ para que se cumpla IC1.

\end{solution}

\part Escriba la condición para que el lector tipo 2 prefiera comprar inmediatamente el libro de tapa dura en vez de esperar hasta $t$ para comprar el libro de tapa blanda (IC2), teniendo en cuenta que su disposición a pagar cae con el tiempo.
\needspace{5\baselineskip}
\begin{solution}
Para que el lector tipo 2 prefiera comprar el libro de tapa dura inmediatamente, el excedente neto al comprarlo debe ser mayor que el de esperar por la versión en tapa blanda:
\[
  IC_2:\quad V_2|_{t=0} - P_2 \geq V_2|_{t} - P_1, \quad\text{es decir,}\quad 50 - P_2 \geq 50e^{-\delta t} - P_1 .
\]

\end{solution}

\part Suponga ahora que solo son restricciones activas (se cumplen con igualdad) IR1 e IC2. Escriba entonces las ganancias de la editorial y reemplace las restricciones para expresarlas solo en función de $t$ y los parámetros del problema.
\needspace{5\baselineskip}
\begin{solution}
Las ganancias de la editorial, considerando que los ingresos por la venta de los libros de tapa blanda ocurren en el momento $t$ y que los costos marginales de producción son 0, son $\pi = P_1e^{-rt} + P_2$. Reemplazando IR1 ($P_1 = 30$) e IC2 ($P_2 = 80 - 50e^{-\delta t}$):
\[
  \pi = 30e^{-rt} + 80 - 50e^{-\delta t}.
\]

\end{solution}

\part Obtenga la condición de primer orden para el $t$ óptimo y comente su respuesta en el contexto de discriminación de segundo grado.
\needspace{5\baselineskip}
\begin{solution}
Para maximizar las ganancias, derivamos $\pi$ respecto de $t$ e igualamos a cero:
\[
\begin{gathered}
  \frac{\partial\pi}{\partial t} = -30r\,e^{-rt} + 50\delta\,e^{-\delta t} = 0 \\
  \Longrightarrow\; t^* = \frac{\ln\!\bigl(\frac{3r}{5\delta}\bigr)}{r - \delta}.
\end{gathered}
\]
Este resultado indica que la editorial opta por retrasar el lanzamiento del libro de tapa blanda para capturar más valor de los lectores impacientes, lo cual es un ejemplo de discriminación de segundo grado en precios, donde se ofrece el mismo producto en diferentes versiones y tiempos para extraer mayor excedente del consumidor.

\end{solution}
\end{parts}

%-------------------------------------------------------------------- Comente: competencia y discriminación
\needspace{8\baselineskip}
\tema{Comente: discriminar bajo competencia}
\question \textbf{Comente.} En cualquier ambiente de competencia, las empresas van a preferir competir discriminando precios que aplicando un precio uniforme.
\needspace{5\baselineskip}
\begin{solution}
Falso. En monopolio siempre será preferible discriminar a fijar un precio uniforme. En el caso de competencia imperfecta, depende de la asimetría entre las firmas respecto de cuál es el mercado donde cobrarían un mayor precio si pudieran discriminar (como en el caso de AeroWeak, que prefiere que AeroStrong \emph{no} discrimine).

\end{solution}

%-------------------------------------------------------------------- Parque de diversiones
\needspace{8\baselineskip}
\tema{Tarifa en dos partes y tarjeta prepago: parque de diversiones}
\question Un parque de diversiones puede diferenciar dos tipos de clientes: niños y adultos. Los niños tienen una curva de demanda $Q_1 = 1000 - P$ y los adultos tienen una curva de demanda $Q_2 = 500 - \frac{P}{2}$. El costo marginal del parque es \$20, igual al costo medio. Se supone que la demanda por niños y adultos está conformada por 100 niños y 100 adultos.
\begin{parts}
\part Aplicar una tarifa en dos partes, cobrando una entrada al parque diferente a cada tipo de cliente y una tarifa igual para cada juego a que entre. Determine la utilidad obtenida por el monopolio.
\needspace{5\baselineskip}
\begin{solution}
La tarifa en dos partes consiste en fijar un precio por unidad igual al costo marginal (en este caso, \$20 por juego) y luego cobrar una tarifa de entrada que capture todo el excedente del consumidor. A $P = 20$, las cantidades demandadas por cada tipo son $Q_1 = 980$ (niños) y $Q_2 = 490$ (adultos). Los excedentes son
\[
  EC_{\text{niños}} = \tfrac{1}{2}\cdot 980\cdot 980 = 480\,200, \qquad EC_{\text{adultos}} = \tfrac{1}{2}\cdot 490\cdot 980 = 240\,100 .
\]
Por lo tanto, las tarifas de entrada son iguales al excedente de cada tipo (por persona): $\frac{480\,200}{100} = 4802$ para niños y $\frac{240\,100}{100} = 2401$ para adultos. El precio por juego es igual al costo marginal, por lo que toda la ganancia proviene de las tarifas de entrada. La utilidad total del monopolio es
\[
  \pi = 480\,200 + 240\,100 = 720\,300 .
\]

\end{solution}

\part El parque está analizando implementar una tarjeta de prepago con derecho a entrar a una cantidad de juegos diferente por tipo de cliente. ¿Cuánto debería cobrar por la tarjeta para cada tipo de cliente y cuántos juegos debería permitir?
\needspace{5\baselineskip}
\begin{solution}
Se calcula la disposición a pagar total de cada grupo con $P = 20$ (excedente más el gasto en juegos): niños, $480\,200 + 20\cdot 980 = 499\,800$; adultos, $240\,100 + 20\cdot 490 = 249\,900$. Por persona, la tarjeta cuesta $4998$ para niños (por $9{,}8$ juegos) y $2499$ para adultos (por $4{,}9$ juegos). La utilidad total es
\[
  \pi = 499\,800 + 249\,900 - 20\cdot 1470 = 720\,300 ,
\]
la misma que con la tarifa en dos partes.

\end{solution}

\part Compare las utilidades de los puntos anteriores con un monopolio que no discrimina tarifas (no cobra entrada, pero cobra una tarifa única por cada juego). ¿Cuál es la utilidad en ese caso?
\needspace{5\baselineskip}
\begin{solution}
La demanda agregada es $Q = Q_1 + Q_2 = 1500 - \frac{3}{2}P$, es decir, $P = 1000 - \frac{2}{3}Q$. El monopolio maximiza $\pi = (1000 - \tfrac{2}{3}Q)Q - 20Q$, con CPO $1000 - \tfrac{4}{3}Q - 20 = 0$, de donde $Q = 735$ y $P = 1000 - \tfrac{2}{3}\cdot 735 = 510$. Entonces $Q_{\text{niños}} = 490$, $Q_{\text{adultos}} = 245$ y la utilidad total es
\[
  \pi = 510\cdot 735 - 20\cdot 735 = 360\,150 ,
\]
bastante menor que los $720\,300$ con discriminación.

\end{solution}
\end{parts}

%-------------------------------------------------------------------- Discriminación perfecta tres tipos
\needspace{8\baselineskip}
\tema{Discriminación perfecta vs precio uniforme (tres tipos de consumidores)}
\question Considere un monopolista que cuenta con $N$ consumidores divididos en 3 tipos según sus demandas inversas: Tipo 1: $P = 10 - q$; Tipo 2: $P = 6 - 2q$; Tipo 3: $P = 3 - q$. El costo marginal del monopolista es constante e igual a 1. No hay costos fijos y el monopolista sabe que el porcentaje de consumidores del tipo 1 es $\lambda$ y del tipo 2 es $(1/2 - \lambda)$, con $0 < \lambda < 1/2$.
\begin{parts}
\part Encuentre las tarifas óptimas y los beneficios si el monopolista puede discriminar perfectamente a los distintos grupos de consumidores.
\needspace{5\baselineskip}
\begin{solution}
Con discriminación perfecta, el monopolista puede cobrar a cada tipo su excedente del consumidor como tarifa fija más el costo marginal por unidad ($p = c = 1$). Las cantidades son $q_1 = 9$, $q_2 = 2{,}5$, $q_3 = 2$, y los excedentes (tarifas fijas):
\[
  T_1 = \tfrac{1}{2}\cdot 9\cdot 9 = 40{,}5, \qquad T_2 = \tfrac{1}{2}\cdot 2{,}5\cdot 5 = 6{,}25, \qquad T_3 = \tfrac{1}{2}\cdot 2\cdot 2 = 2 .
\]
La proporción de consumidores tipo 3 es $1 - \lambda - (1/2 - \lambda) = 1/2$. Los beneficios son
\[
  \pi = N\bigl[40{,}5\lambda + 6{,}25(1/2 - \lambda) + 2\cdot\tfrac{1}{2}\bigr] = N\bigl[34{,}25\lambda + 4{,}125\bigr].
\]

\end{solution}

\part ¿Qué sucede con los beneficios cuando aumenta o disminuye $\lambda$? ¿Por qué se debe este resultado?
\needspace{5\baselineskip}
\begin{solution}
$\frac{\partial\pi}{\partial\lambda} = 34{,}25N > 0$: cuando $\lambda$ aumenta, los beneficios también aumentan. Esto se debe a que al aumentar $\lambda$ hay más consumidores tipo 1 (que tienen un mayor excedente) y menos tipo 2 (que generan menos excedente).

\end{solution}

\part Suponga que el gobierno decide prohibir este tipo de discriminación y obliga al monopolista a cobrar un precio único y uniforme para todos los consumidores. ¿Cuál es el precio que cobrará el monopolista? \emph{(Hint: suponga de aquí en adelante que $\lambda = 1/4$.)}
\needspace{5\baselineskip}
\begin{solution}
Con $\lambda = 1/4$ las proporciones son $\frac14$, $\frac14$ y $\frac12$. Las demandas individuales son $q_1 = 10 - p$, $q_2 = \frac{6-p}{2}$, $q_3 = 3 - p$. Si vende a los tres tipos ($p \leq 3$), la demanda agregada es
\[
  Q(p) = N\Bigl[\tfrac{1}{4}(10 - p) + \tfrac{1}{4}\cdot\tfrac{6-p}{2} + \tfrac{1}{2}(3 - p)\Bigr] = N\,\frac{38 - 7p}{8}.
\]
La CPO de $\max (p - 1)\frac{38 - 7p}{8}$ es $38 - 7p - 7(p - 1) = 0$, es decir $p = \frac{45}{14} \approx 3{,}21$. Como $p > 3$, los consumidores tipo 3 no compran. Si vende solo a los tipos 1 y 2 ($3 < p < 6$): $Q(p) = N\,\frac{26 - 3p}{8}$, y la CPO de $\max (p - 1)\frac{26 - 3p}{8}$ es $26 - 3p - 3(p - 1) = 0$, de donde
\[
  p = \frac{29}{6} \approx 4{,}83, \qquad \pi = N\Bigl(\tfrac{29}{6} - 1\Bigr)\frac{26 - 14{,}5}{8} = \frac{529}{96}N \approx 5{,}51N .
\]
Por lo tanto, solo compran los consumidores tipo 1 y tipo 2, a un precio de $4{,}83$. (Los beneficios son mucho menores que los $12{,}7N$ de la discriminación perfecta.)

\end{solution}

\part ¿Resulta la prohibición de discriminación una buena política óptima desde el punto de vista social?
\needspace{5\baselineskip}
\begin{solution}
La política del gobierno no es una buena política porque deja fuera del mercado a los individuos tipo 3. La discriminación perfecta permite que todos los tipos con disposición a pagar mayor al costo marginal puedan consumir, maximizando así el bienestar social. Sin discriminación, se genera una pérdida irrecuperable de eficiencia al excluir a los consumidores tipo 3 del mercado (y además los consumidores de los tipos 1 y 2 pagan un precio mayor al costo marginal).

\end{solution}
\end{parts}

%-------------------------------------------------------------------- Comente: 3er grado puede disminuir beneficios
\needspace{8\baselineskip}
\tema{Comente: discriminar en tercer grado vs precio lineal}
\question \textbf{Comente.} La discriminación de precios por indicadores (o de tercer grado) puede disminuir los beneficios del monopolista si lo comparamos con precios lineales. Verdadero o falso. Justifique.
\needspace{5\baselineskip}
\begin{solution}
Falso. El monopolista discriminador siempre puede fijar el mismo precio para cada grupo de consumidores y obtener los beneficios de los precios lineales del monopolista clásico. Por lo tanto, discriminar nunca puede dar menores beneficios que el precio uniforme.

\end{solution}

%-------------------------------------------------------------------- Tercer grado: 300-4p, 240-5p
\needspace{8\baselineskip}
\tema{Discriminación de tercer grado y precio uniforme impuesto por el gobierno}
\question Una empresa monopolista vende su producto para dos mercados distintos. Las demandas en esos mercados son $Q_1 = 300 - 4P$ y $Q_2 = 240 - 5P$. La función de costos de la empresa es $C(Q) = 10Q$.
\begin{parts}
\part Suponga que el monopolista puede practicar discriminación de tercer grado (es decir, cobrar precios distintos $P_1$ y $P_2$). Calcule los precios óptimos en cada uno de los mercados, las cantidades vendidas y el beneficio total del monopolista.
\needspace{5\baselineskip}
\begin{solution}
Se resuelve $\max\ (P_1 - 10)(300 - 4P_1) + (P_2 - 10)(240 - 5P_2)$. Las CPO son $300 - 4P_1 - 4(P_1 - 10) = 0 \Rightarrow P_1 = 42{,}5$ y $240 - 5P_2 - 5(P_2 - 10) = 0 \Rightarrow P_2 = 29$. Las cantidades vendidas son $Q_1 = 300 - 4\cdot 42{,}5 = 130$ y $Q_2 = 240 - 5\cdot 29 = 95$. Los beneficios son
\[
  \pi_1 = (42{,}5 - 10)\cdot 130 = 4225, \qquad \pi_2 = (29 - 10)\cdot 95 = 1805, \qquad \pi = 6030 .
\]

\end{solution}

\part Con el objetivo de disminuir la pérdida irrecuperable de eficiencia, promoviendo un mayor bienestar social, el gobierno decide obligar al monopolista a cobrar un precio único para los dos mercados. Calcule el precio óptimo, las cantidades vendidas y el beneficio total del monopolista bajo las nuevas reglas.
\needspace{5\baselineskip}
\begin{solution}
La demanda agregada es $Q = 540 - 9P$, y se resuelve $\max\ (P - 10)(540 - 9P)$: $540 - 9P - 9(P - 10) = 0 \Rightarrow P = 35$. Las cantidades son $Q_1 = 300 - 4\cdot 35 = 160$ y $Q_2 = 240 - 5\cdot 35 = 65$ (total $225$), y el beneficio total es $\pi = (35 - 10)\cdot 225 = 5625$.

\end{solution}

\part ¿Quién gana y quién pierde con la medida del gobierno?
\needspace{5\baselineskip}
\begin{solution}
Con la medida del gobierno, el bienestar del consumidor del mercado 1 aumenta (compra más, de 130 a 160, a menor precio, de 42,5 a 35): su excedente pasa de $2112{,}5$ a $3200$. El bienestar del consumidor del mercado 2 disminuye (compra menos, de 95 a 65, a mayor precio, de 29 a 35): su excedente pasa de $902{,}5$ a $422{,}5$. El monopolista pierde (sus beneficios caen de $6030$ a $5625$). Notar que la producción total es la misma en ambos casos ($225$), pero el bienestar total sube de $9045$ a $9247{,}5$, pues con un precio uniforme el bien se asigna mejor entre los consumidores.

\end{solution}
\end{parts}

%-------------------------------------------------------------------- Movistel
\needspace{8\baselineskip}
\tema{Discriminación de segundo grado: planes de internet (Movistel)}
\question Considere el siguiente problema. Hay dos tipos de consumidores de internet: los que juegan en línea y les importa la velocidad, y los que usan internet de manera general y les importa menos la velocidad. Suponga que hay 1.000 consumidores de cada tipo. La compañía proveedora de internet Movistel (un monopolista) ha estudiado la heterogeneidad de los consumidores y quiere ofrecer dos opciones de velocidad de bajada: un plan de 10 megas por segundo y un plan de 100 megas por segundo. Obviamente, todos los consumidores prefieren mayor velocidad pero, dada su heterogeneidad, están dispuestos a pagar cantidades distintas por cada opción. Los que juegan en línea están dispuestos a pagar \$10 por el plan de 100 megas y \$1 por el plan de 10 megas. Los que usan internet de manera general están dispuestos a pagar \$5 por el plan de 100 megas y \$0,75 por el plan de 10 megas. Suponga que la empresa no observa la disponibilidad a pagar de cada consumidor (aunque conoce las propiedades de la demanda). Suponga que el costo marginal es igual a 0.
\begin{parts}
\part Suponga que Movistel ofrece un único plan. ¿Qué plan ofrecerá y a qué precio? Calcule los beneficios de la empresa.
\needspace{5\baselineskip}
\begin{solution}
Hay que probar las 4 posibilidades:
\begin{enumerate}
  \item Plan de 100 megas, compra solo el tipo 1: precio 10 y beneficio $10\,000$.
  \item Plan de 100 megas, compran los tipos 1 y 2: precio 5 y beneficio $10\,000$.
  \item Plan de 10 megas, compra solo el tipo 1: precio 1 y beneficio $1000$.
  \item Plan de 10 megas, compran los tipos 1 y 2: precio $0{,}75$ y beneficio $1500$.
\end{enumerate}
De modo que el monopolista es indiferente entre ofrecer el plan de 100 megas a un precio de 10 o de 5, ya que en ambos casos obtiene un beneficio de $10\,000$.

\end{solution}

\part Suponga que Movistel ofrece un menú de precios (es decir, distintos planes). Escriba las condiciones de incentivos y participación para ambos grupos de consumidores. Calcule los precios óptimos y los beneficios.
\needspace{5\baselineskip}
\begin{solution}
Las condiciones de participación (PC) e incentivos (IC) para ambos grupos vienen dadas por:
\begin{align*}
  PC_1:&\quad 10 - p_{100} \geq 0, & IC_1:&\quad 10 - p_{100} \geq 1 - p_{10}, \\
  PC_2:&\quad 0{,}75 - p_{10} \geq 0, & IC_2:&\quad 0{,}75 - p_{10} \geq 5 - p_{100}.
\end{align*}
Los precios óptimos son $p_{10} = 0{,}75$ (se extrae todo el excedente del tipo 2) y $p_{100} = 9{,}75$ (el máximo que hace que el tipo 1 no prefiera el plan de 10 megas: $10 - p_{100} = 1 - 0{,}75$). El beneficio obtenido es $1000\cdot 9{,}75 + 1000\cdot 0{,}75 = 10\,500$.

\end{solution}

\part Determine si Movistel ofrece el menú de precios.
\needspace{5\baselineskip}
\begin{solution}
Sí, porque el beneficio obtenido ($10\,500$) es mayor que sin diferenciación ($10\,000$).

\end{solution}
\end{parts}

\end{unidad}

\unidadtitulo{Venta atada, empaquetamiento y disuasión de entrada}
\begin{unidad}
%-------------------------------------------------------------------- Comente: venta atada vs empaquetamiento
\needspace{8\baselineskip}
\tema{Comente: venta atada vs empaquetamiento}
\question \textbf{Comente.} Una firma que es monopolio en dos productos, con el objetivo de maximizar sus beneficios, siempre preferirá realizar una venta atada por sobre una venta empaquetada (empaquetamiento) de sus productos, ya que, al limitar al máximo las opciones de las y los consumidores, es capaz de extraerles el máximo excedente posible.
\needspace{5\baselineskip}
\begin{solution}
Falso. Preferirá el empaquetamiento por sobre la venta atada, ya que en el primero tiene más instrumentos para poder llevarse un mayor excedente de los consumidores.

\end{solution}

%-------------------------------------------------------------------- Comente: beneficio consumidor
\needspace{8\baselineskip}
\tema{Comente: ¿benefician al consumidor la venta atada y el empaquetamiento?}
\question \textbf{Comente.} La venta atada y el empaquetamiento son prácticas que benefician al consumidor.
\needspace{5\baselineskip}
\begin{solution}
Incierto. Efectivamente hay un conjunto de consumidores que se ven beneficiados: quienes habrían comprado ambos bienes incluso sin una venta atada o sin empaquetamiento. Sin embargo, existe otro grupo de consumidores que se ven perjudicados por estas prácticas, en específico quienes en el caso inicial solo habrían comprado uno de los dos bienes. En el caso de venta atada, parte de dicho grupo de consumidores se verá obligado a comprar el paquete de ambos bienes y pagar un precio mayor que el precio inicial. En el caso de empaquetamiento, aquel grupo de consumidores tiene la opción de solo comprar el bien que más valoren, pero a un precio mayor que el precio inicial.

\end{solution}

%-------------------------------------------------------------------- Comente: packs
\needspace{8\baselineskip}
\tema{Comente: ofertas en formato «pack»}
\question \textbf{Comente.} Las ofertas en formato «pack» que realizan las firmas no tienen racionalidad económica, pues tendrían mayor beneficio vendiendo los productos individualmente.
\needspace{5\baselineskip}
\begin{solution}
Falso. Existen razones económicas para realizar estas ofertas. El empaquetamiento otorga una mayor flexibilidad a la empresa, pues dispone de más instrumentos u ofertas con que segmentar a los consumidores. Otras razones son la economía de escala producto de la eficiencia del bien (por ejemplo, telecomunicaciones).

\end{solution}

%-------------------------------------------------------------------- Comente: per se
\needspace{8\baselineskip}
\tema{Comente: la venta atada como práctica anticompetitiva \emph{per se}}
\question \textbf{Comente.} La venta conjunta de productos, en particular la venta atada, al imponer que la adquisición de un producto deba ser en conjunto con otro(s) producto(s), sin opción de comprarlo por separado, siempre es perseguida y sancionada por las agencias de libre competencia, dado que es una práctica anticompetitiva \emph{per se}.
\needspace{5\baselineskip}
\begin{solution}
Las ventas conjuntas de productos pueden tener distintas justificaciones. En primer lugar, pueden ocurrir por un tema de eficiencia de costos para la firma. En segundo lugar, puede ser una manifestación del poder de mercado de la empresa en un mercado oligopólico. En tercer lugar, puede ser el resultado de una práctica anticompetitiva por parte de la firma para excluir competidores. Debido a las distintas explicaciones que puede tener una venta conjunta, dicha estrategia no es una práctica anticompetitiva \emph{per se}, sino que está sujeta a la \emph{regla de la razón}. Es decir, se debe evaluar caso a caso con el objetivo de dilucidar si dicha práctica tiene una justificación como una estrategia competitiva o bien como una práctica anticompetitiva. Por lo tanto, la afirmación es falsa.

\end{solution}

%-------------------------------------------------------------------- Comente: crítica de Chicago
\needspace{8\baselineskip}
\tema{Comente: crítica de Chicago (bienes complementarios)}
\question \textbf{Comente.} Existen 2 bienes perfectamente complementarios: $A$ y $B$. Una firma produce ambos bienes, donde en el mercado del bien $A$ es un monopolio y el del bien $B$ es de competencia perfecta. Además, todos los consumidores tienen la misma valoración por los bienes y todas las firmas tienen los mismos costos para producir el bien $B$. En esta situación, el monopolio tiene incentivos a desplegar venta atada para monopolizar el otro mercado y aumentar sus beneficios.
\needspace{5\baselineskip}
\begin{solution}
Falso. Si operara como monopolio en el mercado del bien $B$, no obtiene ganancia alguna. Esto se conoce como la \emph{crítica de Chicago}, la cual señala que es irrelevante monopolizar 2 bienes en vez de 1 cuando estos son perfectos complementos. Esto es posible porque los individuos solo obtienen utilidad al consumir simultáneamente ambos bienes (el monopolista ya extrae todo el excedente con el precio del bien $A$).

\end{solution}

%-------------------------------------------------------------------- Comente: fallo créditos
\needspace{8\baselineskip}
\tema{Comente: fallo sobre ventas atadas y empaquetamiento en créditos}
\question \textbf{Comente.} Hace unos años, un fallo penalizó las ventas atadas en el mercado de los créditos. Ese mismo fallo determinó que no había problemas con el empaquetamiento, ya que esas prácticas eran beneficiosas para los consumidores.
\needspace{5\baselineskip}
\begin{solution}
El empaquetamiento es una práctica donde se cobra un precio descontado si se consumen dos productos. A pesar de ello, esta estrategia permite cobrarle más a los consumidores de mayor valoración, por lo que es una forma de discriminar precios. Es así como no es claro si efectivamente beneficia a los consumidores: los consumidores de mayor valoración por un bien en particular y baja por el otro se verán perjudicados, ya que tendrán un precio mayor al caso sin empaquetamiento.

\end{solution}

%-------------------------------------------------------------------- Comente: venta atada y entrada
\needspace{8\baselineskip}
\tema{Comente: venta atada como disuasión de entrada}
\question \textbf{Comente.} Ventas atadas es un mecanismo de disuasión de entrada. Verdadero o falso. Justifique.
\needspace{5\baselineskip}
\begin{solution}
Verdadero. Además de extraer excedente al consumidor, usar ventas atadas es una forma de bloquear la entrada de competidores. El entrante obtiene menores beneficios si el incumbente utiliza ventas atadas, lo que reduce las probabilidades de entrada.

\end{solution}

%-------------------------------------------------------------------- Venta atada y empaquetamiento (uniforme)
\needspace{8\baselineskip}
\tema{Venta separada, venta atada y empaquetamiento (valoraciones uniformes)}
\question Un monopolista vende dos productos independientes $X$ e $Y$. El costo marginal de producción de cada bien es cero. Existe un conjunto de consumidores que normalizamos a una cantidad total de 1, los que valoran los bienes en $V_x$ y $V_y$, respectivamente. Las valoraciones se distribuyen uniformemente en el espacio $[0,1]\times[0,1]$.
\begin{parts}
\part Obtenga la función de beneficios a maximizar y grafique las demandas si es que el monopolista decide vender los productos por separado. Adicionalmente, señale en el espacio $V_x$--$V_y$ qué individuos consumen de cada tipo de bien por separado, de ambos bienes y de ningún bien.
\needspace{5\baselineskip}
\begin{solution}
El monopolista maximiza
\[
  \max_{p_x,p_y}\ \pi = (1 - p_x)p_x + (1 - p_y)p_y ,
\]
por lo que $p_x = p_y = \frac{1}{2}$ y $\pi = \frac{1}{2}$. Gráficamente, compran el bien $X$ quienes tienen $V_x \geq p_x$, y el bien $Y$ quienes tienen $V_y \geq p_y$:
\begin{center}
\cuadrado{
  \draw[dashed] (0.5,0) -- (0.5,1);
  \draw[dashed] (0,0.5) -- (1,0.5);
  \node at (0.25,0.25) {Ninguno};
  \node at (0.75,0.25) {Solo $X$};
  \node at (0.25,0.75) {Solo $Y$};
  \node at (0.75,0.75) {$X$ e $Y$};
  \node[below] at (0.5,0) {$p_x$};
  \node[left] at (0,0.5) {$p_y$};
}
\end{center}

\end{solution}

\part Obtenga la función de beneficios a maximizar y grafique las demandas si es que el monopolista decide vender los productos mediante venta atada. Adicionalmente, señale en el espacio $V_x$--$V_y$ qué individuos consumen de cada tipo de bien por separado, de ambos bienes y de ningún bien. En el gráfico muestre qué consumidores se ven beneficiados y qué consumidores se ven perjudicados con esta estrategia con respecto al anterior. ¿Qué nuevos individuos consumen con respecto al caso anterior? ¿Qué individuos dejan de consumir con respecto al caso anterior?
\needspace{5\baselineskip}
\begin{solution}
Si el monopolista decide vender los productos mediante venta atada, solo se vende el paquete $X+Y$ a un precio $p_{xy}$ y compran quienes cumplen $V_x + V_y \geq p_{xy}$. La función a maximizar es
\[
\begin{gathered}
  \max_{p_{xy}}\ \pi = \Bigl(1 - \frac{p_{xy}^2}{2}\Bigr)p_{xy} \\
  \;\Rightarrow\; 1 - \frac{3}{2}p_{xy}^2 = 0 \\
  \;\Rightarrow\; p_{xy} = \sqrt{\tfrac{2}{3}} \approx 0{,}82, \\
  \qquad \pi \approx 0{,}544 > 0{,}5 .
\end{gathered}
\]
\begin{center}
\cuadrado{
  \draw[dashed] (0.5,0) -- (0.5,1);
  \draw[dashed] (0,0.5) -- (1,0.5);
  \draw[thick] (0,0.8165) -- (0.8165,0);
  \node[font=\scriptsize,align=center] at (0.3,0.3) {Ninguno};
  \node[font=\scriptsize,align=center] at (0.72,0.72) {$X$ e $Y$\\(paquete)};
}
\end{center}
Con respecto al caso anterior (líneas punteadas):
\begin{itemize}
  \item \emph{Beneficiados:} quienes compraban ambos bienes por separado ($V_x, V_y \geq \frac12$) y que ahora pagan $0{,}82$ en vez de $1$.
  \item \emph{Nuevos consumidores:} quienes tienen $V_x, V_y < \frac12$ pero $V_x + V_y \geq 0{,}82$ (compran el paquete y antes no compraban nada).
  \item \emph{Dejan de consumir:} quienes antes compraban solo un bien (alta valoración por uno, baja por el otro) y cumplen $V_x + V_y < 0{,}82$. Los que sí siguen comprando pagan ahora más que antes por el bien que valoran y además se ven obligados a comprar el otro: están perjudicados.
\end{itemize}

\end{solution}

\part Obtenga la función de beneficios a maximizar y grafique las demandas si es que el monopolista decide vender los productos mediante empaquetamiento. Adicionalmente, señale en el espacio $V_x$--$V_y$ qué individuos consumen de cada tipo de bien por separado, de ambos bienes y de ningún bien. ¿Cuál de los tres escenarios le generará mayores beneficios al monopolista? Fundamente su respuesta.
\needspace{5\baselineskip}
\begin{solution}
En el caso en el que el monopolista aplique empaquetamiento (precios $p_x$, $p_y$ y $p_{xy}$), debe considerar las restricciones de participación y autoselección para cada uno de los precios.

\emph{Demanda que solo consume $X$.} Participación: $V_x - p_x \geq 0$. Autoselección: $V_x - p_x \geq V_x + V_y - p_{xy}$, es decir, $V_y \leq p_{xy} - p_x$. Juntando ambas:
\[
  D_x = (1 - p_x)(p_{xy} - p_x), \qquad \text{y de igual forma} \qquad D_y = (1 - p_y)(p_{xy} - p_y).
\]
\emph{Demanda por el paquete.} Participación: $V_x + V_y \geq p_{xy}$. Autoselección: $V_x \geq p_{xy} - p_y$ y $V_y \geq p_{xy} - p_x$. Juntando las restricciones:
\[
  D_{xy} = (1 - p_{xy} + p_y)(1 - p_{xy} + p_x) - \tfrac{1}{2}(p_x + p_y - p_{xy})^2 .
\]
Por lo tanto, la función a maximizar es
\[
\begin{gathered}
  \max\ \pi = p_x(1 - p_x)(p_{xy} - p_x) + p_y(1 - p_y)(p_{xy} - p_y) \\ + p_{xy}\Bigl[(1 - p_{xy} + p_y)(1 - p_{xy} + p_x) - \tfrac{1}{2}(p_x + p_y - p_{xy})^2\Bigr].
\end{gathered}
\]
Resolviendo numéricamente (caso simétrico) se obtiene $p_x = p_y \approx \frac{2}{3}$ y $p_{xy} \approx 0{,}86$, con $\pi \approx 0{,}549$. Gráficamente:
\begin{center}
\cuadrado{
  \draw (0.6667,0) -- (0.6667,0.1967) -- (1,0.1967);
  \draw (0,0.6667) -- (0.1967,0.6667) -- (0.1967,1);
  \draw (0.1967,0.6667) -- (0.6667,0.1967);
  \node[font=\scriptsize] at (0.28,0.28) {Ninguno};
  \node[font=\scriptsize] at (0.83,0.09) {Solo $X$};
  \node[font=\scriptsize] at (0.1,0.9) {Solo $Y$};
  \node[font=\scriptsize,align=center] at (0.7,0.7) {Paquete\\$X$ e $Y$};
}
\end{center}
El escenario que mayores beneficios reporta al monopolista es este último ($0{,}549 > 0{,}544 > 0{,}5$), debido a que, al tener más instrumentos disponibles, el monopolista puede crear mejores ofertas para sus distintos grupos de consumidores.

\end{solution}
\end{parts}

%-------------------------------------------------------------------- Bundling con tabla
\needspace{8\baselineskip}
\tema{Bundling puro y mixto con tres tipos de consumidores}
\question Considere la siguiente tabla, que contiene las valoraciones para los bienes $X$ e $Y$. Suponga que el costo marginal de $X$ es 1 y el costo marginal de $Y$ también es 1.
\begin{center}
\begin{tabular}{lcc}
  & Producto $X$ & Producto $Y$ \\ \hline
  Consumidor tipo 1 & 4 & 3 \\
  Consumidor tipo 2 & 3 & 3 \\
  Consumidor tipo 3 & 0 & 4
\end{tabular}
\end{center}
\begin{parts}
\part Encuentre el precio de monopolio óptimo en caso de ventas atadas o \emph{bundling} puro, es decir, precio de $X+Y$.
\needspace{5\baselineskip}
\begin{solution}
Bajo ventas atadas el monopolio vende $X+Y$ o no vende nada.
\begin{itemize}
  \item Si $P(X+Y) = 7$, solo compra el consumidor tipo 1 y los beneficios son $\pi = 7 - 1 - 1 = 5$.
  \item Si $P(X+Y) = 6$, compran los consumidores tipo 1 y 2, y los beneficios son $\pi = 2\cdot(6 - 1 - 1) = 8$.
  \item Si $P(X+Y) = 4$, compran los 3 tipos de consumidores, y los beneficios son $\pi = 3\cdot(4 - 1 - 1) = 6$.
\end{itemize}
Luego, la mejor estrategia es $P(X+Y) = 6$ con beneficios $\pi = 8$.

\end{solution}

\part Encuentre el precio de monopolio óptimo en caso de empaquetamiento o \emph{bundling} mixto, es decir, vender por separado y en conjunto. ¿Cuál de las estrategias, ventas atadas puras o mixtas, conlleva el mayor beneficio para el monopolista?
\needspace{5\baselineskip}
\begin{solution}
Se maximizan los beneficios vendiendo el paquete a un precio $P(X+Y) = 6$ y el bien $Y$ a $P(Y) = 4$. Note que el precio de $X$ solo es irrelevante, pues ningún tipo de consumidor compra el bien $X$ solo (se puede fijar, por ejemplo, un precio alto). El consumidor tipo 3 compra solo $Y$ (que valora en 4), y los tipos 1 y 2 compran el paquete. Los beneficios son
\[
  \pi = 2\cdot(6 - 1 - 1) + (4 - 1) = 11 > 8 .
\]
Por lo tanto, el \emph{bundling} mixto conlleva el mayor beneficio para el monopolista.

\end{solution}
\end{parts}

%-------------------------------------------------------------------- Entrada y venta atada
\needspace{8\baselineskip}
\tema{Entrada y venta atada: Agro S.A. y Siembra S.A.}
\question La empresa agrícola Agro S.A. siembra, cosecha y vende manzanas y naranjas, y actualmente posee el monopolio de ambos mercados. Se sabe además que las demandas que se enfrentan en estos mercados son de la forma $Q(P) = 4 - P$ en cada mercado (se puede ver así que las máximas disposiciones a pagar son de 4 en cada mercado). El costo marginal de producción tanto para manzanas como para naranjas es de 1. Agro se ha enterado de que Siembra S.A. está evaluando entrar al mercado de las manzanas, quien de hacerlo lo haría con un costo marginal igual a cero, pero debe enfrentar un costo de entrada $F = 2$ para hacerlo.
\begin{parts}
\part Evalúe si una estrategia en donde Agro S.A. vende una canasta de manzanas y naranjas podría disuadir la entrada de Siembra S.A.
\needspace{5\baselineskip}
\begin{solution}
El \emph{timing} del juego es:
\begin{itemize}
  \item $T=1$: Siembra decide si realiza su entrada al mercado o no.
  \item $T=2$: Agro y Siembra compiten en precios.
\end{itemize}
Resolviendo por inducción hacia atrás, en $T=2$ hay 2 posibles casos.

\textbf{Caso 1: Agro acomoda la entrada (no hay venta atada).} Siembra entra en el mercado de las manzanas quedándose con éste, al ser más eficiente y al estar compitiendo en precios. Agro sigue como monopolio en el mercado de las naranjas. Calculamos los beneficios de cada empresa:
\[
  \max_Q\ \pi_A = Q(4 - Q) - Q \;\Rightarrow\; 4 - 2Q - 1 = 0 \;\Rightarrow\; Q = \tfrac{3}{2}, \qquad \pi_A = \tfrac{9}{4}.
\]
Siembra cobra un precio $P = 1 - \epsilon$ y vende $Q = 3$, por lo que sus beneficios son $\pi_S = 3$.

\textbf{Caso 2: Agro impide la entrada con venta atada.} Para que la canasta de manzanas y naranjas sea preferida, la utilidad asociada (excedente: disposición máxima a pagar $V$ menos el precio) tiene que ser igual o mayor a la utilidad de comprar las manzanas por sí solas: $U_C \geq U_M$, es decir
\[
  V_N + V_M - P_C \geq V_M - P_M \quad\Longleftrightarrow\quad P_C \leq V_N + P_M .
\]
El precio de la canasta tiene que ser menor a la disposición a pagar por las naranjas más el precio de las manzanas. Debido a la competencia en precios en el mercado de las manzanas, Siembra cobrará a lo más su costo marginal, que es $0$. Por lo que el precio de la canasta tiene que ser menor a la disposición máxima a pagar por las naranjas (4). La demanda de la canasta es entonces $Q(P_C) = 16 - \frac{P_C^2}{2}$. Agro fija un precio de la canasta que cumple esta condición, de modo que el precio de las manzanas de Siembra debe ser $0$ para competir contra la canasta (la utilidad de comprar manzanas, para el consumidor, es la mayor posible dados los costos de producción de Siembra), y sus beneficios son $0$.

\emph{Nota:} la pauta original calculaba el precio óptimo de la canasta con la CPO de $\pi_A = (P_C - 1)\bigl(16 - \frac{P_C^2}{2}\bigr)$, obteniendo $P_C \approx 3{,}62$. Notar que el costo de producir la canasta es 2 (1 por cada bien); con ese costo la CPO lleva a $P_C = 4$, el borde del rango. En ambos casos se cumple la condición $P_C \leq 4$.

\textbf{En $T=1$ Siembra decide si entra.}
\begin{itemize}
  \item En el caso 1, donde Agro se acomoda: $\pi_S - F = 3 - 2 = 1 > 0$. En este caso sí habrá entrada.
  \item En el caso 2, donde Agro realiza una venta atada para disuadir la entrada: $\pi_S - F = 0 - 2 = -2 < 0$. No habrá entrada de Siembra en el mercado.
\end{itemize}
Por lo tanto, la venta atada sí disuade la entrada. Queda como propuesto ver si a Agro le conviene o no disuadir la entrada, comparando los beneficios asociados a cada caso.

\end{solution}
\end{parts}

%-------------------------------------------------------------------- Hotel Isla de Pascua
\needspace{8\baselineskip}
\tema{Bundling en un hotel (habitación y desayuno)}
\question Usted es el manager del hotel más grande de la Isla de Pascua y debe elegir los paquetes que ofrece a sus huéspedes. La tabla muestra la disposición a pagar de los turistas de tipo 1 y de tipo 2, tanto para una habitación (H) como para el desayuno (D). Además, se tiene información del número esperado de huéspedes de cada tipo y el costo marginal de proveer cada servicio.
\begin{center}
\begin{tabular}{lccc}
  Tipo de huésped & Habitación (H) & Desayuno (D) & N.º esperado de huéspedes \\ \hline
  Tipo 1 & 100 dólares & 5 dólares & 200 \\
  Tipo 2 & 60 dólares & 10 dólares & 800
\end{tabular}
\end{center}
Costo marginal: 40 dólares para H y 2 dólares para D.
\begin{parts}
\part Calcule el precio de la habitación $P_H$ y del servicio de desayuno $P_D$ que maximiza los beneficios del hotel. Encuentre los beneficios de vender cada servicio separadamente.
\needspace{5\baselineskip}
\begin{solution}
\textbf{Habitación.} El hotel puede poner el precio más bajo ($P_H = 60$) y vender a todos (1000 huéspedes): $\pi_H = (60 - 40)\cdot 1000 = 20\,000$. Si el hotel pone el precio más alto, $P_H = 100$, vende solo a los de tipo 1: $\pi_H = (100 - 40)\cdot 200 = 12\,000$. Por lo tanto, $P_H = 60$.

\textbf{Desayuno.} Si el precio es $P_D = 5$, vende a todos: $\pi_D = (5 - 2)\cdot 1000 = 3000$. Si el precio es $P_D = 10$, solo compran los de tipo 2: $\pi_D = (10 - 2)\cdot 800 = 6400$. El hotel obtiene mayor utilidad con $P_D = 10$.

Los beneficios totales son $\pi = 20\,000 + 6400 = 26\,400$ dólares.

\end{solution}

\part Suponga que el hotel ahora ofrece una habitación con desayuno incluido. Encuentre el precio del paquete $P_{HD}$ que maximiza los beneficios del hotel. Calcule los beneficios que obtiene el hotel al ofrecer el paquete.
\needspace{5\baselineskip}
\begin{solution}
Las disposiciones a pagar por el paquete son $105$ (tipo 1) y $70$ (tipo 2), y el costo marginal del paquete es $42$. El hotel puede vender el paquete a $P_{HD} = 70$ dólares, en cuyo caso le vende a todos y los beneficios son $\pi_{HD} = (70 - 42)\cdot 1000 = 28\,000$. Si el precio del paquete es $P_{HD} = 105$, vende solo a los que están dispuestos a pagar ese precio y los beneficios son $\pi_{HD} = (105 - 42)\cdot 200 = 12\,600$.

\end{solution}

\part ¿Qué le conviene más al hotel? ¿Ofrecer ventas atadas u ofrecer una habitación y el desayuno de manera separada?
\needspace{5\baselineskip}
\begin{solution}
La estrategia óptima es ofrecer la habitación con desayuno (paquete) a un precio $P_{HD} = 70$ dólares, con beneficios $\pi_{HD} = 28\,000 > 26\,400$.

\end{solution}
\end{parts}

%-------------------------------------------------------------------- Comente: bundling mixto vs puro
\needspace{8\baselineskip}
\tema{Comente: bundling puro vs bundling mixto}
\question \textbf{Comente.} Un monopolio multiproducto siempre preferirá realizar una venta atada por sobre una venta empaquetada (\emph{bundling} mixto) de sus productos, ya que, al limitar al máximo las opciones de las y los consumidores, es capaz de extraerles el máximo excedente posible.
\needspace{5\baselineskip}
\begin{solution}
Falso. Por lo general, el \emph{bundling} mixto tiene mayor capacidad de llevarse un mayor excedente de los consumidores: al vender en conjunto y por separado tiene más herramientas para extraer excedente. Considere también que al hacer \emph{bundling} mixto se relajan las restricciones de participación para ciertos individuos que no valoran mucho un bien relativo al otro. Por esto mismo, el \emph{bundling} puro puede llevar a perder consumidores; el uso de venta atada se justificaría, por ejemplo, para evitar la entrada de una firma.

\end{solution}

%-------------------------------------------------------------------- Comente: por qué ventas atadas
\needspace{8\baselineskip}
\tema{Comente: ¿por qué las empresas usan ventas atadas?}
\question \textbf{Comente.} Si ofrecer ventas atadas reduce el número de productos de una empresa, ¿por qué las empresas utilizan este mecanismo de precios?
\needspace{5\baselineskip}
\begin{solution}
Las empresas utilizan ventas atadas como un mecanismo de discriminación de precios, pues puede permitir aumentar utilidades al incorporar clientes que valoran mucho uno de los productos y por tanto compran el conjunto (y pagan por un bien que por separado no hubieran comprado). Cabe destacar que pueden perder clientes que no están dispuestos a comprar el paquete.

Por otra parte, las ventas atadas son un mecanismo de disuasión de entrada: la competencia que quiere entrar con un producto no es competitiva al no vender en paquete, y el incumbente fija precios de manera que a los clientes les convenga el paquete y el entrante se queda sin mercado.

\end{solution}

%-------------------------------------------------------------------- Venta atada Internet y TV
\needspace{8\baselineskip}
\tema{Venta atada, precios individuales y Doble Play (Internet y TV)}
\question Una empresa de telecomunicaciones vende dos productos: Internet y Televisión de Pago. Existen tres tipos de consumidores, cuyas máximas disposiciones a pagar por cada producto, para el consumidor representativo del grupo, son las que se indican a continuación. Existe igual número de consumidores en cada tipo. Considere un costo de producción igual a 0.
\begin{center}
\begin{tabular}{cccc}
  Consumidor & Internet & TV Pago & Suma \\ \hline
  A & 120 & 70 & 190 \\
  B & 90 & 100 & 190 \\
  C & 40 & 145 & 185
\end{tabular}
\end{center}
\begin{parts}
\part Si la empresa fija precio por productos individuales, ¿cuáles serían los precios fijados en cada uno de ellos?
\needspace{5\baselineskip}
\begin{solution}
Se venden los productos por separado, fijando precios que maximicen los beneficios.
\begin{itemize}
  \item \emph{Internet.} Hay tres opciones: que solo compren los de alta valoración ($P = 120$, beneficio $120\cdot 1$), que compren los de alta y media ($P = 90$, beneficio $90\cdot 2 = 180$), o que compren todos ($P = 40$, beneficio $40\cdot 3 = 120$). Se decide por $P_{Int} = 90$, sirviendo a $2/3$ del mercado.
  \item \emph{TV Pago.} $P = 145$: $145$; $P = 100$: $200$; $P = 70$: $210$. El precio debe apelar a todos los grupos, $P_{TV} = 70$.
\end{itemize}
Los beneficios son $\pi = \pi_{Int} + \pi_{TV} = 90\cdot 2 + 70\cdot 3 = 390$.

\end{solution}

\part La empresa decide realizar una venta atada, es decir, ofrecer los dos bienes en conjunto sin posibilidad de compra individual. ¿Cuál sería el precio que se cobraría por los dos bienes? ¿Le conviene esta estrategia a la empresa respecto de la venta individual?
\needspace{5\baselineskip}
\begin{solution}
Imponemos la restricción de participación para cada grupo: la valoración conjunta debe superar al precio para que haya venta. Las valoraciones conjuntas son $190$ (A), $190$ (B) y $185$ (C). Ante el \emph{trade-off} entre volumen y margen, toca decidir entre $190\cdot 2 = 380$ o $185\cdot 3 = 555$, donde claramente se impone $P = 185$ con beneficios $555$.

En comparación con el caso anterior hay mayores utilidades ($555 > 390$). Esto se debe a que los que valoraban mucho el Internet (grupo A) tienen que comprar también TV, mientras que los que valoraban más la TV (grupo C) tienen que comprar también Internet.

\end{solution}

\part La empresa decide vender los productos individuales, pero si un suscriptor desea comprar ambos productos (Doble Play), se le realiza un descuento respecto de la compra de ambos bienes por separado. Obtenga los tres precios óptimos que aplicaría la empresa (dos precios de productos separados y el precio Doble Play). Obtenga los beneficios totales de la empresa y compárelos con las dos preguntas anteriores. Comente.
\needspace{5\baselineskip}
\begin{solution}
Para planificar el \emph{bundling} mixto representamos todas las valoraciones por producto y paquete. Si alguien compra los bienes por separado, tienen que ser los que más lo valoran: $P_{Int} = 120$ (grupo A) y $P_{TV} = 145$ (grupo C). El precio del paquete tiene que considerar las restricciones de incentivo para que los de mayor valoración por TV e Internet sigan comprando los bienes separados.

El arte está en determinar el precio del paquete. Al fijar $P_{Int+TV} = 190 - \epsilon$, los grupos A y B compran el paquete: capturamos al consumidor intermedio B y sacamos mayor excedente de A, que no valoraba tanto la TV. Con este \emph{bundling} mixto tenemos $\pi = 190\cdot 2 + 145 = 525$.

Vemos que la venta atada pura ($555$) es mejor en este caso. Esto puede explicarse por la gran diferencia en valoración entre bienes y grupos, por lo que el paquete tiende a sacar mayor excedente al juntar el bien muy valorado con el menos valorado. Adicionalmente, la valoración conjunta en los tres grupos es similar, por lo que bajar 5 unidades el precio para capturar al restante $1/3$ es altamente beneficioso. (Notar que el \emph{bundling} mixto óptimo nunca puede ser peor que el puro, pues puede replicarlo fijando precios individuales muy altos; aquí el óptimo mixto coincide con no vender por separado.)

\end{solution}
\end{parts}

%-------------------------------------------------------------------- Comente: menú restaurante
\needspace{8\baselineskip}
\tema{Comente: menú y platos a la carta}
\question \textbf{Comente.} Cuando una firma como un restaurante ofrece un menú que incluye varios platos por un precio único por el conjunto, pero también ofrece los platos por separado a la carta, entonces todos los consumidores siempre estarán mejor que si no se ofrece el menú, ya que siempre es mejor tener más opciones que menos.
\needspace{5\baselineskip}
\begin{solution}
Falso. El menú de opciones le permite al restaurante subir los precios de los platos a la carta, extrayendo más excedente de los consumidores que valoran mucho alguno de estos platos en forma individual pero no los otros del menú, sin perder la demanda de los consumidores con gustos intermedios por todos los platos. Así, hay algunos consumidores que ganan al ofrecerse el menú (los de gustos intermedios) y otros que pierden (los de gustos extremos por algún plato en particular).

\end{solution}

%-------------------------------------------------------------------- Triple pack
\needspace{8\baselineskip}
\tema{Triple pack vs servicios por separado (Internet, TV y telefonía)}
\question Una empresa de telecomunicaciones ofrece tres servicios: Internet, Televisión y Telefonía. Existen tres tipos de consumidoras, clasificadas como A, B y C. La máxima disposición a pagar por categoría de servicio y según tipo de consumidora es la siguiente:
\begin{center}
\begin{tabular}{cccc}
  Consumidoras & Internet & Televisión & Telefonía \\ \hline
  A & 13 & 6 & 3 \\
  B & 7 & 12 & 16 \\
  C & 6 & 25 & 11
\end{tabular}
\end{center}
Suponga que hay un mismo número de consumidoras por cada categoría. Todos los costos fijos y variables son cero.
\begin{parts}
\part Si la empresa vende solamente un único plan que incluye todos los servicios, ¿cuál sería el precio que cobra por el triple pack? ¿Cuánto sería el beneficio obtenido?
\needspace{5\baselineskip}
\begin{solution}
Las valoraciones del triple pack son $22$ (A), $35$ (B) y $42$ (C). Normalizando el número de consumidoras de cada tipo a 1:
\[
  P = 42 \Rightarrow 1\text{ compradora}: 42; \qquad
  P = 35 \Rightarrow 2: 70; \qquad
  P = 22 \Rightarrow 3: 66 .
\]
Por lo tanto, la firma fijará $P = 35$ para el triple pack y obtendrá beneficios de $70$.

\end{solution}

\part La empresa está considerando vender los servicios por separado. ¿Cuál sería el precio cobrado por cada servicio? ¿Le conviene cambiar su estrategia de precios respecto al plan único? Explique.
\needspace{5\baselineskip}
\begin{solution}
Analizamos precios, cantidades y beneficios de cada servicio:
\begin{itemize}
  \item \emph{Internet:} $P=13 \to 13$; $P=7 \to 14$; $P=6 \to 18$. Conviene $P = 6$ (beneficio $18$).
  \item \emph{Televisión:} $P=25 \to 25$; $P=12 \to 24$; $P=6 \to 18$. Conviene $P = 25$ (beneficio $25$).
  \item \emph{Telefonía:} $P=16 \to 16$; $P=11 \to 22$; $P=3 \to 9$. Conviene $P = 11$ (beneficio $22$).
\end{itemize}
Sumando los beneficios, el total es $18 + 25 + 22 = 65$. Al compararlo con el beneficio del triple pack ($70$), concluimos que no le conviene a la empresa cambiar su estrategia de precios. La razón es que al ofrecer un paquete único, la firma puede capturar mayor disposición a pagar de consumidores con preferencias diversas, mientras que al desagregar los servicios se pierde esta ventaja.

\end{solution}
\end{parts}

%-------------------------------------------------------------------- HP impresoras y toner
\needspace{8\baselineskip}
\tema{Bundling puro y mixto: impresoras y cartuchos de tóner (HP)}
\question HP vende dos productos: impresoras y cartuchos de tóner. Hay tres consumidores en el mercado. El consumidor 1 necesita imprimir con mucha frecuencia, por lo que valora mucho ambos bienes; el consumidor 2 necesita imprimir con menos frecuencia, por lo que valora menos ambos bienes; y el consumidor 3 ya tiene impresora, por lo que solo valora el tóner, y lo valora mucho porque ya ha realizado el gasto en la impresora. Las valoraciones de los consumidores aparecen en la siguiente tabla:
\begin{center}
\begin{tabular}{lcc}
  & Impresora & Tóner \\ \hline
  Consumidor tipo 1 & 8 & 4 \\
  Consumidor tipo 2 & 6 & 3 \\
  Consumidor tipo 3 & 0 & 8
\end{tabular}
\end{center}
Para simplificar, el costo marginal de producir ambos bienes es cero. Calcule los precios óptimos y los beneficios de HP en las siguientes situaciones:
\begin{parts}
\part HP vende ambos bienes de forma separada.
\needspace{5\baselineskip}
\begin{solution}
Hay dos posibilidades para el precio de la impresora: $P(I) = 8$ (solo compra el consumidor 1; $\pi = 8$) o $P(I) = 6$ (compran los consumidores 1 y 2; $\pi = 12$). Hay tres posibilidades para el precio del tóner: $P(T) = 8$ (solo compra el consumidor 3; $\pi = 8$), $P(T) = 4$ (compran 1 y 3; $\pi = 8$) o $P(T) = 3$ (compran los tres; $\pi = 9$). Luego, los precios óptimos son $P(I) = 6$ y $P(T) = 3$, y las utilidades totales son $\pi = 12 + 9 = 21$.

\end{solution}

\part HP aplica ventas atadas y solo vende los dos bienes de forma conjunta. ¿Qué preferiría HP: vender ambos bienes separadamente o de forma conjunta?
\needspace{5\baselineskip}
\begin{solution}
Las valoraciones del paquete son $12$, $9$ y $8$ (consumidores 1, 2 y 3). Hay tres posibilidades: $P(I+T) = 12$ (compra solo el 1; $\pi = 12$), $P(I+T) = 9$ (compran 1 y 2; $\pi = 18$), $P(I+T) = 8$ (compran los tres; $\pi = 24$). Por lo tanto, el precio óptimo es $P(I+T) = 8$ y las utilidades son $\pi = 24$. HP preferirá vender los bienes en conjunto, ya que las utilidades son mayores ($24 > 21$).

\end{solution}

\part HP también puede aplicar empaquetamiento y vender los bienes tanto de forma conjunta como separada. ¿Cuál será la elección final de HP? ¿Vender los bienes de forma separada, aplicar ventas atadas o empaquetamiento?
\needspace{5\baselineskip}
\begin{solution}
Con $P(I+T) = 9$, $P(I) = 8$ y $P(T) = 8$: los consumidores 1 y 2 compran el paquete (el consumidor 1 obtiene un excedente de $3$ con el paquete y $0$ con la impresora sola; el consumidor 2 no puede comprar la impresora a $8$) y el consumidor 3 compra el tóner. Las utilidades son $\pi = 9 + 9 + 8 = 26$. (Cualquier precio de la impresora entre $6$ y $8$ funciona; debe ser al menos $6$, porque si fuera menor el consumidor 2 compraría solo la impresora, que le da un excedente positivo en lugar de cero.)

Notar que el esquema $P(I+T) = 12$, $P(I) = 6$, $P(T) = 8$ no es compatible en incentivos: con $P(I) = 6$, el consumidor 1 prefiere comprar solo la impresora (excedente $2$) al paquete (excedente $0$), por lo que HP obtendría solo $6 + 6 + 8 = 20$.

Por consiguiente, la elección final de HP será el empaquetamiento (vender por separado y en conjunto), ya que las utilidades ($\pi = 26$) son mayores que si vendiera solo por separado ($\pi = 21$) o ventas atadas ($\pi = 24$).

\end{solution}
\end{parts}

\end{unidad}

\unidadtitulo{Entrada y amenazas creíbles}
\begin{unidad}
%-------------------------------------------------------------------- Comente: bloquear entrada
\needspace{8\baselineskip}
\tema{Comente: ¿siempre conviene bloquear la entrada?}
\question \textbf{Comente.} En los modelos de entrada, siempre la mejor estrategia de la incumbente será bloquear la entrada.
\needspace{5\baselineskip}
\begin{solution}
Falso. La firma tiene que analizar ambas decisiones en torno a la entrada de una nueva firma al mercado. Si el beneficio de acomodar la entrada es mayor al beneficio de bloquearla, la firma incumbente decidirá acomodar la entrada. En este sentido, es muy importante que la estrategia elegida sea creíble.

\end{solution}

%-------------------------------------------------------------------- Comente: mercados desafiables
\needspace{8\baselineskip}
\tema{Comente: mercados desafiables}
\question \textbf{Comente.} La teoría de mercados desafiables nos diría que, ante pequeñas barreras de entrada, un mercado monopólico u oligopólico podría llegar a un equilibrio más similar al de competencia perfecta.
\needspace{5\baselineskip}
\begin{solution}
Verdadero. La existencia de pocas barreras de entrada genera una amenaza de competencia por parte de otras firmas, forzando a las incumbentes a comportarse de una manera más competitiva (es decir, no hay poder de mercado independiente de la concentración).

\end{solution}

%-------------------------------------------------------------------- Entrada con capacidad
\needspace{8\baselineskip}
\tema{Entrada con restricciones de capacidad}
\question Suponga que un mercado está caracterizado por la siguiente función de demanda lineal: $Q = 12 - P$. La firma incumbente ($i=1$) posee un costo marginal $c = 2$. Existe un potencial entrante ($i=2$) que posee el mismo costo marginal que la incumbente, pero si entra debe pagar un costo fijo de $F = 2$. Suponga que tanto la incumbente como el entrante poseen una restricción de capacidad de $k_i = 2$. En $t=1$ la firma entrante decide si ingresa o no al mercado; en caso de que ingrese, desembolsa $F$. Si hay entrada, en $t=2$ las firmas compiten à la Bertrand. Determine el precio, las cantidades y los beneficios de equilibrio si existiera entrada. ¿Existirá entrada en este mercado?
\needspace{5\baselineskip}
\begin{solution}
Dado que hay restricciones de capacidad, la decisión de la entrante en $t=2$ se reduce a ser \emph{demandante residual} o \emph{recortador de precios}.

\textbf{Demandante residual.} La demanda residual es la demanda descontada por la cantidad que alcance a vender la incumbente: $Q_r = 12 - 2 - P = 10 - P$. Se resuelve
\[
  \max_{P_2}\ \pi_2 = (10 - P_2)(P_2 - 2)
  \;\Rightarrow\; 10 - 2P_2 + 2 = 0
  \;\Rightarrow\; P_2 = 6, \quad Q_r = 4 .
\]
Dado que $Q_r = 4 > k_2 = 2$, la cantidad pasa a ser $2$, y el precio se ajusta para que la demanda residual sea igual a la capacidad: $10 - P_2 = 2$, es decir $P_2 = 8$. Los beneficios son $\pi_2 = (8 - 2)\cdot 2 = 12$.

\textbf{Recortar precios.} Si la entrante recorta el precio de la incumbente ($P_2 = P_1 - \epsilon$) vende toda su capacidad, con beneficios $\pi_2 = (P_1 - 2)\cdot 2$. Conviene recortar precios en lugar de ser demandante residual si $(P_1 - 2)\cdot 2 > 12$, es decir, si $P_1 > 8$. Por lo tanto, siempre que la incumbente ponga un precio mayor a 8, la firma 2 entra recortando precios, de modo que el precio de equilibrio es $P^* = 8$.

\textbf{Equilibrio.} Con $P^* = 8$, la cantidad de mercado es $Q = 12 - 8 = 4 = k_1 + k_2$: cada firma vende $q_1 = q_2 = 2$ y obtiene beneficios $\pi_1 = \pi_2 = (8-2)\cdot 2 = 12$.

\textbf{Decisión de entrada ($t=1$).} La entrante obtiene $12 - F = 12 - 2 = 10 > 0$ al entrar, y su costo de oportunidad es 0, por lo que sí habrá entrada.

\end{solution}

%-------------------------------------------------------------------- Amenazas creíbles
\needspace{8\baselineskip}
\tema{Amenazas creíbles a la entrada (capacidad)}
\question Suponga una firma incumbente que se ve enfrentada a la entrada de una nueva firma. La demanda del mercado es $P = 20 - Q$. Además, las firmas tienen costos marginales iguales a cero y tienen una capacidad fija. La firma entrante, en caso de entrar, empieza con una capacidad igual a 4. La incumbente no tiene costos de aumentar su capacidad. El precio estará dado por las capacidades ofrecidas, es decir, la cantidad será igual a la capacidad de cada firma y luego con eso se fija el precio. El juego funciona de la siguiente manera:
\begin{itemize}
  \item $t=1$: la firma entrante decide si entrar o no. En caso de no entrar, la firma entrante tiene un proyecto alternativo donde obtiene beneficios iguales a 16.
  \item $t=2$: la firma incumbente decide si actuar agresivamente o no. En caso de actuar agresivamente, elige una capacidad que minimiza las utilidades de la firma entrante. En caso de no actuar agresivamente, maximiza sus utilidades.
\end{itemize}
\begin{parts}
\part Plantee el juego y el árbol de decisión.
\needspace{5\baselineskip}
\begin{solution}
Para resolver el juego hay que calcular los pagos de las distintas alternativas.
\begin{itemize}
  \item \emph{La entrante no entra:} recibe $16$ por su proyecto alternativo y la incumbente es monopolio, con $\pi_I = 100$ ($Q = 10$, $P = 10$).
  \item \emph{La entrante ingresa y la incumbente actúa agresivamente:} la incumbente elige una capacidad que minimiza los beneficios de la entrante, es decir, que lleva el precio a cero: $P = 20 - k_I - 4 = 0$, de modo que $k_I = 16$, y ambas obtienen beneficios iguales a $0$.
  \item \emph{La entrante ingresa y la incumbente no actúa agresivamente:} la incumbente maximiza sus beneficios dado $k_E = 4$: $\max_{k_I}\ (20 - k_I - 4)k_I$, de donde $k_I = 8$, $P = 8$, $\pi_E = 8\cdot 4 = 32$ y $\pi_I = 8\cdot 8 = 64$.
\end{itemize}
\begin{center}
\begin{tikzpicture}[x=1cm,y=1cm,font=\small]
  \node (e) at (0,0) {Entrante};
  \draw (e) -- (-2.8,-1.3) node[midway,above left] {No entra} node[below,align=center] {$\pi_E = 16$\\$\pi_I = 100$};
  \node (i) at (2.8,-1.3) {Incumbente};
  \draw (e) -- (i) node[midway,above right] {Entra};
  \draw (i) -- (1.2,-3) node[midway,left] {Agresiva} node[below,align=center] {$\pi_E = 0$\\$\pi_I = 0$};
  \draw (i) -- (4.4,-3) node[midway,right] {No agresiva} node[below,align=center] {$\pi_E = 32$\\$\pi_I = 64$};
\end{tikzpicture}
\end{center}

\end{solution}

\part Encuentre el equilibrio y explique el resultado. ¿Qué rol juega el beneficio del proyecto alternativo de la firma entrante?
\needspace{5\baselineskip}
\begin{solution}
El equilibrio se resuelve por inducción hacia atrás. En la última etapa, la incumbente preferirá no actuar agresivamente, ya que obtiene mayores beneficios así ($64 > 0$). Sabiendo esto, la firma entrante preferirá ingresar al mercado, ya que la competencia no agresiva le reporta un beneficio mayor que su proyecto alternativo ($32 > 16$). A la incumbente le gustaría amenazar con actuar agresivamente para evitar la entrada y obtener un beneficio de $\pi_I = 100$, pero no puede hacer esta amenaza creíble. La razón es que, una vez que se produce la entrada, actuar agresivamente no es la mejor opción para la incumbente.

El beneficio del proyecto alternativo de la entrante ($16$) representa el costo de oportunidad de la firma al decidir entrar o no: si la firma entra y obtiene beneficios nulos, habría perdido el beneficio alternativo de 16.

\end{solution}

\part Suponga ahora que la firma incumbente puede elegir su capacidad antes de la entrada. ¿Cuál será la mejor opción de la incumbente en este caso?
\needspace{5\baselineskip}
\begin{solution}
En este caso la incumbente puede fijar una capacidad para desincentivar la entrada de la competencia. Sin embargo, esta estrategia solo es conveniente si los beneficios son mayores que permitir la entrada, lo que anteriormente generaba una utilidad de $\pi_I = 64$. El nivel de capacidad que desincentiva la entrada es tal que los beneficios de la entrante sean menores que su proyecto alternativo:
\[
  \pi_E = (20 - k_I - 4)\cdot 4 < 16 \quad\Longrightarrow\quad k_I > 12 .
\]
Por lo tanto, la incumbente debe fijar una capacidad de al menos 13 unidades. Si lo hace, los beneficios serán los monopólicos menos los costos de inversión en capacidad (la pauta supone un costo de $3$ por unidad de capacidad):
\[
  \pi_I = 100 - 3\cdot 13 = 61 .
\]
Dado que estos beneficios son menores que $\pi_I = 64$, la mejor opción de la incumbente será permitir la entrada y escoger su capacidad posterior para maximizar sus beneficios.

\end{solution}
\end{parts}

%-------------------------------------------------------------------- Demasiadas empresas
\needspace{8\baselineskip}
\tema{¿Demasiadas empresas en la industria? (costos fijos y Cournot)}
\question \textbf{Comente.} Mientras más empresas haya en una industria más eficiente será, puesto que un número reducido de estas podría llevar a mayor poder de mercado reduciendo el bienestar de los consumidores. Verdadero o falso.
\needspace{5\baselineskip}
\begin{solution}
Depende. Si bien la intuición es correcta, dado que a mayor cantidad de empresas podría reducirse el índice de Lerner (y con ello el poder de mercado), en la realidad hay fricciones que considerar. En gran parte de las industrias los costos fijos (de entrada o recurrentes) que enfrenta cada firma son considerables: mientras más firmas haya, más costos fijos se pagarán, lo cual impacta negativamente la eficiencia (estática). En la misma línea, la presencia de economías de escala debería presionar a la baja los precios. Por lo tanto hay un \emph{trade-off} entre la competencia dada por la cantidad de firmas y la ineficiencia debida a costos fijos y economías de escala.

\end{solution}

%-------------------------------------------------------------------- Cournot con costo fijo
\needspace{8\baselineskip}
\tema{Número de firmas en Cournot con costo fijo y bienestar}
\question Considérese una industria homogénea con $N$ empresas perfectamente simétricas. Las empresas deciden de forma simultánea la cantidad que llevarán al mercado, y tienen un costo de producción $C = cq + F$, donde $c$ es el costo marginal y $F$ un costo fijo. La demanda del mercado se caracteriza por $P = 1 - Q$, siendo $P$ el precio de mercado y $Q$ la producción total.
\begin{parts}
\part Plantee la maximización de beneficios de una empresa $i$ que elige $q_i$.
\needspace{5\baselineskip}
\begin{solution}
\[
  \max_{q_i}\ \pi_i = \Bigl(1 - q_i - \sum_{j\neq i}q_j\Bigr)q_i - cq_i - F .
\]

\end{solution}

\part Derive las condiciones de primer orden e imponga la simetría para llegar a la cantidad de equilibrio.
\needspace{5\baselineskip}
\begin{solution}
La CPO es $1 - 2q_i - \sum_{j\neq i}q_j - c = 0$, es decir, $q_i = \frac{1 - c - \sum_{j\neq i}q_j}{2}$. Imponiendo $q_i = q_j = q^*$:
\[
  q^* = \frac{1 - c}{N + 1}.
\]

\end{solution}

\part Encuentre, con la cantidad de equilibrio, el precio de equilibrio, el cual debe depender del número de firmas en el mercado. ¿Qué efecto tiene $N$ sobre $P$ y sobre el excedente del consumidor?
\needspace{5\baselineskip}
\begin{solution}
$Q^* = Nq^* = \frac{N(1-c)}{N+1}$, de modo que
\[
  P^* = 1 - \frac{N(1-c)}{N+1} = \frac{1 + Nc}{N+1}.
\]
A mayor $N$, menor $P^*$ y mayor es el excedente de los consumidores.

\end{solution}

\part ¿Cuál es el efecto de un mayor $N$ sobre los beneficios de las empresas y el excedente de los productores?
\needspace{5\baselineskip}
\begin{solution}
Los beneficios de cada firma, incluyendo el costo fijo, son
\[
  \pi_i = (P^* - c)q^* - F = \frac{(1-c)^2}{(N+1)^2} - F ,
\]
que son decrecientes en $N$. El excedente de los productores (suma de los beneficios de las $N$ empresas) es
\[
  EP = \frac{N(1-c)^2}{(N+1)^2} - NF .
\]

\end{solution}

\part ¿Qué ocurre con el excedente del productor encontrado cuando $N$ aumenta?
\needspace{5\baselineskip}
\begin{solution}
El excedente del productor de la industria disminuye con $N$ (el primer término decrece a partir de $N = 1$ y el segundo término, $-NF$, decrece linealmente). Observe que a medida que $N$ aumenta y finalmente tiende a infinito, el excedente del productor y el bienestar total pueden volverse negativos: si bien el excedente del consumidor aumenta, el costo fijo total $NF$ crece sin límite y el bienestar total tiende a $-\infty$. De ahí la idea de que puede haber demasiadas empresas en una industria.

\end{solution}
\end{parts}

%-------------------------------------------------------------------- Comente: N exógeno
\needspace{8\baselineskip}
\tema{Comente: número de empresas exógeno en modelos de entrada}
\question \textbf{Comente.} En los modelos de entrada, el número de empresas es exógeno y no depende de las condiciones del mercado.
\needspace{5\baselineskip}
\begin{solution}
Falso. En los modelos de entrada el número de empresas depende del tipo de competencia, de los costos fijos y marginales; por lo tanto, la cantidad de firmas es endógena.

\end{solution}

%-------------------------------------------------------------------- Comente: condición de equilibrio de N
\needspace{8\baselineskip}
\tema{Comente: condición de equilibrio del número de firmas}
\question \textbf{Comente.} La condición para que un número $N$ de empresas esté en equilibrio es que las ganancias sean no negativas para esas empresas, pero negativas si entra una empresa adicional.
\needspace{5\baselineskip}
\begin{solution}
Verdadero. El equilibrio de Nash para la cantidad de firmas se da en el punto en que $\pi^*(N) \geq 0$ y $\pi^*(N+1) < 0$: la entrada de una nueva firma le implicaría utilidades negativas.

\end{solution}

%-------------------------------------------------------------------- Comente: Bertrand con costos fijos
\needspace{8\baselineskip}
\tema{Comente: Bertrand con costos fijos positivos}
\question \textbf{Comente.} En competencia à la Bertrand con costos fijos positivos, el equilibrio suele implicar múltiples empresas compitiendo en precios.
\needspace{5\baselineskip}
\begin{solution}
Falso. En Bertrand con costos marginales simétricos el precio llegaría al costo marginal y los beneficios serían cero. Dados costos fijos positivos, entonces ninguna firma obtiene utilidades no negativas si hay dos o más firmas. \emph{Nota:} solo habrá una firma, que cobra más que el costo marginal pero no el precio que pondría un monopolista normal. Si cobra el precio de monopolio podría atraer la entrada, por lo que, siguiendo la idea de mercados desafiables, la firma cobrará el costo medio.

\end{solution}

%-------------------------------------------------------------------- Comente: diferenciación y entrada
\needspace{8\baselineskip}
\tema{Comente: bienes muy diferenciados y entrada}
\question \textbf{Comente.} Si los bienes están muy diferenciados, la entrada de nuevas empresas puede generar un mayor bienestar social.
\needspace{5\baselineskip}
\begin{solution}
Verdadero. Productos muy diferenciados tienen valor agregado y el \emph{business stealing} es más difícil. La entrada de nuevas firmas implica una mayor variedad para los consumidores, generando mayor excedente del consumidor. Además, mayor competencia induce a menores precios y por tanto mayor excedente del consumidor.

\end{solution}

%-------------------------------------------------------------------- Comente: Sutton
\needspace{8\baselineskip}
\tema{Comente: costos endógenos de entrada (Sutton)}
\question \textbf{Comente.} Según el modelo de Sutton, mayores costos endógenos de entrada están asociados con menor concentración en la industria y por tanto menos competencia.
\needspace{5\baselineskip}
\begin{solution}
Falso. Los costos endógenos tienen un factor estratégico que tiende a concentrar la industria. Un ejemplo de industria con altos costos de entrada endógenos es el \emph{retail}: la publicidad es un factor estratégico que deciden las firmas y que detiene a posibles firmas entrantes de ingresar al mercado. Este es un mercado que, a pesar de estar concentrado, tiene un nivel de competitividad alto.

\end{solution}

%-------------------------------------------------------------------- Comente: mercados desafiables y economías de escala
\needspace{8\baselineskip}
\tema{Comente: mercados desafiables y economías de escala}
\question \textbf{Comente.} El modelo de mercado desafiable contradice la idea de que las economías de escala necesariamente implican poder de mercado.
\needspace{5\baselineskip}
\begin{solution}
Verdadero. Baumol demuestra que las economías de escala no implican poder de mercado si el mercado es desafiable. La amenaza creíble de entrada puede mantener los precios competitivos incluso con un solo productor. \emph{Nota:} si se aprovecha la economía de escala entonces el precio bajaría, pero eso no implica que la amenaza de entrada sea menos creíble.

\end{solution}

%-------------------------------------------------------------------- Comente: entrada de firma eficiente
\needspace{8\baselineskip}
\tema{Comente: entrada de una firma muy eficiente y concentración}
\question \textbf{Comente.} El mercado de la atrapada de dos se compone de firmas simétricas que compiten à la Cournot y enfrentan una demanda lineal. El pasado cuatrimestre irrumpió una nueva firma mucho más eficiente. El mercado ahora está más concentrado y el precio es menor. ¿Cómo se explicaría esta aparente contradicción?
\needspace{5\baselineskip}
\begin{solution}
Si la empresa entrante es mucho más eficiente, el precio tiende a la baja. Por otro lado, se rompe la simetría: para todas las firmas antiguas los beneficios se verán reducidos, puesto que con Cournot asimétrico la firma más eficiente acapara más mercado. Esta reducción de beneficios podría llevar a un menor número de firmas. Por lo tanto, el precio sería menor pero con una mayor concentración de mercado (concentración y competencia no se mueven necesariamente en direcciones opuestas).

\end{solution}

%-------------------------------------------------------------------- Costos de entrada exógenos y tamaño del mercado
\needspace{8\baselineskip}
\tema{Costos de entrada exógenos, tamaño de mercado y concentración}
\question Note que cuando hablamos de costos de entrada no nos referimos exactamente a lo mismo que costos fijos. Explique cómo se relacionan el tamaño del mercado y la concentración (1/$N$) bajo costos de entrada exógenos, distinguiendo entre competencia à la Bertrand y à la Cournot.
\needspace{5\baselineskip}
\begin{solution}
Bajo Bertrand, la cantidad de firmas será siempre una: una firma adicional bajo Bertrand implicaría llevar el precio a un nivel no sostenible (utilidades negativas). En Bertrand con costos fijos, la entrada de una nueva firma solo es posible si logra tener un costo medio más bajo que el de la incumbente; si entra una más eficiente, la incumbente sale del mercado. La concentración no baja con el tamaño del mercado.

En el caso de Cournot, la concentración de mercado disminuye mientras más grande sea el mercado. Pero esto no es una relación 1:1, ya que los costos fijos hacen que el tamaño de mercado tenga un menor efecto sobre la cantidad de firmas. Por ejemplo: si aumenta la población, aumenta la demanda por peluquerías y por aeropuertos; si aumenta el número de peluquerías en 500 nuevos locales, sería absurdo pensar que ahora habrá 500 nuevos aeropuertos. Sujeto a estar en un cartel, los beneficios serán mayores dada la menor competencia, lo que atrae más entrada y menor $1/N$. El gráfico (concentración vs.\ tamaño de mercado) muestra que mayores niveles de competencia podrían estar relacionados a menores niveles de concentración.

\end{solution}

%-------------------------------------------------------------------- Derivación N Cournot
\needspace{8\baselineskip}
\tema{Derivación del número de firmas en Cournot con demanda lineal}
\question Los modelos de entrada piensan la entrada de $N$ firmas como un equilibrio de Nash que cumple: (1) $\pi^*(N) \geq 0$ y (2) $\pi^*(N+1) < 0$. Para caracterizar esta condición es necesario plantear un tipo de competencia y una función de demanda. La demanda del mercado es $Q = S\,q_d$, donde $S$ es el tamaño de mercado y $q_d = a - bP$ la cantidad demandada per cápita. La función de costos es $C_i(q_i) = cq_i + F$, con $c$ el costo marginal y $F$ un costo fijo.
\begin{parts}
\part Llegue a la cantidad producida por cada empresa bajo competencia por cantidades. \emph{(Indicación: llegue a la demanda inversa y defina $\alpha = a/b$, $\beta = 1/(Sb)$.)}
\needspace{5\baselineskip}
\begin{solution}
Despejando, $P = \frac{a}{b} - \frac{Q}{Sb} = \alpha - \beta Q$. Con $N$ firmas, cada una maximiza $q_i\bigl(\alpha - \beta\sum_j q_j - c\bigr) - F$, con CPO $\alpha - c - \beta\sum_{j\neq i}q_j - 2\beta q_i = 0$. Por simetría,
\[
  q^* = \frac{\alpha - c}{(1+N)\beta} = S\,\frac{a - bc}{1+N}.
\]

\end{solution}

\part Dada la cantidad de Cournot, encuentre la función de utilidades de la firma. \emph{(Indicación: exprese primero $Q$ y luego úselo para expresar $P(Q) - c$ y sustitúyalo en $\pi_i$.)}
\needspace{5\baselineskip}
\begin{solution}
La cantidad total es $Q^* = \frac{NS(a - bc)}{1+N}$, de modo que $P(Q^*) = \alpha - \beta Q^* = \frac{a}{b} - \frac{N(a-bc)}{b(1+N)}$ y
\[
  P(Q^*) - c = \frac{a - bc}{b(1+N)}, \qquad
  \pi^*(N) = \bigl(P(Q^*) - c\bigr)q^* - F = \frac{S(a - bc)^2}{b(1+N)^2} - F .
\]

\end{solution}

\part Encuentre el número de firmas del mercado dada la condición de entrada. ¿Qué tan elástica es la cantidad de firmas al tamaño de mercado?
\needspace{5\baselineskip}
\begin{solution}
La condición de entrada es $\pi^*(N) \geq 0$, es decir,
\[
  \frac{S(a-bc)^2}{b(1+N)^2} \geq F
  \quad\Longleftrightarrow\quad
  N^* = (a - bc)\sqrt{\frac{S}{bF}} - 1 .
\]
$N^*$ es creciente en el tamaño de mercado $S$, pero la elasticidad es menor a 1: un aumento de 1\% del mercado se traduce en menos de 1\% más de firmas (en efecto, $\varepsilon_{N,S} = \frac{1}{2}\frac{N^*+1}{N^*} < 1$ para $N^* > 1$).

\end{solution}

\part ¿Qué representa $b$? ¿Cuál es su efecto en el número de firmas? Justifique.
\needspace{5\baselineskip}
\begin{solution}
$b$ es la sensibilidad de la demanda per cápita al precio (cuánto cae la cantidad demandada cuando sube el precio), es decir, la pendiente de la demanda y, por lo tanto, mide cuán elástica es. $\frac{\partial N^*}{\partial b} < 0$: a mayor sensibilidad al precio, el margen de las firmas es menor, los beneficios caen y entran menos firmas.

\end{solution}
\end{parts}

\end{unidad}

\unidadtitulo{Restricciones verticales y doble marginalización}
\begin{unidad}
%-------------------------------------------------------------------- Comente: doble marginalización
\needspace{8\baselineskip}
\tema{Comente: doble marginalización y cómo eliminarla}
\question \textbf{Comente.} Explique qué es el problema de doble marginalización y por qué ocurre. Además, describa las herramientas o estrategias que pueden utilizarse para eliminar dicho problema.
\needspace{5\baselineskip}
\begin{solution}
La doble marginalización se produce cuando tanto el productor como el distribuidor realizan su maximización de beneficios de manera independiente, y cada uno aplica un margen sobre su costo marginal. La consecuencia es una subproducción en relación a la situación de integración vertical (monopolio): el consumidor termina pagando un precio mayor, se transa una cantidad menor y los beneficios totales (la suma de los beneficios del productor y del distribuidor) son menores que en la situación integrada.

Dentro de las soluciones al problema se encuentran la fijación de un precio recomendado de venta, la tarifa en dos partes y la competencia en la distribución. Estas soluciones se diseñan con el fin de que el productor obtenga las rentas monopólicas:
\begin{itemize}
  \item \emph{Precio de reventa recomendado:} se fija igual al precio monopólico, para extraer toda la renta; prohíbe al distribuidor cobrar un margen y, por tanto, elimina el problema.
  \item \emph{Tarifa en dos partes (cargo fijo y variable):} el productor fija $w = c_{mg}$, sin obtener ningún margen. Luego el distribuidor maximiza y obtiene el $p$ y $q$ monopólicos. Finalmente, el productor le quita todo el excedente a través de la tarifa fija.
  \item \emph{Competencia en la distribución:} da como solución $P = w$, con lo cual el margen en la distribución se hace cero y el problema desaparece.
\end{itemize}

\end{solution}

%-------------------------------------------------------------------- Comente: competencia aguas abajo
\needspace{8\baselineskip}
\tema{Comente: competencia entre distribuidores y doble marginalización}
\question \textbf{Comente.} En el contexto de una relación vertical donde el productor tiene todo el poder de negociación, la competencia aguas abajo entre distribuidores idénticos (igual función de costos) siempre disipará las rentas de los distribuidores, lo que eliminará completamente el problema de doble marginalización.
\needspace{5\baselineskip}
\begin{solution}
Falso, ya que no siempre disipará las rentas de los distribuidores. Si los distribuidores aguas abajo compiten en precios, se llegará a que el precio minorista es $P = W + C$, por lo que el productor se llevará todo el beneficio y no habrá doble marginalización. En cambio, si hay poder de mercado aguas abajo (por ejemplo, competencia en cantidades o diferenciación), los distribuidores podrán fijar un precio mayor que su costo unitario, pudiendo así conseguir beneficios, y el problema subsiste.

\end{solution}

%-------------------------------------------------------------------- Comente: precio recomendado
\needspace{8\baselineskip}
\tema{Comente: precios de reventa recomendados}
\question \textbf{Comente.} Los precios recomendados de reventa (o venta minorista) establecidos por parte de un productor son siempre perjudiciales para los consumidores.
\needspace{5\baselineskip}
\begin{solution}
No siempre. Por ejemplo, si la fijación de precio por parte del productor tiene como propósito eliminar la doble marginalización, entonces el consumidor es favorecido con dicha recomendación de precio (pues el precio baja y la cantidad aumenta respecto del caso con doble marginalización).

\end{solution}

%-------------------------------------------------------------------- Comente: poder de compra
\needspace{8\baselineskip}
\tema{Comente: poder de compra en el productor o en el distribuidor}
\question \textbf{Comente.} El consumidor prefiere que el poder de compra esté concentrado en el productor en vez del distribuidor.
\needspace{5\baselineskip}
\begin{solution}
Falso. Cuando el poder de compra está concentrado en el productor, se obtiene el resultado de doble marginalización. Cuando el poder de compra se concentra en el distribuidor, se obtiene un resultado normal de monopolio. Dado que la doble marginalización es un peor resultado que el monopolio para el consumidor, se debería preferir que el poder de compra se concentre en el distribuidor.

\end{solution}

%-------------------------------------------------------------------- Productor y distribuidor (Q=a-bP)
\needspace{8\baselineskip}
\tema{Relación vertical con demanda lineal general}
\question Un fabricante $M$ vende su producto a través de un distribuidor independiente. El costo de fabricación del bien es $c$, y la demanda por el producto final es $Q(P) = a - bP$. El costo variable de distribución es $d$.
\begin{parts}
\part Suponga que la transacción entre $M$ y $D$ ocurre a través de un precio mayorista $W$ únicamente, mientras que el precio minorista $P$ lo fija el distribuidor. Obtenga el precio mayorista y minorista resultante si todo el poder negociador lo tiene $M$.
\needspace{5\baselineskip}
\begin{solution}
El juego se resuelve mediante inducción hacia atrás. La firma productora prevé los beneficios que obtendrá la distribuidora a cualquier $W$ dado, y luego incorpora la cantidad entregada al mercado en su propia maximización.

\emph{Distribuidora:} $\max_P\ \pi_D = (P - W - d)(a - bP)$, con CPO $a - 2bP + bW + bd = 0$:
\[
  P^* = \frac{a + bW + bd}{2b}, \qquad Q(P^*) = \frac{a - bW - bd}{2}.
\]
\emph{Productora:} $\max_W\ \pi_P = (W - c)\,\frac{a - bW - bd}{2}$, con CPO $a - bd - 2bW + bc = 0$:
\[
  W^* = \frac{a - bd + bc}{2b}, \qquad P(W^*) = \frac{3a + bd + bc}{4b}.
\]
Este es el precio que pondría la firma distribuidora.

\end{solution}

\part Resuelva la pregunta (i) pero asumiendo ahora que todo el poder negociador lo tiene $D$.
\needspace{5\baselineskip}
\begin{solution}
Como la firma distribuidora tiene todo el poder negociador, escogerá un $W$ tal que maximice sus beneficios (es decir, el mínimo posible). Además, como la firma productora debe participar en el mercado, el $W$ escogido debe asegurar $\pi_P^* \geq 0$. Esto implica $W = c$. De esta forma, el precio final que pondrá la firma distribuidora será
\[
  P^* = \frac{a + bc + bd}{2b}.
\]

\end{solution}

\part Suponga que todo el poder negociador reside en $M$, pero esta empresa no puede forzar a $D$ a aceptar un contrato donde también se fije el precio minorista $P$. Obtenga el acuerdo óptimo entre $M$ y $D$ en base a los precios $W$ y $P$, que induce a $D$ a aceptarlo (en caso de no acuerdo, se aplica el contrato de (i)).
\needspace{5\baselineskip}
\begin{solution}
Ambas partes ganarían más si $P = P_M$, donde
\[
  P_M = \arg\max_P\ (P - c - d)(a - bP) = \frac{a + bc + bd}{2b}.
\]
La negociación ocurre entonces sobre $W$, el precio mayorista. Para que $D$ acepte el trato, $W$ debe dejarle al menos el mismo beneficio que la opción en (i), al que llamamos $\pi_D^0$. Así, se debe cumplir
\[
  (P_M - W - d)\,Q(P_M) \geq \pi_D^0 .
\]
$M$ elige el $W$ más alto que cumple esta restricción con igualdad, y se queda con el resto de la renta monopólica.

\end{solution}
\end{parts}

%-------------------------------------------------------------------- Productor y distribuidor numérico
\needspace{8\baselineskip}
\tema{Doble marginalización: contratos y esfuerzo de venta}
\question Una empresa monopolista $M$ vende su producto a través de un distribuidor $D$. La demanda por el bien final es $Q(P) = 120 - P$, donde $P$ es el precio minorista del bien. El costo variable de producir el bien es $C_M = 20$ y el costo de distribución es $C_D = 10$. Inicialmente, todo el poder negociador está en $M$.
\begin{parts}
\part Si el trato entre ambas partes se realiza mediante un precio mayorista unitario $W$, ¿cuál sería el precio que $M$ fijaría a $D$? Identifique el fenómeno de la doble marginalización en este caso.
\needspace{5\baselineskip}
\begin{solution}
El distribuidor, dado $W$, maximiza $\pi_D = (P - W - 10)(120 - P)$. Derivando respecto a $P$ e igualando a cero: $130 + W - 2P = 0$, de donde
\[
  P = \frac{130 + W}{2}, \qquad Q = \frac{110 - W}{2}.
\]
Así, el monopolista maximiza $\pi_M = (W - 20)\,\frac{110 - W}{2}$. Derivando: $\frac{110 - W}{2} - \frac{W - 20}{2} = 0$, lo que da
\[
  W = 65, \qquad P = 97{,}5, \qquad Q = 22{,}5 .
\]
El fenómeno de doble marginalización se observa en que ambas firmas fijan un precio por encima de sus costos unitarios: $C_M = 20$ para el monopolista y $W + C_D = 75$ para el distribuidor, obteniendo beneficios positivos: $\pi_M = 45\cdot 22{,}5 = 1012{,}5$ y $\pi_D = 22{,}5^2 = 506{,}25$ (total $1518{,}75$, menos que los $2025$ de la integración vertical).

\end{solution}

\part Proponga y desarrolle dos soluciones contractuales que permitan eliminar el problema de la doble marginalización.
\needspace{5\baselineskip}
\begin{solution}
Dado que el poder negociador está en el monopolista, este podría ofrecer un contrato $(P, W)$, fijando el precio minorista y mayorista, de modo que la firma actúe como si estuviera verticalmente integrada y fijara $W = P - C_D$. Entonces el monopolista maximiza $\max_P\ (P - 20 - 10)(120 - P)$, con CPO $150 - 2P = 0$:
\[
  P = 75, \qquad W = 75 - 10 = 65 .
\]
Así, el monopolista se lleva todo el beneficio, $\pi_M = (65 - 20)\cdot 45 = 2025$, y el distribuidor nada: $\pi_D = (75 - 65 - 10)\cdot 45 = 0$.

Otra solución es una \emph{tarifa en dos partes} $(T, W)$, donde $T$ es la tarifa fija que le cobra al distribuidor y $W$ el precio por unidad. Así, $W = C_M = 20$ y $T$ es el beneficio que obtendría el distribuidor: $T = (75 - 30)\cdot 45 = 2025$. Nuevamente $\pi_M = 2025$ y $\pi_D = 0$.

\end{solution}

\part Suponga ahora que el distribuidor puede ejercer un esfuerzo de venta $e$, de modo tal que la demanda por el bien se incrementa: $Q(P,e) = 120 - P + e$, donde el costo por ejercer el esfuerzo es $c(e) = \frac{e^2}{2}$. Por simplicidad, ahora asuma que $C_D = 0$. Plantee un contrato de precio recomendado minorista $(W, P)$. ¿El productor logra maximizar el beneficio de la industria? Además, ¿el productor logra quedarse con dicho beneficio?
\needspace{5\baselineskip}
\begin{solution}
Por inducción hacia atrás, el distribuidor elegirá cuánto esforzarse dados $(W,P)$ fijados por el productor:
\[
  \max_e\ \pi_D = (P - W)(120 + e - P) - \frac{e^2}{2}
  \;\Rightarrow\; P - W - e = 0
  \;\Rightarrow\; e = P - W .
\]
Así, el distribuidor se esforzará más mientras mayor sea la distancia entre $P$ y $W$. Pero de la parte (b) sabemos que el monopolista tiene incentivos a fijar en este caso $P = W$ (para no dejarle margen al distribuidor), y entonces el esfuerzo será $0$. Si permitiera que el precio minorista se distanciara del mayorista, se podría capturar un mayor nivel de beneficios de forma conjunta. Por lo tanto, el productor \emph{no} logra maximizar el beneficio de la industria, y solo se queda con el beneficio ya encontrado antes (con $C_D = 0$: $\pi_M = 2500$, con $P = W = 70$). El beneficio conjunto máximo, en cambio, resultaría de elegir $P$ y $e$ para maximizar $(P - 20)(120 - P + e) - \frac{e^2}{2}$, lo que da $e = P - 20$, $P = 120$, $e = 100$ y un beneficio de $5000$.

\end{solution}
\end{parts}

%-------------------------------------------------------------------- Delegación a distribuidores
\needspace{8\baselineskip}
\tema{Cournot con distribuidores (precios de transferencia)}
\question Suponga que dos firmas enfrentan una demanda $Q = 1 - P$, compiten en cantidades y no enfrentan costos de producción ni de distribución.
\begin{parts}
\part ¿Cuál es el equilibrio en este mercado?
\needspace{5\baselineskip}
\begin{solution}
Cada firma maximiza $\pi_i = (1 - q_i - q_j)q_i$. Derivando: $1 - 2q_i - q_j = 0$, por lo tanto $q_i = \frac{1 - q_j}{2}$. Por simetría, reemplazando una función de reacción en la otra se tiene $q_i = q_j = \frac{1}{3}$. Por lo tanto,
\[
  Q = \frac{2}{3}, \qquad P = 1 - \frac{2}{3} = \frac{1}{3}, \qquad \pi_i = \pi_j = \frac{1}{9}.
\]

\end{solution}

\part Suponga ahora que cada firma decide vender su producto a través de un distribuidor. En términos prácticos, cada firma le vende a su distribuidor todas las unidades que desee a un precio $w_i$ ($i = 1,2$), y luego las firmas distribuidoras compiten en cantidades. ¿Cuál es el nuevo equilibrio de mercado (precio, precio de transferencia $w$, cantidad, beneficios, cantidad agregada y beneficios agregados)?
\needspace{5\baselineskip}
\begin{solution}
El \emph{timing} es el siguiente: en $T=1$, el productor le vende toda la cantidad del producto al distribuidor a un precio $w_i$; en $T=2$, los distribuidores compiten en cantidades. En $T=2$ cada distribuidor maximiza $\pi_{di} = (1 - q_i - q_j - w_i)q_i$, con CPO $1 - 2q_i - q_j - w_i = 0$, es decir, $q_i = \frac{1 - q_j - w_i}{2}$ (y análogamente para $j$). Sumando ambas expresiones:
\[
  Q = \frac{2 - w_i - w_j}{3}, \qquad P = 1 - Q = \frac{1 + w_i + w_j}{3}, \qquad q_i = \frac{1 - 2w_i + w_j}{3}.
\]
Ahora, en $T=1$, cada productor maximiza $\pi_{pi} = w_i q_i = w_i\,\frac{1 - 2w_i + w_j}{3}$. Derivando: $\frac{1 + w_j - 4w_i}{3} = 0$. Intersectando las funciones de reacción se tiene $w_i = w_j = \frac{1}{3}$. Así,
\[
  q_i = q_j = \frac{2}{9}, \qquad Q = \frac{4}{9}, \qquad P = 1 - \frac{4}{9} = \frac{5}{9},
\]
\[
  \pi_{di} = \pi_{dj} = \Bigl(\frac{5}{9} - \frac{1}{3}\Bigr)\frac{2}{9} = \frac{4}{81}, \qquad
  \pi_{pi} = \pi_{pj} = \frac{1}{3}\cdot\frac{2}{9} = \frac{2}{27}.
\]
Los beneficios agregados (productores y distribuidores) son $2\bigl(\frac{4}{81} + \frac{6}{81}\bigr) = \frac{20}{81}$.

\end{solution}

\part ¿Ocurre un fenómeno de doble marginalización en este caso?
\needspace{5\baselineskip}
\begin{solution}
La doble marginalización surge del hecho de que, al maximizar sus beneficios, tanto el productor como el distribuidor aplican un margen sobre su respectivo costo marginal. Como se puede ver en el desarrollo del ejercicio, la doble marginalización sí ocurre: el precio de transferencia es $w = \frac{1}{3} > 0$ y el precio minorista ($\frac{5}{9}$) es mayor que el de Cournot sin distribuidores ($\frac{1}{3}$).

\end{solution}
\end{parts}

\end{unidad}

\unidadtitulo{Subastas}
\begin{unidad}
%-------------------------------------------------------------------- Comente: maldición del ganador
\needspace{8\baselineskip}
\tema{Comente: maldición del ganador en subastas}
\question \textbf{Comente.} Si las empresas participan en una subasta inglesa o de segundo precio, ¿por qué entonces ocurriría la maldición del ganador?
\needspace{5\baselineskip}
\begin{solution}
La maldición del ganador de una subasta o licitación ocurre cuando el ganador ofrece un valor mayor al valor real del bien que se adjudica, aun cuando pague el segundo precio. Bajo asimetrías de información respecto del valor real del bien subastado, los oferentes sobrestiman el valor; luego, el ganador (aquel con la estimación más optimista) paga un precio mayor al valor o retorno del bien, lo que se traduce en pérdidas.

\end{solution}

\end{unidad}

\end{document}
